% The next line makes it so it compiles the main file if I hit the compile buttonf
% when editing this file in TeXworks.
% !TEX root = ../AIDMA.tex
\chapter{Sets, Functions, and Relations}\label{chapter:sets}

\section{Sets}\label{section:sets}

\subsection{Definitions}\label{section:setDefs}


\begin{df}
\index{set}
\index{element, of a set}
\index{a sets elementof@$\in$ (element of set)}
\index{a sets elementofnot@$\not\in$ (not element of set)}
Sets

\begin{itemize}
\item
A {\bf set} is an unordered collection of
objects. 
\item The objects in the set are called the {\bf elements} of the set.
\item
 If $a$ belongs to the set $A$, then we write $a\in
A$, read ``$a$ is an element of $A$.''
\item
If $a$ does not belong to the set $A$, we write $a\not\in A$, read ``$a$ is  
not an element of $A$.'' 
\item
Generally speaking, repeated elements in a set are ignored.
\end{itemize}
\end{df}

\begin{rem}
The symbol $\in$ should be read as {\bf is an element of}, not {\em exists in}.
\end{rem}

\begin{exa}\label{setsexa1}
The sets $A=\{1, 2, 3\}$, $B=\{3, 2, 1\}$,  and $C=\{1, 1, 1, 2, 2, 3\}$ actually represent the same set
since repeated values are ignored and the order elements are listed
does not matter.  Notice that $1\in A$ and $3\in A$, but $4\not\in A$.

Let $D=\{0,1,2,3,4,5,6,7,8,9\}$ be the set of decimal
digits. Then $4\in D$ but $11\nin D$.  

%The cardinality of each of these sets is 3.
\end{exa}

Notice that the elements in a set are listed between curly braces.  
Thus, $\{1,2,3\}$ is a set (where order does not matter and duplicates are ignored), 
but $[1,2,3]$ is a list (where order {\em does} matter and duplicates are allowed).
Also, 1, 2, 3 is just a list of three numbers whereas $\{1,2,3\}$ is the set containing
the numbers 1, 2, and 3. 


\begin{df}
\index{cardinality, set}
\index{set!cardinality}
\index{set!size}
\index{a sets card@$\vert A \vert$ (set cardinality)}
Cardinality

\begin{itemize}
\item
The number of elements in a set $A$, also known as the 
the {\bf cardinality} of $A$, will be denoted by $|A|$.
\item
If $|A|$ is finite, we call $A$ a {\bf finite set}\index{finite set}\index{set!finite}.
\item
If the set $A$ has infinitely many elements, we write $|A| = \infty$
and we refer to $A$ as an {\bf infinite set}\index{infinite set}\index{set!infinite}.
\end{itemize}
\end{df}


\begin{exa}
If $A$, $B$, $C$, and $D$ are the sets from Example~\ref{setsexa1},
 then $|A|=3$, $|B|=3$, $|C|=3$, and $|D|=10$.
\end{exa}

\begin{exer}
Give the set of prime numbers less than $10$. What is its cardinality?

\noindent
\answerline
\begin{answer}
The prime numbers less than $10$ are $2$, $3$, $5$, and $7$.
But the problem asked for the {\em set of} prime numbers less than $10$.
Therefore, the answer is $\{2, 3, 5, 7\}$.
If you were asked to {\em list} the prime numbers less than $10$,
an appropriate answer would have been $2,3,5,7$ (but that is not what was asked).
The cardinality of the set is 4. That is, $|\{2, 3, 5, 7\}|=4$.
\end{answer}
\end{exer}


Since we cannot list every element of an infinite set, we need a way of expressing the set so that it is 
clear what elements it contains. If the elements of the set follow some pattern, it is common to list the
first several elements and then conclude with $\ldots$, indicating that the pattern continues.
There is no ``right'' number of elements to list when using this notation, but there needs to be enough
so that the pattern is evident. Often 3-5 elements suffices.

\begin{exa}
The set of positive integers can be expressed as $\BBZ^+=\{1, 2, 3, \ldots\}$. 
Notice that $|\BBZ^+|=\infty$.

The set of positive integers that are a multiple of $5$ can be expressed as $\{5, 10, 15, 20, \ldots\}$.
Hopefully it is clear that $|\{5, 10, 15, 20, \ldots\}|=\infty$.

The set of integer multiples of $5$ can be expressed as 
$\{\ldots, -15, -10, -5, 0, 5, 10, 15, \ldots\}.$
Hopefully it is clear that this is also an infinite set.
\end{exa}


\begin{df}
We say two sets are {\bf equal} if they contain the same elements.
That is $\forall x(x\in A \leftrightarrow x\in B)$.
If $A$ and $B$ are equal sets, we write $A=B$. %\\[-3ex] % bad box break.
\end{df}

\begin{rem} We will normally denote sets by capital letters, like  $A, B, S ,\naturals$, etc. 
Elements will be denoted by lowercase letters, like $a, b, r$, etc.
\end{rem}

\begin{exer}
Let $A=\{1,2,3,4,5,6\}$, $B=\{1,2,3,4,5,4,3,2,1\}$,  
$C=\{6,3,4,5,1,3,2\}$.
Then $|A|=\blankline{.5}$, $|B|=\blankline{.5}$, and $|C|=\blankline{.5}$.

Which of $A$, $B$, and $C$ represent the same sets? \blanklinefill
\begin{answer}
6; 5; 6; $A$ and $C$ represent the same set. That is, $A=C$.
\end{answer}
\end{exer}


\begin{df}
The following notation is pretty standard, and we will follow it in this book.

\noindent
\begin{center}
\begin{tabular}{ll}
$\naturals= \{0,1,2, 3, \ldots\}$ &  the {\bf natural numbers}.\index{natural numbers}
\index{a sets known Nat@$\naturals$ (natural numbers)}\\
$\BBZ = \{\ldots -2, -1, 0, 1,2, \ldots\}$& the {\bf integers}.\index{integers}
\index{a sets known Z@$\BBZ$ (integers)}\\
$\BBZ^+ = \{1,2,3, \ldots\}$& the {\bf positive  integers}.\index{positive integers}
\index{a sets known Z P@$\BBZ^+$ (positive integers)}\\
$\BBZ^- = \{-1,-2,-3, \ldots\}$& the {\bf negative integers}.\index{negative integers}
\index{a sets known Z Q@$\BBZ^-$ (negative integers)}\\
$\BBQ$& the {\bf rational numbers}.\index{rational numbers}
\index{a sets known Rational@$\BBQ$ (rational numbers)}\\
$\BBR$& the {\bf real numbers}.\index{real numbers}
\index{a sets known Real@$\BBR$ (real numbers)}\\
$\BBC$& the {\bf complex numbers}.\index{complex numbers}
\index{a sets known C@$\BBC$ (complex numbers)}\\
$\varnothing=\{\}$ &  the {\bf empty set} or {\bf null set}.
\index{empty set}\index{null set}\index{set!empty}
\index{a sets known ZZ@$\varnothing$ (empty set)}
\index{a sets known ZZ@$\{\}$ (empty set)}\\
\end{tabular}
\end{center}
\end{df}

\begin{rem}
There is no universal agreement of the definition of $\naturals$.  
Although here it is defined as
$\{0,1,2, 3, \ldots\}$, it is sometimes defined as $\naturals = \BBZ^+$.
The only difference is whether or not $0$ is included.
I prefer the definition given here because then we have a notation for the positive integers ($\BBZ^+$)
as well as the non-negative integers ($\naturals$).
\end{rem}


\begin{exa}
Notice that $|\naturals| = |\BBZ|= |\reals|=\infty$.  But this may be a bit misleading.  
Do all of these sets have the same number of elements?  Believe it or not, it turns out that 
$\naturals$ and $\BBZ$ do, but that $\reals$ has many more elements than both of these.
If it seems strange to talk about whether or not two infinite sets have the same number of elements, don't
worry too much about it.  We probably won't bring it up again.
\end{exa}

\begin{exer}
(a)  $|\BBC|=\blankline{.75}$,  
(b) $|\BBZ^+|=\blankline{.75}$, 
(c)  $|\varnothing|=\blankline{.75}$
\begin{answer}
$\infty$; $\infty$.  You might think it is $\infty/2$, but you can't do arithmetic with $\infty$ since it
isn't a number. Without getting too technical, although $\BBZ^+$ seems to have about half as many elements
as $\BBZ$, it actually doesn't.  It has the exact same number: $\infty$. ; $0$.
\end{answer}
\end{exer}

Another other common notation that is used express sets is called {\em set builder notation}. 
It is easier to understand through examples than by giving a formal definition.

\begin{exa}
Let $S$ be the set of the squares of integers.  
We can express this as
$S = \{n^2|n\in \BBZ\}$ or $S = \{n^2:n\in \BBZ\}$.  
We call this {\em set builder notation}.
We read the $:$ or $|$ as ``such that.''
Thus, $S$ is the set containing numbers of the form $n^2$ such that $n$ is an integer.
\end{exa}


\begin{exa}
Consider the set of integers that are one less than a positive power of 2. 
This set contains elements such as $1=2^1-1$, $3=2^2-1$, and $255=2^8-1$.
We can express this set in set builder notation as
$\{2^n-1|n\in\BBZ^+\}$.

We could also use the notation $\{2^1-1,2^2-1,2^3-1,\ldots\}$, since the pattern is evident.
However, it would be unwise to write it as $\{1,3,7,15,\ldots\}$ since it may or may not be
evident what the pattern is when expressed in this way.
\end{exa}


\begin{exer}
Use two different notations to express the set of even integers.

\noindent
\answerline
\begin{answer}
$\{2a : a\in\integers \}$ and $\{\ldots, -4, -2, 0, 2, 4, \ldots\}$.
\end{answer}
\end{exer}

\vspace*{-.1in} % Getting a bad break after example on bottom of page

\begin{exa}
Let $T$ be the set of all integers that can be expressed as the sum of the square of two positive integers.
This set contains elements such as $2=1^2+1^2$, $5=1^2+2^2$, and $25=3^2+4^2$.
Then $T=\{n^2+m^2|n,m\in\BBZ^+\}$.

In this case, expressing the set as something like $\{2, 5, 8, 9,\ldots \}$ does not make sense at all because there
is no way of discerning a pattern.
\end{exa}

\vspace*{-.1in} % Getting a bad break after example on bottom of page

\begin{exa}
Use set builder notation to express $\BBC$, the set of complex numbers.

{\bf Solution:} 
$\BBC = \{a+bi : a,b\in\reals \}$.
\end{exa}

\begin{exer}
Use set builder notation to express $\BBQ$, the set of all rational numbers.
\answerline
\begin{answer}
$\BBQ = \{a/b : a,b\in\integers,b\neq 0 \}$.
\end{answer}
\end{exer}

\begin{df}
\index{subset}
Subsets

\begin{itemize}
\item
If every element in $A$ is also in $B$, we say that 
$A$ is a {\bf subset} of $B$ and we write this as 
$A \subseteq B$.  
\index{subset!proper}
\index{a sets opers sub@$\subseteq$ (subset)}
\item
If $A\subseteq B$ and there is some $x\in B$ such that $x\not\in A$,
then we say A is a {\bf proper subset} of $B$, denoting it by $A\subset B$.
\index{proper subset}
\index{a sets opers sub prop@$\subset$ (proper subset)}
\item
If there is some $x \in A$ such that $x
\not\in B$, then $A$ is not a subset of $B$, which we write as $A
\not\subseteq B.$ 
\index{a sets opers sub not@$\not\subseteq$ (not a subset)}
\end{itemize}
\end{df}

\begin{rem}
Some authors use $\subset$ to mean the same thing as $\subseteq$. 
You will need to consider the context in order to interpret it correctly.
\end{rem}


\begin{exa}
Let $S = \{1, 2, \ldots, 20\}$, that is, the set of integers
between $1$ and $20$, inclusive. 
Let $E = \{2, 4, 6, \ldots , 20\}$, the set of all even integers between $2$ and $20$, inclusive. 
Notice that $E\subseteq S$.
Let $P = \{2, 3, 5, 7, 11, 13, 17, 19\}$, the set of primes less than $20$. 
Then $P\subseteq S$, but $P\not\subseteq E$ and $E\not\subseteq P$.
\end{exa}


\begin{exer}
Let  $S = \{n^2|n\in \BBZ\}$ and  $A = \{ 1, 4, 9, 16\}$.
Answer each of the following, including a brief justification.
\begin{lenum}
\item
Is $A\subseteq S$?\blanklinefill
\item
Is $A\subset S$? \blanklinefill
\item
Is $S\subseteq S$? \blanklinefill
\item
Is $S\subset S$? \blanklinefill
\item
Is $S\subset A$? \blanklinefill
\end{lenum}
\begin{answer}
(a) Yes.
(b) Yes.  $A$ is a proper subset since $25$, for instance, is in $S$ but not in $A$.
(c) Yes.  Every set is a subset of itself.
(d) No.  No subset is a {\em proper subset} of itself.
(e) No.  $25\in S$, but $25\not\in A$.
\end{answer}
\end{exer}

\newpage % bad break. Force it.

\begin{exer}
Let $A$ be the set of integers divisible by $6$, $B$ be the set of integers divisible by $2$, and
$C$ be the set of integers divisible by $3$.
Answer each of the following, giving a brief justification.
\begin{enumerate}[(a)]
\item Is $A\subseteq B$?\blanklinefill
\item Is $A\subseteq C$?\blanklinefill
\item Is $B\subseteq A$?\blanklinefill
\item Is $B\subseteq C$?\blanklinefill
\item Is $C\subseteq A$?\blanklinefill
\item Is $C\subseteq B$?\blanklinefill
\end{enumerate}
\begin{answer}
(a) yes.  Any number that is divisible by 6 is divisible by 2.;
(b) yes. Any number that is divisible by 6 is divisible by 3.; 
(c) no.  $4\in B$, but $4\not\in A$.; 
(d) no.  $4\in B$, but $4\not\in C$.; 
(e) no.  $3\in C$, but $3\not\in A$.; 
(f) no.  $3\in C$, but $3\not\in B$.
\end{answer}
\end{exer}


\begin{exa}
The set $$S=\{\mathrm{Roxan, Jacquelin, Sean, Fatimah, Wakeelah, Ashley, Ruben,
Leslie, Madeline}\}$$
is the set of students in a particular course. 
This set can be split into two subsets: the
set $ F=\{\mathrm{Roxan, Jacquelin,
Fatimah, Wakeelah, Ashley, Madeline}\}$
of females in the class, and the set
$M=\{\mathrm{ Sean, Ruben, Leslie}\}$ of males in the class.  Thus we have $F\subseteq S$ and
$M \subseteq S$. Notice that it is {\em not true} that $F\subseteq
M$ or that $M \subseteq F$. Put another way, $F\not\subseteq M$ and $M \not\subseteq F$
\end{exa}



\begin{exa}\label{exa:subsets3element}
Find all the subsets of $\{a, b, c\}$.
\begin{sol}
They are  $\varnothing, \{a\}, \{b\}, \{c\}, \{a, b\},
\{b, c\}, \{a, c\}$, and  $\{a, b, c\}$.
\end{sol}
Notice that there are 8 subsets.  Also notice that $8=2^3$.  As we will see shortly, that is not a coincidence.
\end{exa}

Notice that we wrote $\varnothing$ and not $\{\varnothing\}$ in the previous example.  
It turns out that
$\varnothing\neq \{\varnothing\}$.  $\varnothing$ is the empty set--that is, the set that has no elements.
$\{\varnothing\}$ is the set containing the empty set.  
Thus, $\{\varnothing\}$ is a set containing the single element $\varnothing$. 
You can use either $\varnothing$ or $\{\}$ to denote the empty set, but not $\{\varnothing\}$. 

\newpage % Got a bad break.

\begin{exer}\label{exa:subsets4element}
Find all the subsets of $\{a, b, c, d\}$.
\vspace*{2.25in}
\begin{answer}
 We will use the result of example
\ref{exa:subsets3element}.
A subset of $\{a, b, c, d\}$
either contains $d$ or it does not.
Since the subsets of $\{a, b, c\}$ do not contain $d$, we simply list all the subsets of $\{a, b, c\}$ and then to each one of them  we add $d$. This gives

$\begin{array}{lllllll} S_1 & = & \varnothing & \hspace{1cm} &  S_9 & = & \{d\} \\
S_2 & = & \{a\} & \hspace{1cm} &  S_{10} & = & \{a, d\} \\
S_3 & = & \{b\}& \hspace{1cm} &  S_{11} & = & \{b, d\} \\
S_4 & = & \{c\} & \hspace{1cm} &  S_{12} & = & \{c, d\} \\
S_5 & = & \{a, b\}& \hspace{1cm} &  S_{13} & = & \{a, b, d\} \\
S_6 & = & \{b, c\}& \hspace{1cm} &  S_{14} & = & \{b, c, d\}  \\
S_7 & = & \{a, c\} & \hspace{1cm} &  S_{15} & = & \{a, c, d\} \\
S_8 & = & \{a, b, c\} & \hspace{1cm} &  S_{16} & = & \{a, b, c, d\}  \\
\end{array}   $
\end{answer}
\end{exer}


\begin{df}
\index{set!power}
\index{power set}
\index{a sets opers powerset@$P(A)$ (power set)}
The {\em power set} of a set is the set of all subsets of a set. The
power set of a set $A$  is denoted by $P(A)$.
\end{df}

\begin{exa}\label{exa:powerset}
If $A=\{a, b, c\}$,  example \ref{exa:subsets3element} implies that
$P(A)=\{  \varnothing, \{a\}, \{b\}, \allowbreak\{c\}, \allowbreak\{a, b\}, \allowbreak\{b, c\},\allowbreak \{a, c\},\allowbreak \{a, b, c\} \}$.
Notice that the solution is a set, the elements of which are also sets.

An {\em incorrect answer} would be 
$\{  \varnothing, a , b, c , \{a, b\} , \{b, c\} , \{a, c\}, \{a, b, c\} \}$.
This is incorrect because $a$ is not the same thing as $\{a\}$ (the set containing $a$).  
$\{a\}\in P(A)$, but $a\not\in P(A)$.  This is a subtle but important distinction.
\end{exa}


\begin{exer}\label{exa:pset4element}
Find $P(\{a, b, c, d\})$.
\vspace*{2.25in}
\begin{answer}
Based on the answer to Exercise~\ref{exa:subsets4element}, we have that
$P(\{a, b, c, d\})=\{ \varnothing,
 \{a\}, \{b\},  \allowbreak\{c\}, \allowbreak\{a, b\},\allowbreak\{b, c\}, 
\allowbreak\{a, c\}, \allowbreak\{a, b, c\},\allowbreak  \{d\},\allowbreak \{a, d\}, 
 \allowbreak\{b, d\},\allowbreak \{c, d\}, \allowbreak \{a, b, d\},\allowbreak \{b, c, d\}, 
\allowbreak\{a, c, d\}, \allowbreak \{a, b, c, d\} 
\}$.
Notice that a list of these 16 sets not separated by commas and not enclosed in $\{\}$ is
not correct.  It may have the correct {\em content}, but it is not in the proper {\em form}.
\end{answer}
\end{exer}

We will prove the following theorem in the next section after we have developed the appropriate notation to do so.
\begin{thm}\label{thm:powerSetSize}
Let $A$ be a set with $n$ elements.  Then $|P(A)|=2^n$. 
%\begin{pf}
%We use induction\footnote{We will cover induction more fully and formally later.
%But since this use of induction is pretty intuitive, 
%especially in light of Example~\ref{exa:subsets4element}, it serves as a useful foreshadowing of things to come.}
% and the idea from the solution to Exercise~\ref{exa:subsets4element}. 
% Clearly if $|A|= 1$, $A$
%has $2^1 = 2$ subsets: $\varnothing$ and $A$ itself.\\
%Assume every
%set with $n - 1$ elements has $2^{n-1}$ subsets. Let $A$ be a set
%with $n$ elements.  Choose some $x\in A$.  Every subset of $A$
%either contains $x$ or it doesn't.  
%Those that do not contain $x$  are subsets of $A\setminus \{x\}$.  Since
%$A\setminus \{x\}$ has $n - 1$ elements, the induction hypothesis implies that it has
%$2^{n-1}$ subsets.
%Every subset that does contain $x$ corresponds to one of the subsets of $A\setminus \{x\}$ with the element $x$ added.  That is, for each subset $S \subseteq A\setminus \{x\}$,  $S\cup \{x\}$ is a subset of $A$ containing $x$.
%Clearly there are $2^{n-1}$ such new subsets.
%Since this accounts for all subsets of $A$, $A$ has $2^{n-1} + 2^{n-1} = 2^n$ subsets.  
%\end{pf}
\end{thm}

\begin{exer}
Let $A$ be a set with 4 elements.
\begin{lenum}
\item
$|P(A)|=\blankline{1.75}$.
\item
$|P(P(A))|=\blankline{1.75}$.
\item
$|P(P(P(A)))|=\blankline{1.75}$.
\end{lenum}
\begin{answer}
(a) By Theorem~\ref{thm:powerSetSize}, $|P(A)|=2^4=16$.
(b) Similarly, $|P(P(A))|=2^{16}=65536$.
(c) This is just getting a bit ridiculous, but the answer is $|P(P(P(A)))|=2^{65536}$.
\end{answer}
\end{exer}

\begin{exer}
If one element is added to a finite set $A$, how much larger is the power set of $A$ after the element is added
(relative to the size of the power set before it is added)? Explain your answer.

\noindent
\answerlinethree
\begin{answer}
Applying Theorem~\ref{thm:powerSetSize}, it is not too hard to see that the power set will be
twice as big after a single element is added.
\end{answer}
\end{exer}

%%%%%%%%%%%%%%%%%%%%%%%%%%%%%%%%%%%%%%%%%%%
\newpage
\subsection{Set Operations}\label{section:setOps}
%%%%%%%%%%%%%%%%%%%%%%%%%%%%%%%%%%%%%%%%%%%

\index{set!operations}
We can obtain new sets by performing operations on other sets.  
In this section we discuss the common set operations.
{\em Venn diagrams}\index{Venn diagram} 
are often used as a pictorial representation of the relationships
between sets.  We provide Venn diagrams to help visualize the set operations.
In our Venn diagrams, the region(s) in the darker color represent the elements
of the set of interest.


% Had to do this formatting by hand to get it to work.
% Not happy about it, but what can you do?
\begin{df}
\rule{0in}{0in}

\noindent
\begin{minipage}{.62\textwidth} 
The {\bf union} of two sets  $A$ and $B$ is
the set containing elements from either $A$
or $B$.  More formally, 
$$A\cup B = \{x:x\in A \ {\mathrm{or}}\ x\in B\}.$$ 
\index{set!union}
\index{union, set}
\index{a sets opers basic@$\cup$ (union)}
Notice that in this case the {\bf or} is an  {\bf inclusive or}.  
That is, $x$ can be in $A$, or it 
can be in $B$, or it can be in both.
\end{minipage}
\hfill
\begin{minipage}[t]{.2\textwidth}
\psset{unit=1pc} 
\begin{pspicture}(-3, -1)(3,4)
\pscustom[fillstyle=solid,fillcolor=white]{\psline(-3.5,2.5)(-3.5,-2.5)(3.5,-2.5)(3.5,2.5)(-3.5,2.5)}
\pscircle[fillstyle=solid,fillcolor=lightgray](-1,0){2}
\pscircle[fillstyle=solid,fillcolor=lightgray](1,0){2} 
\pscustom[fillstyle=solid,fillcolor=lightgray]{\psarc(1,0){2}{120}{240}\psarc(-1,0){2}{300}{60}}
\uput[d](-2,.5){$A$} 
\uput[d](2,.5){$B$}
\uput[l](-3.5,0){$A\cup B$}
\end{pspicture}
%\meinecaption{1}{$A\cup B$} 
%\label{fig:a_union_b}
%\end{minipage}
\end{minipage}
\end{df}


\begin{exa}
Let $A = \{1, 2, 3, 4, 5, 6\}$, and $B = \{1, 3, 5, 7\}$. Then
$A\cup B = \{1, 2, 3, 4, 5, 6, 7\}$.
\end{exa}

\begin{exer}
Let $A$ be the set of even integers and $B$ be the set of odd integers.
Then $A\cup B$= \blanklinefill
\begin{answer}
$\BBZ$, or the set of (all) integers.
\end{answer} 
\end{exer}

\begin{df}
\rule{0in}{0in}

\noindent
\begin{minipage}{.62\textwidth}
The {\bf intersection} of two sets $A$
and $B$ is the set containing elements that are
in  both $A$ and $B$.  More formally, 
\index{set!intersection}
\index{intersection, set}
\index{a sets opers basic@$\cap$ (intersection)}
$$A\cap B = \{x:x\in A \ {\mathrm{and}}\ x\in B\}.$$
\end{minipage}
\hfill
\begin{minipage}[t]{.2\textwidth}
\psset{unit=1pc} 
\begin{pspicture}(-3, -.5)(3,3)
\pscustom[fillstyle=solid,fillcolor=white]{\psline(-3.5,2.5)(-3.5,-2.5)(3.5,-2.5)(3.5,2.5)(-3.5,2.5)}
\pscircle[fillstyle=solid,fillcolor=white](-1,0){2}
\pscircle[fillstyle=solid,fillcolor=white](1,0){2}
\uput[d](-2,.5){$A$} 
\uput[d](2,.5){$B$}
\uput[l](-3.5,0){$A\cap B$}
\pscustom[fillstyle=solid,fillcolor=lightgray]{
\psarc(1,0){2}{120}{240}
\psarc(-1,0){2}{300}{60}
}
\end{pspicture}
\end{minipage}
\end{df}

\begin{exa}
Let $A = \{1, 2, 3, 4, 5, 6\}$, and $B = \{1, 3, 5, 7, 9\}$. Then
$A\cap B = \{1, 3, 5\}$.
\end{exa}

\begin{exer}
Let $A$ be the set of even integers and $B$ be the set of odd integers.
Then $A\cap B$= \blanklinefill
\begin{answer}
$\varnothing$.
\end{answer} 
\end{exer}

\newpage
\begin{df}
\rule{0in}{0in}

\noindent
\begin{minipage}{.62\textwidth}
The {\bf difference} (or {\bf set-difference})
\index{difference, set}
\index{set!difference}
\index{a sets opers basic@$\setminus$ (set-difference)}
of sets $A$ and $B$
is the set containing elements from $A$ that are not in $B$.  More formally, 
$$A\setminus B = \{x:x\in A \ {\mathrm{and}}\ x\not\in B\}.$$
The set difference of $A$ and $B$ is sometimes denoted by $A-B$.
\end{minipage}
\hfill
\begin{minipage}[t]{.2\textwidth}
\psset{unit=1pc}
\begin{pspicture}(-3, -1)(3,4)
\pscustom[fillstyle=solid,fillcolor=white]{\psline(-3.5,2.5)(-3.5,-2.5)(3.5,-2.5)(3.5,2.5)(-3.5,2.5)}
\pscircle[fillstyle=solid,fillcolor=lightgray](-1,0){2}
\pscircle[fillstyle=solid,fillcolor=white](1,0){2}
\pscustom[fillstyle=solid,fillcolor=white]{\psarc(1,0){2}{120}{240}\psarc(-1,0){2}{300}{60}}
\uput[d](-2,.5){$A$} 
\uput[d](2,.5){$B$}
\uput[l](-3.5,0){$A\setminus B$}
\end{pspicture}
\end{minipage}
\end{df}

\begin{exa}
Let $A = \{1, 2, 3, 4, 5, 6\}$, and $B = \{1, 3, 5, 7, 9\}$. Then  
$A\setminus B = \{2, 4, 6\}$ and
$B\setminus A = \{7, 9\}$.
\end{exa}


\begin{exer}
Let $A$ be the set of even integers and $B$ be the set of odd integers.
Then $A\setminus B$= \blankline{2} and $B\setminus A$= \blankline{2}.
\begin{answer}
$A$; $B$.
\end{answer} 
\end{exer}


We can now prove Theorem~\ref{thm:powerSetSize}.
\begin{exa}
Let $A$ be a set with $n$ elements.  Then $|P(A)|=2^n$. 
\begin{pf}
We use induction\footnote{We will cover induction more fully and formally later.
But since this use of induction is pretty intuitive, 
especially in light of Example~\ref{exa:subsets4element}, it serves as a useful foreshadowing of things to come.}
 and the idea from the solution to Exercise~\ref{exa:subsets4element}. 
 Clearly if $|A|= 1$, $A$
has $2^1 = 2$ subsets: $\varnothing$ and $A$ itself.

Assume every
set with $n - 1$ elements has $2^{n-1}$ subsets. Let $A$ be a set
with $n$ elements.  Choose some $x\in A$.  Every subset of $A$
either contains $x$ or it doesn't.  
Those that do not contain $x$  are subsets of $A\setminus \{x\}$.  Since
$A\setminus \{x\}$ has $n - 1$ elements, the induction hypothesis implies that it has
$2^{n-1}$ subsets.
Every subset that does contain $x$ corresponds to one of the subsets of $A\setminus \{x\}$ with the element $x$ added.  That is, for each subset $S \subseteq A\setminus \{x\}$,  $S\cup \{x\}$ is a subset of $A$ containing $x$.
Clearly there are $2^{n-1}$ such new subsets.
Since this accounts for all subsets of $A$, $A$ has $2^{n-1} + 2^{n-1} = 2^n$ subsets.  
\end{pf}
\end{exa}

\begin{df}
\rule{0in}{0in}

\noindent
\begin{minipage}{.62\textwidth}
\index{set!complement}\index{complement, set}
\index{a sets opers basic@$\overline{A}$ (complement of $A$)}
Let $A \subseteq U$. The {\bf complement} of  $A$ with respect
to $U$ is just the set difference $U\setminus A$.  More formally, 
$$\overline{A} = \{x\in U: x \not\in A\} =
U\setminus A.$$
In words, $\overline{A}$ is the set of everything not in $A$.
Other common notations for set complement include $A^c$ and $A^\prime$.
\end{minipage}
\hfill
\begin{minipage}[t]{.25\textwidth}
\psset{unit=1pc}
\begin{pspicture}(-3, -.25)(3,4.1)
\pscustom[fillstyle=solid,fillcolor=lightgray]{\psline(-3.5,2.5)(-3.5,-2.5)(3.5,-2.5)(3.5,2.5)(-3.5,2.5)}
\uput[l](-1.75,1.5){$\overline{A}$}
\pscircle[fillstyle=solid,fillcolor=white](0,0){2} 
\uput[u](0,0){$A$}
\uput[r](4,2){$U$}
\end{pspicture}
%\meinecaption{1}{$\overline{A}$} 
%\label{fig:a_complement}
\end{minipage}
\end{df}

\begin{rem}
\index{universe}
\index{universal set}
\index{set!universe}
Often the set $U$, which is called the {\bf universe} or {\bf universal set}, 
is implied and we just use $\overline{A}$ to denote the complement.  
We usually follow this convention here.  
Further, when talking about several sets, we will usually assume they have the
same universal set.
\end{rem}


\begin{exa}
Let $U = \{0,1,2,3,4,5,6,7,8,9\}$ be the universal set of decimal digits and
$A = \{0,2,4,6,8\}\subset U$ be the set of even digits. Then
$\overline{A} = \{1,3,5,7,9\}$ is the set of odd digits.
\end{exa}


\begin{exer}
Let $A$ be the set of even integers and $B$ be the set of odd integers, and 
let the universal set be $U=\BBZ$. 
Then $\overline{A}$= \blankline{1.5} and $\overline{B}$= \blankline{1.5}.
\begin{answer}
$B$; $A$.
\end{answer} 
\end{exer}

It should not be too difficult to convince yourself that the following theorem
is true.

\begin{thm}
Let $A$ be a subset of some universal set $U$.  Then 
\begin{eqnarray*}
\overline{A}\cap A &=& \varnothing,\mbox{ and}\\
\overline{A}\cup A &=& U.
\end{eqnarray*}
\end{thm}




 The various intersecting regions for two and three sets can be
seen in Figures \ref{fig:two_sets} and \ref{fig:three_sets}.

%I changed a bunch of stuff here to get these to position correctly.
%I have no idea what I may have messed up. Since I don't know what this is doing.
%CAC, 7/26/13
\begin{figure}[tbh]
\centering
\begin{minipage}{8cm}
 \psset{unit=2.2pc}
\begin{pspicture}(-4, 0)(3,3)
\pscustom[fillstyle=solid,fillcolor=white]{\psline(-3.5,2.5)(-3.5,-2.5)(3.5,-2.5)(3.5,2.5)(-3.5,2.5)}
 \pscustom[fillstyle=solid,fillcolor=blue]{\psarc(1,0){2}{120}{240}\psarcn(-1,0){2}{300}{60}} %AB'
\pscustom[fillstyle=solid,fillcolor=yellow]{\psarc(-1,0){2}{300}{60}\psarcn(1,0){2}{120}{240}}%A'B
\pscustom[fillstyle=solid,fillcolor=white]{\psarc(1,0){2}{120}{240}\psarc(-1,0){2}{300}{60} }%AB
{\footnotesize
 \uput[l](0.6,0){ $A\cap B$} 
\uput[l](-1.25,0.5){\textcolor{white}{$A\cap\overline{B}$}}
\uput[r](1.25,0.5){$\overline{A}\cap B$}
\uput[ur](1.85,1.65){$\overline{(A\cup B)}$}
\uput[d](-2.5,-1.5){$A$} 
\uput[d](2.5,-1.5){$B$}
}
\end{pspicture}
\meinecaption{2}{Venn diagram for two sets.}\label{fig:two_sets}

\end{minipage}
\hfill
\begin{minipage}{8cm}
 \psset{unit=2.2pc}
\begin{pspicture}(-4, 0)(4,3.5)
\pscustom[fillstyle=solid,fillcolor=white]{\psline(-3.5,4)(-3.5,-2.5)(3.5,-2.5)(3.5,4)(-3.5,4)}
 \pscustom[fillstyle=solid,fillcolor=blue]{\psarc(1,0){2}{180}{240}
\psarcn(-1,0){2}{300}{120} \psarc(0,1.7320508076){2}{180}{240}
 }
 %AB'C'
\pscustom[fillstyle=solid,fillcolor=yellow]{\psarc(1,0){2}{240}{60}
\psarcn(0,1.7320508076){2}{0}{300} \psarc(-1,0){2}{300}{360}}%A'BC'
\pscustom[fillstyle=solid,fillcolor=green]{
\psarc(1,0){2}{180}{240}\psarc(-1,0){2}{300}{360}
 \psarcn(0,1.7320508076){2}{300}{240}}%ABC'
 \pscustom[fillstyle=solid,fillcolor=white]{\psarc(0,1.7320508076){2}{240}{300}\psarc(-1,0){2}{0}{60}
 \psarc(1,0){2}{120}{180}}%ABC
\pscustom[fillstyle=solid,fillcolor=violet]{\psarcn(1,0){2}{180}{120}
\psarc(-1,0){2}{60}{120} \psarc(0,1.7320508076){2}{180}{240} }%AB'C
 \pscustom[fillstyle=solid,fillcolor=red]{\psarc(-1,0){2}{60}{120}
 \psarcn(0,1.7320508076){2}{180}{0}\psarc(1,0){2}{60}{120} }%A'B'C
 \pscustom[fillstyle=solid,fillcolor=orange]{\psarcn(-1,0){2}{60}{0}
 \psarc(0,1.7320508076){2}{300}{360}\psarc(1,0){2}{60}{120}}%A'BC
{\footnotesize
\uput[u](.1,2.5){\textcolor{white}{$\overline{A}\cap \overline{B}\cap C$}} 
\uput[d](0,.75){$A\cap B\cap C$} 
\uput[r](.05,1.65){ $\overline{A}\cap B\cap C$} 
\uput[l](-.2,1.7){\textcolor{white}{$A\cap \overline{B}\cap C$}} 
\uput[dl](-1.1,0){\textcolor{white}{$A\cap \overline{B}\cap \overline{C}$}}
\uput[d](0,-.25){$A\cap B \cap\overline{C}$}
\uput[dr](1.1,0){$\overline{A}\cap B \cap \overline{C}$} 
\uput[ur](1.3,3.1){$\overline{(A\cup B \cup C)}$}
} 
\uput[d](-2.5,-1.5){$A$} 
\uput[d](2.5,-1.5){$B$}
\uput[u](-1.75,3){$C$}
\end{pspicture}\meinecaption{2}{Venn diagram for three sets.}\label{fig:three_sets}
\end{minipage}
\end{figure}

\begin{df}
\index{set!disjoint}\index{set!mutually exclusive}\index{disjoint, set}
Two sets  $A$ and $B$ are {\bf disjoint} or {\bf mutually exclusive}
if $A\cap B = \varnothing$.  That is, they have no elements in common.
\end{df}

\begin{absolutelynopagebreak}
\begin{exa}
Let $A$ be the set of prime numbers, $B$ be the set of perfect squares, and
$C$ be the set of even numbers.
Then $A$ and $B$ are clearly disjoint since if a number is a perfect square,
it cannot possibly be prime (although $0$ and $1$ are not prime for different
reasons than the rest of the elements of $B$).
On the other hand, $A$ and $C$ are not disjoint since they both
contain $2$, and $B$ and $C$ are not disjoint because they both contain $4$.
\end{exa}
\end{absolutelynopagebreak}

\begin{exer}
Let $A$ be the set of even integers and $B$ be the set of odd integers.
Are $A$ and $B$ disjoint?  Explain.

\noindent
\answerlinetwo
\begin{answer}
Since no integer is both even and odd, $A$ and $B$ are disjoint.
\end{answer} 
\end{exer}

{\em Set identities} can be used to show that two sets are the same.
Table~\ref{table:setIdentities} gives some of the most common set identities.
In these identities, $U$ is the universal set.  We won't provide proofs for most of these,
but we will present a few examples and a technique that will allow you to verify that they 
are correct in Section~\ref{section:setProofs}.
 
\begin{table}[bht]
\centering
$
\begin{array}{|l|l|}
\hline
Name&Identity\\
\hline
\hline
commutativity	& A\cup B=B\cup A\\
				& A\cap B = B\cap A\\
\hline
associativity	&A\cup(B\cup C)=(A\cup B)\cup C\\
				&A\cap(B\cap C)=(A\cap B)\cap C\\
\hline
distributive&	A\cap(B\cup C) = (A\cap B) \cup (A\cap C)\\
			&	A\cup(B\cap C) = (A\cup B)\cap(A\cup C)\\
\hline
identity&		A \cup \varnothing  = A\\
			&		A\cap U = A\\
\hline
complement&A\cup \overline{A} = U\rule{0in}{.17in}\\ % rule gives enough space on top o
		& A\cap\overline{A} = \varnothing\rule{0in}{.17in} \\
\hline
domination&	  	 A\cup U \ = U\\
			&	A\cap \varnothing  = \varnothing \\
\hline
idempotent&		A\cup A=A\\
			&	 A\cap A = A\\
\hline
complementation	& \overline{(\overline{A})}=A	\rule{0in}{.2in}\\
\hline
De Morgan's	&	 \overline{A \cup B} = \overline{A}\cap\overline{B}\rule{0in}{.17in}\\
			&	\overline{A\cap B} = \overline{A}\cup\overline{B}\rule{0in}{.17in}\\
\hline
absorption & A\cup(A\cap B) = A\\
& A\cap(A\cup B) = A\\
\hline
\end{array}
$
\caption{Set Identities}\label{table:setIdentities}
\end{table}

These identities may look somewhat familiar.
They are
essentially the same as the logical equivalences presented in Table~\ref{table:logicLaws}.
In fact, if we equate $T$ to $U$, $F$ to $\varnothing$, $\vee$ to $\cup$, 
$\wedge$ to $\cap$, and $\neg$ to $\bar{ }$ (complement), 
the laws are identical.
This is because logic operations and sets are both what we call {\em Boolean algebras}.
We won't go into detail about this connection, but in case you run into the concept
in the future, you heard it here first!



Sometimes you need to find the number of elements in the union of several sets.
This is easy if the sets do not intersect.  If they do intersect, 
more care is needed to make sure no elements are missed or counted more than once.
In the following examples we will use Venn diagrams to help us do this correctly.
Later, we will learn about a more powerful tool to do this---{\em inclusion-exclusion}.


\begin{exa}
Of  $40$ people, $28$ smoke and  $16$ chew tobacco. It is also known
that  $10$ both smoke and chew. How many among the $40$ neither
smoke nor chew? 
% Getting multiply defined error.  Is it even used?
%\label{exa:incl_excl_1}
\begin{sol}  
We fill up the Venn diagram below as follows.
Since $|Smoke \cap Chew| = 10$, we put a $10$ in the intersection.
Then we put $28 - 10 = 18$ in the part that $Smoke$ does not overlap
$Chew$ and $16 - 10 = 6$ in the part of $Chew$ that does not overlap
$Smoke$. We have accounted for $10 + 18 + 6 = 34$ people that are in at
least one of the sets. The remaining $40 - 34 = 6$ people outside these
sets don't smoke or chew (and probably don't date girls who do).
\begin{pspicture}(-7, -1.5)(3,2)
\psset{unit=1.5pc} 
\pscustom[fillstyle=solid,fillcolor=white]{\psline(-3.5,2.5)(-3.5,-2.5)(3.5,-2.5)(3.5,2.5)(-3.5,2.5)}
\pscircle(-1,0){2} 
\pscircle(1,0){2}
\rput(0,0){\small 10} 
\rput(-2, 0){\small 18} 
\rput(2, 0){\small 6}
\rput(2.5, 2){\small 6} 
\rput(-2.5, -2){\small Smoke} 
\rput(2.5, -2){\small Chew} 
\end{pspicture}

We truly hope that these numbers are not representative of the number
of people who smoke and/or chew in real life.  It's bad for you.   Don't do it.  Really.
\end{sol}
\end{exa} 

% \vspace*{2cm}
%\begin{figure}[ht]
%\centering
%\begin{minipage}{7cm}
%
%\meinecaption{1}{Example
%\ref{exa:incl_excl_1}.}\label{fig:incl_excl_1}
%\end{minipage}
%\begin{minipage}{7cm}
% \meinecaption{1}{Example
%\ref{exa:3set_incl_excl}.} \label{fig:3set_incl_excl_exa}
%\end{minipage}
%\end{figure}

\begin{exer}
In a group of $30$ people, $8$ speak English, $12$ speak Spanish and
$10$ speak French. It is known that $5$ speak English and Spanish,
$7$ Spanish and French, and $5$ English and French. The number of
people speaking all three languages is $3$. How many people speak
at least one of these languages?

% Some ugly formatting tricks going on here.
\begin{pspicture}(-2, -2)(3,3)
\psset{unit=1.5pc} 
\pscustom[fillstyle=solid,fillcolor=white]{\psline(-4,3.75)(-4,-3)(4,-3)(4,3.75)(-4,3.75)}
\pscircle(-1,0){2} 
\pscircle(1,0){2}
\pscircle(0,1.41421356){2} 
\end{pspicture}
\begin{answer}  
Let $E$ be the set
of all English speakers, $S$ the set of Spanish speakers and $F$ the set of

\noindent
\begin{tabular}{p{.55\textwidth}p{.38\textwidth}}
French speakers in our group.
We fill-up the  Venn diagram (to the right) successively. In the
intersection of all three we put 3. In the region common to $E$ and
$S$ which is not filled up we put $5 - 3 = 2$. In the region common
to $E$ and $F$ which is not already filled up we put $5 - 3 = 2$. In
the region common to $S$ and $F$ which is not already filled up, we
put $7 - 3 = 4.$ In the remaining part of $E$ we put $8 - 2 - 3 - 2
= 1,$ in the remaining part of $S$ we put $12 - 4 - 3 - 2 = 3$, and
in the remaining part of $F$ we put $10 - 2 - 3 - 4 = 1$. Therefore, 
 $1 + 2 + 3 + 4 + 1 + 2 + 3 = 16$ people speak at least one of these languages.
&
\begin{pspicture}(-3, 3)(3,2.5)
\psset{unit=1.5pc} 
\pscustom[fillstyle=solid,fillcolor=white]{\psline(-4.25,4.25)(-4.25,-2.5)(4.25,-2.5)(4.25,4.25)(-4.25,4.25)}
\pscircle(-1,0){2} 
\pscircle(1,0){2}
\pscircle(0,1.41421356){2} 
\rput(2,0){\small 3} 
\rput(-2, 0){\small 1}
\rput(0,0.707106781){\small 3} 
\rput(0, 2.3){\small 1}
\rput(0,-0.9){\small 2} 
\rput(-1.43,1.43){\small 2}
\rput(1.43,1.43){\small 4} 
\rput(-3.8, 0){\small E} 
\rput(3.8,0){\small S} 
\rput(0, 3.8){\small F}
\end{pspicture}
\end{tabular}
\end{answer}
\end{exer} 




\begin{df}
The {\bf Cartesian product}\index{Cartesian product}
\index{a sets opers basicz@$\times$ (Cartesian product)}
 of sets $A$ and $B$ is the set $A\times B=\{(a,b) | a\in A \wedge b\in B\}$.
In other words, it is the set of all ordered pairs of elements from $A$ and $B$.
\end{df}


\begin{exa}
If $A=\{1,2,3\}$ and $B=\{a,b\}$, then\\
$A\times B=\{(1,a), (1,b), (2,a), (2,b), (3,a), (3,b)\},\mbox{ and}$\\
$B\times A=\{(a,1), (a,2), (a,3), (b,1), (b,2), (b,3)\}.$\\
Notice that $A\times B \neq B\times A$.
If $A\neq B$, this is always the case.
\end{exa}

\begin{exer}
Let $A=\{1,2,3,4\}$, and $B=\{3\}$.  Compute $A\times B$.
$$A\times B=\blankline{4}$$
\begin{answer}
$A\times B = \{(1,3), (2,3), (3,3), (4,3)\}$.
\end{answer}
\end{exer}

\begin{df}
If $A$ is a set, then $A^2 = A\times A$, and $A^n = A\times A^{n-1}$.
\end{df}

\begin{exa}
If $B=\{a,b\}$ then 
$$B^2=\{(a,a), (a,b), (b,a), (b,b)\},\mbox{ and}$$
$$B^3=\{(a,a,a), (a,b,a), (b,a,a), (b,b,a), (a,a,b), (a,b,b), (b,a,b), (b,b,b)\}$$
\end{exa}

\begin{exer}
Let $A=\{0,1\}$.  Find $A^2$ and $A^3$.\\[.25in]
$A^2=$\\[.25in]
$A^3=$\\
\begin{answer}
$A^2=\{(0,0), (0,1), (1,0), (1,1)\}$.

\noindent
$A^3=\{(0,0,0), (0,1,0), (1,0,0), (1,1,0), (0,0,1), (0,1,1), (1,0,1), (1,1,1)\}$.
\end{answer}
\end{exer}

It shouldn't be too difficult to convince yourself of the following.
\begin{thm}
If $A$ and $B$ are finite sets with $|A|=n$ and $|B|=m$, then 
$|A\times B|=n\cdot m$.
\end{thm}

\begin{exa}
Let $A$ and $B$ be finite sets with $|A|=100$ and $|B|=5$.
Then $|A\times B| = 100*5=500$, $|A^2|=100*100=10,000$, and $|B^4|=5^4=625$.
\end{exa}

\begin{exer}
Let $A$, $B$, and $C$ be sets with $|A|=10$, $|B|=50$, and $|C|=20$.
Determine the following
\begin{enumerate}[(a)]
\item $|A\times B| = \blankline{2}$
\item $|A\times C| = \blankline{2}$
%\item $|A^2| = \blankline{2}$
\item $|B^3| = \blankline{2}$
\item $|A\times B\times C| = \blankline{2}$
\end{enumerate}
\begin{answer}
(a) $10*50=500$
(b) $10*20=200$
%(c) $10*10=100$
(c) $50*50*50=125,000$
(d) $10*50*20=10,000$
\end{answer}
\end{exer}

\begin{eval}
If $A\times B=\varnothing$, what can we conclude about $A$ and $B$?
\begin{ssolenum}
\item
Assume $A$ and $B$ are not empty. We know the Cartesian product of $A$ and $B$, denoted by $A\times B$, 
is the set of all ordered pairs $(a,b)$, where $a\in A$ and $b\in B$.  Therefore, we can conclude that our
assumption was incorrect because if each set is not empty, $(a,b)$ is in the cross product, but 
$A\times B=\varnothing$, so at least one of the sets must be empty.

\noindent
\evallinetwo
\item
Notice that if $A=\varnothing$ and $B=\varnothing$, $A\times B = \varnothing$.
Therefore, if $A\times B = \varnothing$, then $A=\varnothing$ and $B=\varnothing$.

\noindent
\evallinetwo
\item
We can conclude that both $A$ and $B$ are empty.  I'll prove it by contradiction.
Assume that $A\times B=\varnothing$, but that it is not the case that both $A$ and $B$ are empty. 
Then neither $A$ nor $B$  is empty. 
But then there is some $a\in A$ and some $b\in B$, and $(a,b)\in A\times B$, which 
implies that $A\times B\neq\varnothing$.  This contradicts our assumption.  Therefore
both $A$ and $B$ are empty.

\noindent
\evallinetwo
\item
At least one of $A$ or $B$ is empty by contradiction.
Assume that $A\times B=\varnothing$, but that it is not the case that at least one of $A$ or $B$ is empty.  
Then neither $A$ nor $B$ is empty.
Then there is some $a\in A$ and some $b\in B$.  But then $(a,b)\in A\times B$, which 
implies that $A\times B\neq\varnothing$.  This contradicts our assumption.  Therefore
at least one of $A$ or $B$ is empty.

\noindent
\evallinetwo
\end{ssolenum}
\begin{answer}
\begin{solevalenum}
\item
Although it is on the right track, this solution has several problems.  First, it would be better to make
it more clear that the assumption is that {\em both} $A$ and $B$ are not empty.  
But the bigger problem is the statement `$(a,b)$ is in the cross product'.  
The problem is that $a$ and $b$ are not defined anywhere.
Saying `where $a\in A$ and $b\in B$' earlier does not guarantee that there is such an $a$ or $b$.  
The proof needs to say something along the lines of 
`Since $A$ and $B$ are not empty, then there exist some $a\in A$ and $b\in B$.  Therefore
$(a,b)\in A\times B$\ldots'
\item
This one is way off.  The proof is essentially saying 
`Notice that $p\rightarrow q$. Therefore $q\rightarrow p$.'  But these are not equivalent
statements.  Although it is true that if both $A$ and $B$ are the empty set, then $A\times B$ is
also the empty set, this does not prove that both $A$ and $B$ must be empty in order for $A\times B$
to be empty.  In fact, this isn't the correct conclusion.
\item
The conclusion is incorrect, as is the proof.
The problem is that the negation of
`both $A$ and $B$ are empty' is `it is not the case that both $A$ and $B$ are empty' 
or `at least one of $A$ or $B$ is not empty,' which is not the same thing as
`neither $A$ nor $B$  is empty.'  So although the proof seems to be correct, it is not.
The reason it seems almost correct is that except for this error, the rest of the proof follows
proper proof techniques.  Unfortunately, all it takes is one error to make a proof invalid.
\item
This is a correct conclusion and proof.
\end{solevalenum}
\end{answer}
\end{eval}

\newpage
\subsection{Set Proofs}\label{section:setProofs}


The following theorem can be used to prove set identities.

\begin{thm}\label{sets:containtmentproof}
Two sets $A$ and $B$ are equal if and only if  $A \subseteq
B$ and $B \subseteq A.$
\end{thm}

Let's see this theorem in action.

\begin{exa}\label{setpf1}
Prove that $A\setminus B=A\cap\overline{B}$.
\begin{pf}
Let $x\in A\setminus B$.  Then by definition of difference, 
$x\in A$ and $x\not\in B$.  
But if $x\not\in B$, then $x\in\overline{B}$ by definition of complement.
Since $x\in A$ and $x\in \overline{B}$, $x\in A\cap\overline{B}$ by 
definition of intersection.  Since whenever $x \in A\setminus B$, $x\in A\cap\overline{B}$,
we have shown that  $A\setminus B\subseteq A\cap\overline{B}$.

Now assume that $x\in A\cap\overline{B}$.  Then $x\in A$ and $x\in \overline{B}$ by
definition of intersection.  By definition of complement, $x\not\in B$.
But if $x\in A$ and $x\not\in B$, then $x\in A\setminus B$ by definition of
difference.  Since whenever $x\in A\cap\overline{B}$, $x\in A\setminus B$,
we have that  $A\cap\overline{B}\subseteq A\setminus B$.

Since we have shown that $A\setminus B\subseteq A\cap\overline{B}$ and that
$A\cap\overline{B}\subseteq A\setminus B$, by Theorem~\ref{sets:containtmentproof}
$A\setminus B=A\cap\overline{B}$.
\end{pf}

That was the long, drawn-out version of the proof.  The purpose of all of the
detail is to make the technique clear.  Here is a proof without any
extraneous details.
\begin{pf}
We will prove this by showing set containment both ways.

Let $x\in A\setminus B$.
Then $x\in A$ and $x\not\in B$.  This implies that $x\in\overline{B}$.
Therefore $x\in A\cap\overline{B}$.  
Since $A\setminus B$ implies $x\in A\cap\overline{B}$,
$A\setminus B\subseteq A\cap\overline{B}$.

Now assume that $x\in A\cap\overline{B}$.  Then $x\in A$ and $x\in \overline{B}$.  
Then $x\not\in B$, and therefore $x\in A\setminus B$.
Since $x\in A\cap\overline{B}$ implies $x\in A\setminus B$,
$A\cap\overline{B}\subseteq A\setminus B$.
\end{pf}
\end{exa}


The proofs in the previous example are called {\em set containment proofs}\index{set!containment proof} since we showed set containment both ways.  
The technique is pretty straightforward:  
Theorem \ref{sets:containtmentproof} tells us that if $X\subseteq Y$ and $ Y\subseteq X$, then
$X=Y$.  Thus, to prove $X=Y$, we just need to show that  $X\subseteq Y$ and $ Y\subseteq X$.
But how do we show that one set is a subset of another? 
This is easy: To show that $X\subseteq Y$, we show that every element from $X$ is also in $Y$.  
In other words, we assume that $x\in X$ and use definitions and logic to show that $x\in Y$.
Assuming we do not use any special properties about $x$ other than the fact that $x\in X$, 
then $x$ is an arbitrary element from $X$, so this shows that $X\subseteq Y$.
Showing that $Y\subseteq X$ uses exactly the same technique.

\begin{rem}
Be careful.  To prove that $X=Y$, you generally need to prove two things: $X\subseteq Y$ and
$Y\subseteq X$.  Do not forget to do both.  On the other hand, if you are asked to prove that
$X\subseteq Y$, you do not need to (and should not) show that $Y\subseteq X$.
\end{rem}

Let's see another example of this type of proof.  This proof will provide a few more details
than necessary in order to further explain the technique.

\begin{exa}
Prove the first De Morgan's Laws: Given sets $A$ and $B$, 
$\overline{(A\cup B)} = \overline{A}\cap \overline{B}$.
\begin{pf}
Let $x \in \overline{(A\cup B)}$. 
Then $x \not\in A \cup B$ (by definition of complement).
Thus $x \not\in A$ and $x \not\in B$ (by definition of union), which is the same thing as 
$x \in \overline{A}$ and $x \in \overline{B}$ (by definition of complement).
But then we have that $x \in \overline{A} \cap \overline{B}$ (by definition of intersection). 
Notice that $x$ was an arbitrary element from $\overline{(A\cup B)}$, 
and we showed that $x\in \overline{A} \cap \overline{B}$.  Therefore, every element
in $\overline{(A\cup B)}$ is also in $\overline{A} \cap \overline{B}$.  In other words,
$\overline{(A \cup B)} \subseteq \overline{A} \cap \overline{B}.$

Now, let $x \in \overline{A} \cap \overline{ B}.$ Then $x\in \overline{A}$ and
$x\in {\overline B}$. This means that $x \not\in A$ and $x \not\in B$
which is the same as $x \not\in A \cup B.$ But this last statement
asserts that $x \in \overline{(A \cup B)}$. Hence $\overline{A} \cap \overline{B}
\subseteq \overline{(A \cup B)}.$

Since we have shown that the two sets contain each other, they are equal by 
Theorem \ref{sets:containtmentproof}.
\end{pf}
\end{exa}


You have already seen a few correct ways to prove that $A\setminus B = A\cap\overline{B}$.  
Can you spot the problem(s) in the following `proofs' of this?
These proofs use the alternative notation of $A-B$ for set difference.

%\newpage
\begin{eval}
Use a set containment proof to prove that if $A$ and $B$ are sets, then $A-B = A\cap\overline{B}$.
\begin{spfenum}
\item
Assume $x\in \{A-B\}$ so $x\in A$ and $x$ is not $\in B$.  This means $x\in A$ and $\overline{B}$.
Therefore $x\in A\cap \overline{B}$.  Thus $A-B=A\cap\overline{B}$.

\noindent
\evallinethree
\item
$\overline{B}$ is the other part of the universal that does not contain any part of $B$.
$A\cup\overline{B}$ means all intersection part of $A$ and the universal that does not contain any part of $B$.
Therefore it returns all elements that are in $A$ but not in $B$ which are $A-B$.  Thus, $A-B=A\cap\overline{B}$.

\noindent
\evallinethree
\item
To prove that $A-B=A\cap\overline{B}$, first let $x\in A-B$.  By definition of the difference of sets, this means
that $x$ is an element of $A$ that is not in $B$, or in other words, $x\in A$ and $x\not\in B$.
This is the same as $x\in A\cap\overline{B}$, thus proving that $A-B\subseteq A\cap\overline{B}$.

Now let $x\in A\cap\overline{B}$.  This means that $x\in A$ and $x\not\in B$, so it is in $A$, but 
not in $B$, which is what we just proved in the previous statement, thus proving that $A-B = A\cap\overline{B}$.

\noindent
\evallinethree
\end{spfenum}
\begin{answer}
\begin{pfevalenum}
\item
This solution has several problems.
\begin{enumerate}
\item
$x\in \{A-B\}$ means `$x$ is an element of the set containing $A-B$,' not
`$x$ is an element of $A-B$.'  What they meant was `$x\in A-B$.'
\item
At the end of the first sentence, `$x$ is not $\in B$' mixes mathematical
notation and English in a strange way.  
This should be either `$x\not\in B$' or `$x$ is not in $B$.'
\item
In the second sentence, 
the phrase  `$x\in A$ and $\overline{B}$' is a strange mixture of math and English 
that is potentially ambiguous.  It should be rephrased as something like 
`$x\in A$ and $x\in\overline{B}$' or `$x$ is in both $A$ and $\overline{B}$.'
\item
Finally, what has been shown here is that $A-B\subseteq A\cap\overline{B}$.
This is only half of the proof.  They still need to prove that 
$A\cap\overline{B}\subseteq A-B$.
\end{enumerate}
\item 
Overall, this proof is very confusing and unclear.  More specifically,
\begin{enumerate}
\item
This is an attempt at working through what each set is by using the definitions.  
That would be fine except for two things.  
First, they were asked to give a set containment proof. 
Second, the wording of the proof is confusing and hard to follow. 
I do not come away from this with a sense that anything has been proven. 
\item
They are not using the terminology properly. The terms `universe' or `universal set' would be appropriate, but not `universal' on its own (used twice).  Similarly, what does the phrase
`all intersection part' mean?
Also, a set doesn't `return' anything.  A set is just a set. It contains elements, but
it doesn't `do' anything.
\end{enumerate}
\item
This proof contains a lot of correct elements.  In fact, the first half is on the right track.
However, they jumped from $x\in A$ and $x\not\in B$ to $x\in A\cap\overline{B}$.  Between
these statements they should say something like  
`$x\not\in B$ is equivalent to $x\in\overline{B}$' since the latter statement is really needed
before they can conclude that $x\in A\cap\overline{B}$.
Also, it would be better if they had `by the definition of intersection' 
before or after the statement $x\in A\cap\overline{B}$.
Finally, it would help clarify the proof if the end was something like 
'We have shown that whenever $x\in A-B$, $x\in A\cap\overline{B}$.  Thus, $A-B\subseteq A\cap\overline{B}$.'

The second half of the proof starts out well, but has serious flaws.
The statement 'This means that $x\in A$ and $x\not\in B$'
should be justified by the definitions of complement and intersection, and might even involve two steps.
This is the same problem they had in the first half of the proof.
More serious is the statement `which is what we just proved in the previous statement'.  What exactly does 
that mean?  It is unclear how `what we just proved' immediately leads us to the conclusion that 
$A-B = A\cap\overline{B}$.  First we need to establish that $x\in A-B$ based on the previous statements (easy).  
Then we can say that $A\cap\overline{B}\subseteq A-B$.  Finally, we can combine this with the first part
of the proof to say that $A-B = A\cap\overline{B}$.

In summary, the first half is pretty good.  It should at least make the connection between 
$x\not\in B$ and $x\in\overline{B}$.  The other suggestions clarify the proof a little,
but the proof would be O.K. if they were omitted.  The second half is another story.  It doesn't really
prove anything, but instead makes a vague appeal to something that was proven before.  Not only is
what they are referring to unclear, but how the proof of one direction is related to the proof of
the other direction is also unclear.
\end{pfevalenum}
\end{answer}
\end{eval}




Sometimes we can do a set containment proof in one step instead of two.  
This only works if every step of the proof is reversible.  We illustrate this idea next.
%(Here, the $ \Leftrightarrow$ means ``if and only if''.  
%Although it looks a lot like it, it is not the logical biconditional operator.)\\

%\newpage

\begin{exa}\label{setpf2}
Prove that $A\setminus(B \cup C) = (A\setminus B) \cap(A\setminus
C)$.
\begin{pf}
We have $$
\begin{array}{lll} x \in A\setminus(B \cup C)
&  \leftrightarrow &  x \in A \wedge x\not \in (B \cup C)\\  
&  \leftrightarrow & (x \in A) \ \wedge \ ( (x \not \in B) \ \wedge\ (x \not \in C)) \\
&  \leftrightarrow & (x  \in A \ \wedge \ x \not \in B) \ \wedge \ (x  \in A \ \wedge \ x \not \in C)\\
&  \leftrightarrow & (x \in A\setminus B) \ \wedge\ (x \in A \setminus C) \\ 
&  \leftrightarrow & x \in (A\setminus B) \cap(A\setminus C).
 \end{array}
$$
\end{pf}
\end{exa}

\begin{rem}
The proof in the previous example works because every step is reversible.
You can only write something like `$\alpha\leftrightarrow \beta$' in a proof if 
$\alpha\rightarrow\beta$ and $\beta\rightarrow\alpha$ are both true.
When attempting to shortcut proofs with this technique, make sure each step truly is
reversible.
\end{rem}

\begin{fib}\label{setpf3}
Use a  set containment proof to show that 
$$(A \cup B) \cap C = (A \cap C) \cup (B \cap C).$$
\begin{sol} We have, 
$$\begin{array}{lll}
\multicolumn{2}{l}{x \in (A \cup B) \cap C}\\
\leftrightarrow  & x\in (A \cup B) \wedge \blankline{1.5}&\mbox{ by def. of intersection}\\
\leftrightarrow & (x\in A\vee \blankline{1.5}) \wedge x\in C &\mbox{ by \blankline{1.5}}\\
\leftrightarrow & (x\in A\wedge x\in C) \vee \blankline{1.75}&\mbox{ by  \blankline{1.5}}\\
\leftrightarrow & \blankline{2}\vee (x\in B \cap C)&\mbox{ by \blankline{1.5}}\\
\leftrightarrow & x \in (A \cap C) \cup (B \cap C).&\mbox{ by \blankline{1.5}}
\end{array}$$
\end{sol}
\begin{answer}
$x\in C$; 
$x\in B$;
definition of union;
$(x\in B \wedge x\in C)$;
distributive law (the logical one);
$(x\in A\cap C) $;
definition of intersection;
definition of union.
\end{answer}
\end{fib}

\newpage
\section{Modular Arithmetic, GCD, Rounding}\label{section:moddiv}

In this section we introduce some notation
that allows us to think about remainders when doing division 
in a different way than you may be used to. 
Then we will discuss the greatest common divisor of two integers,
including a simple algorithm to compute it.
Finally, we will see the floor and ceiling functions, which allow
us to round down or up, depending on our need.

We begin by repeating a few definitions from Section~\ref{section:dProofs} in case
you skipped it.

\begin{df}\index{even}\index{odd}\index{parity}
Recall that:
\begin{itemize} 
\item An {\bf even integer}\index{even} is one of the form $2k$, where $k$ is an integer.
\item An {\bf odd integer}\index{odd} is one of the form $2k + 1$ where $k$ is an integer.
\item Two integers have the same {\bf parity}\index{parity} if they are both even or both odd.
\end{itemize}
\end{df}

\begin{exa}
Since $14=2\cdot 7$, it is even. 
Similarly, since $23=2\cdot 11 + 1$, it is odd.
\end{exa}
\begin{exa}
Since $50$ and $124$ are both even, they have the same parity.
\end{exa}

\begin{df}\index{divides}\index{factor}\index{multiple}
\index{divisor}\index{divisible}
\index{a fdiv@$\vert$ (divides)}
Let $b$ and $a$ be integers with $a\neq 0$.  We say that 
$b$ is {\bf divisible by} $a$ if there exists an integer $c$ such
that $b = ac$. 
If $b$ is divisible by $a$, 
we also say that $b$ is a {\bf multiple} of $a$,
$a$ is a {\bf factor} or {\bf divisor} of $b$, and that 
$a$ {\bf divides} $b$, written as $a|b$.  If $a$ does not divide $b$, we write $a\nmid b$. 
\end{df}

\begin{exa}
Since $6=2\cdot 3$, $2|6$, and $3|6$.  But $4\nmid 6$ since we cannot write $6=4\cdot c$ for any integer $c$.
\end{exa}

\begin{exa}
$100$ is divisible by $25$.
We can say the same thing by saying $25$ divides $100$, which 
we can also write as $25|100$. We can also say that $25$ is a factor of $100$.
\end{exa}

\begin{df}
\index{prime}
\index{composite}
A positive integer $p>1$ is {\bf prime} if its only positive factors are $1$ and $p$. 
A positive integer $c>1$ which is not prime is said to be
{\bf composite}.
\end{df}

\begin{exa}
Since $21=3\cdot 7$, it is composite and therefore not prime.
On the other hand, $17$ has no factors other than $1$ and $17$ so it is prime.
\end{exa}

\begin{rem}
Notice that according to the definitions given above, $1$ is neither prime nor composite.
This is one of the many things that makes $1$ special.
\end{rem}

\begin{df}
\index{factorial}
\index{a fact@"! (factorial)}
For a non-negative integer $n$, the quantity $n!$ (read {\bf ``$n$ factorial''}) is defined as follows.
$0! = 1$ and if $n>0$ then $n!$ is the product of all the integers from $1$ to $n$ inclusive: $$n! = 1\cdot 2 \cdots n.  $$
\end{df}

\begin{exa}
$3! = 1\cdot 2 \cdot 3 = 6$, and 
$5! = 1\cdot 2 \cdot 3 \cdot 4 \cdot 5 = 120$. 
\end{exa}

\begin{df}
%\index{a mod@$\bmod$ (operator)}
\index{mod}
\index{a mod@\% (modulus)}
The  {\bf mod} operator is defined as follows: for 
integers $a$ and $n$ such that $a\geq 0$ and $n>0$, 
%${\boldmath a \bmod n}$ % Don't use math mode here. It gives a warning.
{\bf\em a mod n}
is the integral non-negative remainder
when $a$ is divided by $n$.  Observe that this remainder is one of
the $n$ numbers $$0,\quad  1,\quad  2,\quad  \ldots , \quad n-1.$$
When we are working with the mod operator, we say we are
are doing {\bf modular arithmetic}.
\end{df}

\begin{exa} 
Here are some example computations: \\[-.35in]
\begin{multicols}{3} \columnsep 25pt\multicoltolerance=900
\noindent
$234 \bmod 100 = 34$\\
$38 \bmod 15 = 8$\\
$15 \bmod 38 = 15$\\
$ 1961 \bmod 37 = 0$\\
$ 1966 \bmod 37 = 5$\\
$ 1\bmod 5 = 1$\\
$6\bmod 5 =1 $\\
$11\bmod 5 = 1$\\
$16\bmod 5 = 1$
\end{multicols}
\end{exa}


\begin{exer} Compute the following: \\[-.45in]
\begin{multicols}{3}
\begin{enumerate}[(a)]
\item
$345 \bmod 100 = \blankline{.5}$
\item
$23 \bmod 15 = \blankline{.5}$
\item
$15 \bmod 4 = \blankline{.5}$
\item
$ 15 \bmod 9 = \blankline{.5}$
\item
$ 27 \bmod 9 = \blankline{.5}$
\item
$ 7\bmod 12 = \blankline{.5}$
\item
$19\bmod 12 = \blankline{.5} $
\item
$31\bmod 12 = \blankline{.5}$
\item
$47\bmod 12 = \blankline{.5}$
\end{enumerate}
\end{multicols}
\begin{answer}
(a) 45; (b) 8; (c) 3; (d) 6; (e) 0; (f) 7; (g) 7; (h) 7; (i) 11.
\end{answer}
\end{exer}

\begin{df}\label{df:congr}
\index{congruence modulo n@congruence modulo $n$}
\index{a congruence@$\equiv$ (congruence modulo $n$)}
For integers $a$, $b$, and $n$, where $n>0$,
we say that $a$ is {\bf congruent to } $b$ {\bf modulo $n$} if
$n$ divides $a-b$ (that is, $a-b=kn$ for some integer $k$). 
We write this as $a\equiv b\pmod n$.

There are a few other (equivalent) ways of defining congruence modulo $n$.
\begin{itemize}
\item
$a\equiv b\pmod n$ iff $a$ and $b$ have the same remainder when divided by $n$.
\item
$a\equiv b \pmod n$ iff $a-b$ is a multiple of $n$.
\item
$a\equiv b \pmod n$ iff $a-b=kn$ for some integer $k$.
\end{itemize}

\noindent
If $a-b\neq kn$ for any integer $k$, then $a$ is not congruent to $b$ modulo  $n$, 
and we write this as $a\not \equiv b \pmod n$.
\end{df}

\begin{exa}
Notice that $21-6 = 15 = 3\cdot 5$, so $21\equiv 6 \pmod 5$.
\end{exa}

Notice that if $a\equiv b\pmod n$ and $0\leq b < n$, then 
$b$ is the remainder when $a$ is divided by $n$.

\begin{exer}\label{exa:squares_mod_4_a}
Prove that for every integer $n$, $n^2\pmod 4$ is either $0$ or $1$.
(Hint: Consider the cases when $n$ is even and odd.)
\vspace*{1.5in}
\begin{answer}
Since every integer is either even (of the form $2k$) or odd (of
the form $2k + 1$) we have two possibilities:
$$\begin{array}{lllll}(2k)^2 & = & 4k^2&\equiv&0\pmod 4, \mbox{or}\\
(2k + 1)^2 & = & 4(k^2 + k) + 1&\equiv&1\pmod 4.
  \end{array}$$
Thus, $n^2$ has remainder $0$ or $1$ when divided by $4$.  
\end{answer}
\end{exer}

\begin{exa}\label{exa:squares_mod_4}
 Prove that the sum of two
squares of integers leaves remainder $0$, $1$ or $2$ when divided
by $4$.
\begin{pf}
According to Example \ref{exa:squares_mod_4_a}, 
the squares of integers  have remainder $0$ or $1$
when divided by $4$.  Thus, when we add two squares, the possible remainders when divided by $4$
are $0$ ($0+0$), $1$ ($0+1$ or $1+0$) , and $2$ ($1+1$).
\end{pf}
\end{exa}

\begin{exa}
Prove that $2003$ is not the sum of two squares.
\begin{pf}
Notice that $2003\equiv 3\pmod 4$.  Thus, by Example~\ref{exa:squares_mod_4} we know that $2003$
cannot be the sum of two squares.
\end{pf}
\end{exa}

 We now prove some simple properties of congruences.
\begin{thm}
Let $a, b, c, d \in \BBZ$, and $n,k\in\BBZ^+$ with $a \equiv b$ mod $n$ and
$c \equiv d \bmod n$. Then
\begin{enumerate} \item $a + c \equiv b + d\bmod n$ \item
$a - c \equiv b - d\bmod n$ \item $ac \equiv bd\bmod n$ \item $a^k
\equiv b^k\bmod n$ \item If $f$ is a polynomial with integral
coefficients then $f(a) \equiv f(b)\bmod n$.
\end{enumerate}
\begin{pf} As $a \equiv b \bmod n$ and $c \equiv d\bmod n$, we can find
$k_1, k_2 \in \BBZ$ with $a = b + k_1n$ and $c = d + k_2n$. Thus
$a \pm c = b \pm d + n(k_1 \pm k_2)$ and $ac = bd + n(k_2b +
k_1d)$. These equalities give (1), (2) and (3). Property (4)
follows by successive application of (3), and (5) follows from
(4).
\end{pf}
\end{thm}
%
%Congruences mod $9$ can sometimes be used to check
%multiplications. For example $875961\cdot 2753 \neq 2410520633.$
%For if this were true then
%$$(8 + 7 + 5 + 9 + 6 + 1)(2 + 7 +
%5 + 3) \equiv 2 + 4 + 1 + 0 + 5 + 2 + 0 + 6 + 3 + 3 \  \bmod \ 9.$$
%But this says that  $0\cdot 8 \equiv 8\bmod 9$, which is patently
%false.

\begin{exa} Find the remainder when $6^{1987}$ is divided by
$37$. 
\begin{sol} \ \\
$6^2 \equiv -1$ mod 37. Thus $6^{1987}
\equiv 6\cdot 6^{1986} \equiv 6(6^2 )^{993} \equiv 6(-1)^{993}
\equiv -6 \equiv 31$ mod $37$.
\end{sol}
\end{exa} 

\begin{exer} 
Prove that $7$ divides $3^{2n + 1} + 2^{n + 2}$ for all natural numbers $n$. 
\vspace*{1.5in}
\begin{answer}
Observe that $3^{2n + 1} \equiv 3\cdot 9^n \equiv 3\cdot
2^n$ mod 7 and $2^{n + 2} \equiv 4\cdot 2^n$ mod 7. Hence
$ 3^{2n + 1} + 2^{n + 2} \equiv 3\cdot 2^n+4\cdot 2^n\equiv 7\cdot 2^n \equiv 0 \  \bmod  \ 7,$
for all
natural numbers $n.$
\end{answer}
\end{exer}

%\begin{exa}Prove that $7|(2222^{5555} + 5555^{2222}).$\end{exa}
%Solution: $2222 \equiv 3 \bmod 7$, $5555 \equiv 4 \bmod 7$ and $3^5
%\equiv 5 \bmod 7$. Now $$2222^{5555} + 5555^{2222} \equiv 3^{5555}
%+ 4^{2222} \equiv (3^5)^{1111} + (4^2)^{1111} \equiv 5^{1111} -
%5^{1111} \equiv 0 \bmod 7.$$

\begin{exa} Find the units digit of $7^{7^7}$.
\begin{sol}We must find $7^{7^7}$ mod 10. Now, $7^2 \equiv - 1$ mod
10, and so $7^3 \equiv 7^2 \cdot 7 \equiv -7 \equiv 3$ mod 10 and
$7^4 \equiv (7^2)^2 \equiv 1$ mod 10. Also, $7^2 \equiv 1$ mod 4
and so $7^7 \equiv (7^2)^3 \cdot 7 \equiv 3$ mod 4, which means
that there is an integer $t$ such that $7^7 = 3 + 4t.$ Upon
assembling all this,
$$ 7^{7^7} \equiv 7^{4t + 3} \equiv (7^4)^t \cdot 7^3 \equiv 1^t\cdot 3
\equiv 3 \  \bmod  \ 10. $$Thus the last digit is 3.
\end{sol}
\end{exa}
\begin{exa}
Prove that every year, including any leap year, has at least one
Friday 13th.
\begin{sol}
It is enough to prove that each year has a Sunday the
1st. Now, the first day of a month in each year falls in one of
the following days:
$$\begin{array}{|l|l|l|}
\hline {\rm Month} & {\rm Day\ of\ the\ year} & \bmod 7 \\
\hline {\rm January} &  1 & 1 \\
\hline {\rm February} &  32  & 4\  \\
\hline {\rm March} & 60 \ {\rm or}\ 61&  4\ {\rm or}\ 5\ \\
\hline {\rm April} & 91 \ {\rm or}\ 92& 0 \  {\rm or}\ 1 \\
\hline {\rm May} & 121\ {\rm or} 122 &  2\  {\rm or}\ 3\\
\hline {\rm June} & 152\ {\rm or}\ 153& 5\ {\rm or}\ 6\\
\hline {\rm July}&  182 \ {\rm or} 183\ &  0\ {\rm or}\ 1\\
\hline {\rm August}&  213\ {\rm or}\ 214 &  3\ {\rm or}\ 4\\
\hline {\rm September}& 244 \ {\rm or}\ 245&  6\ {\rm or}\ 0 \\
\hline {\rm October}&  274\ {\rm or}\ 275 &  1\ {\rm or}\ 2\\
\hline {\rm November}&  305\ {\rm or}\ 306&  4\ {\rm or}\ 5\\
\hline {\rm December}&  335\ {\rm or}\ 336&   6\ {\rm or}\ 0\\
\hline
\end{array}$$
(The above table means that, depending on whether the year is a leap
year or not, that March 1st is the $60$th or $61$st day of the year,
etc.) Now, each remainder class modulo $7$ is represented in the
third column, thus each year, whether leap or not, has at least one
Sunday the 1st.
\end{sol}
\end{exa}
%
%\begin{exa} Find infinitely many integers $n$ such that
%$2^n + 27$ is divisible by $7$.\end{exa} Solution: Observe that
%$2^1 \equiv 2, 2^2 \equiv 4, 2^3 \equiv 1, 2^4 \equiv 2, 2^5
%\equiv 4, 2^6 \equiv 1$ mod 7 and so $2^{3k} \equiv 1$ mod 3 for
%all positive integers $k$. Hence $2^{3k} + 27 \equiv 1 + 27 \equiv
%0 \bmod 7$ for all positive integers $k$. This produces the
%infinitely many values sought.
\begin{exer} 
Prove that for any integer $k>0$, $2^k - 5$ never leaves remainder
$1$ when divided by $7$.
\vspace*{1in}
\begin{answer} $2^1 \equiv 2, 2^2
\equiv 4, 2^3 \equiv 1$ mod 7, and this cycle of three repeats.
Thus $2^k - 5$ can leave only remainders 
$4=(2-5)\bmod 7$, $6=(4-5)\bmod 7$, or $3=(1-5\bmod 7)$ upon
division by $7$, none of which are 1.
\end{answer}
\end{exer}

\noindent
The proof of the following is left as an exercise.
Recall that {\bf iff} is shorthand for {\bf if and only if}.
\begin{thm}\label{thm:congruence}
$a\equiv b\pmod n$ iff
%\footnote{{\bf iff} is shorthand for {\bf if and only if}.}
$a\bmod n=b\bmod n$.
\end{thm}

\begin{exa}
Since $1961\bmod{37}=0$ and  $356\bmod{37}=23$, 
and $0\neq 23$, we know that $1961\not\equiv 356 \pmod{37}$
by Theorem~\ref{thm:congruence}.
\end{exa}


\begin{rem}
Our definition of $\bmod$ requires that $n>0$ and $a\geq 0$.
It is possible to define $a \bmod n$ when $a$ is negative.  
Unfortunately, there are two possible ways of doing so based on how you define the
remainder when the dividend is negative. One possible answer is negative and the other
is positive.   They always differ by $n$, so computing one from the other is easy.
\end{rem}
\begin{exa}
Since $-13 = (-2)*5 - 3$ and $-13 = (-3)*5 + 2$, we might consider the remainder of $-13/5$
as either $-3$ or $2$.  Thus,  $-13\bmod 5=-3$ and $-13\bmod 5 = 2$ both seem like reasonable answers.
Fortunately, the two possible answers differ by $5$.  
In fact, you can always obtain the positive possibility by adding $n$ to the negative possibility.
\end{exa}

\begin{exer}
Fill in the missing numbers that are congruent to $1\pmod 4$ (listed in increasing order)\\
\blankline{.5}, -11, \blankline{.5}, -3, 1, 5, \blankline{.5}, \blankline{.5}, 17, \blankline{.5}
\begin{answer}
-15; -7; 9; 13; 21.  Notice that it is every 4th number along the number line, both in the
positive and negative directions.
\end{answer}
\end{exer}

\begin{df}\index{gcd}\index{greatest common divisor}
Let  $a$, $b$ be integers with one of them different from $0$. The
{\bf greatest common divisor} $d$ of $a, b$, denoted by $d = \gcd(a, b)$
is the largest positive integer that divides both $a$ and $b$.
\end{df}

\begin{exa}
Since $300=2^2\cdot 3\cdot 5^2$ and $70=2\cdot 5\cdot 7$, 
$\gcd(300,70)=2\cdot 5=10.$
\end{exa}

Factoring numbers in order to determine the greatest common divisor is
inefficient. There is a better algorithm to compute it based on the fact (which we will not prove) that $\gcd(a,b)=\gcd(b,a\bmod b)$. 
Since $\gcd(a,b)=\gcd(b,a)$, we can assume that $a>b$ when we apply this rule.
Further, once $a\bmod b=0$, we will know that $\gcd(a,b)=b$.

\begin{exa}
We can compute $\gcd(300,70)$ as follows:
$$
\begin{array}{lll}
\gcd(300,70)
&=&\gcd(70,300\bmod 70)=\gcd(70,20)\\
&=&\gcd(20,70\bmod 20)=\gcd(20,10)\\
&=&\gcd(10,20\bmod 10)=\gcd(10,0)=10.
\end{array}$$
\end{exa}

The following procedure is based on this idea.
We will cover algorithms more completely in
Chapter~\ref{ch:algs}, but for now the main
thing to point out is that \verb+int+ mean integer,
and that a \verb+while+ loop executes repeatedly as long as the condition 
is true.

\begin{proc}\index{Euclid's Algorithm}\label{proc:euclid}
{\bf Euclid's Algorithm}
\begin{code}
// Given integers a and b, return their greatest common divisor
int gcd(int a, int b) {
    while (b!=0) {
        int r = a mod b
        a=b
        b=r
    }
    return a
}
\end{code}
\end{proc}

\begin{exa}
We compute $\gcd(300,70)$ again, this time
explicitly using Euclid's algorithm.

\begin{tabular}{|l|l|l|l|}
\hline
step&a&b&r\\
\hline
\hline
1&300&70&20\\
\hline
2&70&20&10\\
\hline
3&20&10&0\\
\hline
4&10&0&0\\
\hline
\end{tabular}

Since $b=0$, we can stop and we know that $\gcd(300,70)=10$.
\end{exa}

\begin{exer}
Compute $\gcd(524,118)$ using Euclid's algorithm.
\vspace*{1.25in}
\begin{answer}
$\gcd(524,118)=2$ as demonstrated here: 

\begin{tabular}{|l|l|l|l|}
\hline
step&a&b&r\\
\hline
\hline
1&524&118&52\\
\hline
2&118&52&14\\
\hline
3&52&14&10\\
\hline
4&14&10&4\\
\hline
5&10&4&2\\
\hline
6&4&2&0\\
\hline
7&2&0&0\\
\hline
\end{tabular}
\end{answer}
\end{exer}

\begin{df}
Two integers are said to be {\bf relatively prime}
if they have no factors in common. That is, $a$ and $b$ are
relatively prime exactly when $\gcd(a,b)=1$.
\end{df}

\begin{exa}
Since we previously saw that $\gcd(300,70)=10$,
300 and 7 are not relatively prime. 
On the other hand, 125 and 16 are relatively prime since they have no common 
factors ($125=5^3$ and $16=2^4$).
\end{exa}

\begin{exer}
Determine whether or not $867$ and $5309$ are relatively prime by
first using Euclid's algorithm to determine their gcd.
\vspace*{1.25in}
\begin{answer}
You should have computed that $\gcd(867,5309)=1$ and therefore
concluded that they are relatively prime.
\end{answer}
\end{exer}

\begin{df}
\index{floor}
\index{ceiling}
\index{a fb@$\lfloor\, \rfloor$ (floor)}
\index{a fc@$\lceil\, \rceil$ (ceiling)}
The {\bf floor} of a real number $x$, written $\lfloor x\rfloor$, 
is the largest integer that is less than or equal to $x$.
The {\bf ceiling} of a real number $x$, written $\lceil x\rceil$, 
is the smallest integer that is greater than or equal to $x$.
\end{df}
\begin{exa}
$\lfloor 4.5 \rfloor = 4$, 
$\lceil 4.5 \rceil = 5$, 
$\lfloor 7 \rfloor = \lceil 7 \rceil = 7$.\\
In general, if $n$ is an integer, then 
$\lfloor n \rfloor = \lceil n \rceil = n$.
\end{exa}

\begin{exer}
Determine each of the following.\\[-.35in]
\begin{multicols}{3}
\begin{enumerate}
\item $\lfloor 9.9 \rfloor = \blankline{.5}$
\item $\lceil 9.9 \rceil = \blankline{.5}$
\item $\lfloor 9.00001 \rfloor = \blankline{.5}$
\item $\lceil 9.00001 \rceil = \blankline{.5}$
\item $\lfloor 9 \rfloor = \blankline{.5}$
\item $\lceil 9 \rceil = \blankline{.5}$
\end{enumerate}
\end{multicols}
\begin{answer}
(1) 9; (2) 10; (3) 9; (4) 10; (5) 9; (6) 9.
\end{answer}
\end{exer}


The following Theorem and Corollary are somewhat 
obvious, but since floors and ceiling can trip people up,
they are useful to have written down explicitly.

\begin{thm}\label{floorThm}
Let $a$ be an integer and $x$ be a real number.  Then $a\leq x$ if and only if $a\leq \lfloor x\rfloor$.
\begin{pf}
If $a\leq \lfloor x\rfloor$, then $a\leq \lfloor x\rfloor\leq x$ is clear.
On the other hand, assume $a \leq x$.  Then $a$ is an integer that is less than or equal to $x$.
Since $\lfloor x\rfloor$ is the largest integer that is less than or equal to $x$,
$a\leq \lfloor x \rfloor$.
\end{pf}
\end{thm}

\begin{cor}\label{floorCor}
Let $a$, $b$, and $c$ be integers.  Then $a\leq b/c$ if and only if $a\leq \lfloor b/c\rfloor$.
\begin{pf}
Since $b/c$ is a real number, this is a special case of Theorem~\ref{floorThm}.
\end{pf}
\end{cor}





%\input{logicSetsBooleanAlg/setExercises.tex}
%\chapter{Functions and Relations}

\newpage
\section{Functions}\label{section:functions}
This section is meant as a review of what you hopefully already learned in an earlier course, probably in high school.
Thus, it is pretty brief.  But we do try to cover all of the important material and provide enough examples to
illustrate the concepts.

\subsection{Definitions}\label{section:functionDefs}

\begin{df}
Let $A$ and $B$ be sets.  Then a {\bf function} $f$ from $A$ to $B$ assigns to each element of $A$
exactly one element from $B$. We write $f: A\to B$ if $f$ is a function from $A$ to $B$.
If $a\in A$ and $f$ assigns to $a$ the value $b\in B$, we write 
$f(a)=b$.  We also say that $f$ {\bf maps} $a$ to $b$.

If $A=B$, we sometimes say $f$ is a function {\bf on} $A$.
\end{df}

\begin{exa}
If $A= B=\naturals$, we can define a function $f: A\to B$ by $f(x)=x^2$.
Then $f(1)=1$, $f(2)=4$, $f(3)=9$, etc.
Although $f(x)$ is defined for all $x\in A$, not every $b\in B$ is mapped to by $f$.
For instance, there is no $a\in A$ for which $f(a)=5$.
\end{exa}

\begin{exa}
Notice that we can define $f(x)=\sqrt{x}$ on the positive real numbers, but we {\em cannot}
define it on the positive integers since $\sqrt{2}$ is not an integer.
Similarly, since $\sqrt{-1}=i\not\in\reals$, we cannot define it on the real numbers.  
We {\em can} let it be a function from $\reals$ to $\complex$, though.
But we won't because this course is complex enough even without complex numbers.
\end{exa}

\begin{df}
Let $f$ be a function from $A$ to $B$.
\begin{enumerate}
\item
We call $A$ the {\bf domain} of $f$.
\item
We call $B$ the {\bf codomain} of $f$.
\item 
The {\bf range} of $f$ is the set $\{b | f(a)=b \mbox{ for some } a\in A\}$.
In other words the range is the subset of $B$ that are actually mapped to by $f$.
\end{enumerate}
\end{df}

\begin{exa}
Let $A= B=\naturals$ and $f:A\to B$ be defined by $f(x)=x^2$.
Then the domain and codomain of $f$ are both $\naturals$, and the range is 
$\{a^2| a\in \naturals\}$, which is a proper subset of the codomain.
\end{exa}

Figure~\ref{func1} gives a pictorial representation of a function.  Notice that in this example 
every element in $A$ has precisely one arrow going from it. 
So if I ask ``what is $f(x)$?'', there is always an answer and it is always unique.
On the other hand, there is a point in $B$ that has two arrows going to it and several points
that have no arrows going to them.  This is fine.
 
Figure~\ref{funcNot} does not represent a function since there are several points in $A$ which
have two arrows going from them and several with no arrows at all.
The problem here is that if I ask ``what is $f(x)$?'', sometimes there is no answer and sometimes
there are multiple answers.  Thus, $f$ would not represent a function. 

\begin{figure}[ht]
\centering
\begin{minipage}{2in}
\input{fig/functionExample.tex}
\footnotesize\caption{Pictorial representation of a function from $A$ to $B$.} \label{func1}
\end{minipage}
\hspace*{1in}
\begin{minipage}{2in}
%LaTeX with PSTricks extensions
%%Creator: 0.48.3.1
%%Please note this file requires PSTricks extensions
\psset{xunit=.5pt,yunit=.5pt,runit=.5pt}
\begin{pspicture}(269.81277466,210.70080566)
{
\newrgbcolor{curcolor}{0.7019608 0.7019608 0.7019608}
\pscustom[linestyle=none,fillstyle=solid,fillcolor=curcolor]
{
\newpath
\moveto(269.8127758,102.66593772)
\curveto(269.8127758,45.96510421)(244.54333166,-0.00000323)(213.3718926,-0.00000323)
\curveto(182.20045354,-0.00000323)(156.9310094,45.96510421)(156.9310094,102.66593772)
\curveto(156.9310094,159.36677123)(182.20045354,205.33187867)(213.3718926,205.33187867)
\curveto(244.54333166,205.33187867)(269.8127758,159.36677123)(269.8127758,102.66593772)
\closepath
}
}
{
\newrgbcolor{curcolor}{0.7019608 0.7019608 0.7019608}
\pscustom[linestyle=none,fillstyle=solid,fillcolor=curcolor]
{
\newpath
\moveto(112.8817668,106.65361272)
\curveto(112.8817668,49.95277921)(87.61232266,3.98767177)(56.4408836,3.98767177)
\curveto(25.26944454,3.98767177)(0.0000004,49.95277921)(0.0000004,106.65361272)
\curveto(0.0000004,163.35444623)(25.26944454,209.31955367)(56.4408836,209.31955367)
\curveto(87.61232266,209.31955367)(112.8817668,163.35444623)(112.8817668,106.65361272)
\closepath
}
}
{
\newrgbcolor{curcolor}{0 0 0}
\pscustom[linestyle=none,fillstyle=solid,fillcolor=curcolor]
{
\newpath
\moveto(34.7512229,153.17648231)
\curveto(34.7512229,149.5059291)(31.89510862,146.53035781)(28.3719131,146.53035781)
\curveto(24.84871758,146.53035781)(21.9926033,149.5059291)(21.9926033,153.17648231)
\curveto(21.9926033,156.84703552)(24.84871758,159.82260681)(28.3719131,159.82260681)
\curveto(31.89510862,159.82260681)(34.7512229,156.84703552)(34.7512229,153.17648231)
\closepath
}
}
{
\newrgbcolor{curcolor}{0 0 0}
\pscustom[linestyle=none,fillstyle=solid,fillcolor=curcolor]
{
\newpath
\moveto(61.5443199,24.24166731)
\curveto(61.5443199,20.5711141)(58.68820562,17.59554281)(55.1650101,17.59554281)
\curveto(51.64181458,17.59554281)(48.7857003,20.5711141)(48.7857003,24.24166731)
\curveto(48.7857003,27.91222052)(51.64181458,30.88779181)(55.1650101,30.88779181)
\curveto(58.68820562,30.88779181)(61.5443199,27.91222052)(61.5443199,24.24166731)
\closepath
}
}
{
\newrgbcolor{curcolor}{0 0 0}
\pscustom[linestyle=none,fillstyle=solid,fillcolor=curcolor]
{
\newpath
\moveto(78.1305369,80.06911331)
\curveto(78.1305369,76.3985601)(75.27442262,73.42298881)(71.7512271,73.42298881)
\curveto(68.22803158,73.42298881)(65.3719173,76.3985601)(65.3719173,80.06911331)
\curveto(65.3719173,83.73966652)(68.22803158,86.71523781)(71.7512271,86.71523781)
\curveto(75.27442262,86.71523781)(78.1305369,83.73966652)(78.1305369,80.06911331)
\closepath
}
}
{
\newrgbcolor{curcolor}{0 0 0}
\pscustom[linestyle=none,fillstyle=solid,fillcolor=curcolor]
{
\newpath
\moveto(36.0270919,61.45996431)
\curveto(36.0270919,57.7894111)(33.17097762,54.81383981)(29.6477821,54.81383981)
\curveto(26.12458658,54.81383981)(23.2684723,57.7894111)(23.2684723,61.45996431)
\curveto(23.2684723,65.13051752)(26.12458658,68.10608881)(29.6477821,68.10608881)
\curveto(33.17097762,68.10608881)(36.0270919,65.13051752)(36.0270919,61.45996431)
\closepath
}
}
{
\newrgbcolor{curcolor}{0 0 0}
\pscustom[linestyle=none,fillstyle=solid,fillcolor=curcolor]
{
\newpath
\moveto(29.6477809,113.29973531)
\curveto(29.6477809,109.6291821)(26.79166662,106.65361081)(23.2684711,106.65361081)
\curveto(19.74527558,106.65361081)(16.8891613,109.6291821)(16.8891613,113.29973531)
\curveto(16.8891613,116.97028852)(19.74527558,119.94585981)(23.2684711,119.94585981)
\curveto(26.79166662,119.94585981)(29.6477809,116.97028852)(29.6477809,113.29973531)
\closepath
}
}
{
\newrgbcolor{curcolor}{0 0 0}
\pscustom[linestyle=none,fillstyle=solid,fillcolor=curcolor]
{
\newpath
\moveto(76.8546749,139.88423321)
\curveto(76.8546749,136.21368)(73.99856062,133.23810871)(70.4753651,133.23810871)
\curveto(66.95216958,133.23810871)(64.0960553,136.21368)(64.0960553,139.88423321)
\curveto(64.0960553,143.55478642)(66.95216958,146.53035771)(70.4753651,146.53035771)
\curveto(73.99856062,146.53035771)(76.8546749,143.55478642)(76.8546749,139.88423321)
\closepath
}
}
{
\newrgbcolor{curcolor}{0 0 0}
\pscustom[linestyle=none,fillstyle=solid,fillcolor=curcolor]
{
\newpath
\moveto(62.8201919,185.07788031)
\curveto(62.8201919,181.4073271)(59.96407762,178.43175581)(56.4408821,178.43175581)
\curveto(52.91768658,178.43175581)(50.0615723,181.4073271)(50.0615723,185.07788031)
\curveto(50.0615723,188.74843352)(52.91768658,191.72400481)(56.4408821,191.72400481)
\curveto(59.96407762,191.72400481)(62.8201919,188.74843352)(62.8201919,185.07788031)
\closepath
}
}
{
\newrgbcolor{curcolor}{0 0 0}
\pscustom[linestyle=none,fillstyle=solid,fillcolor=curcolor]
{
\newpath
\moveto(207.6305039,159.15799431)
\curveto(207.6305039,155.4874411)(204.77438962,152.51186981)(201.2511941,152.51186981)
\curveto(197.72799858,152.51186981)(194.8718843,155.4874411)(194.8718843,159.15799431)
\curveto(194.8718843,162.82854752)(197.72799858,165.80411881)(201.2511941,165.80411881)
\curveto(204.77438962,165.80411881)(207.6305039,162.82854752)(207.6305039,159.15799431)
\closepath
}
}
{
\newrgbcolor{curcolor}{0 0 0}
\pscustom[linestyle=none,fillstyle=solid,fillcolor=curcolor]
{
\newpath
\moveto(218.4753279,20.25398731)
\curveto(218.4753279,16.5834341)(215.61921362,13.60786281)(212.0960181,13.60786281)
\curveto(208.57282258,13.60786281)(205.7167083,16.5834341)(205.7167083,20.25398731)
\curveto(205.7167083,23.92454052)(208.57282258,26.90011181)(212.0960181,26.90011181)
\curveto(215.61921362,26.90011181)(218.4753279,23.92454052)(218.4753279,20.25398731)
\closepath
}
}
{
\newrgbcolor{curcolor}{0 0 0}
\pscustom[linestyle=none,fillstyle=solid,fillcolor=curcolor]
{
\newpath
\moveto(235.0615449,76.08144131)
\curveto(235.0615449,72.4108881)(232.20543062,69.43531681)(228.6822351,69.43531681)
\curveto(225.15903958,69.43531681)(222.3029253,72.4108881)(222.3029253,76.08144131)
\curveto(222.3029253,79.75199452)(225.15903958,82.72756581)(228.6822351,82.72756581)
\curveto(232.20543062,82.72756581)(235.0615449,79.75199452)(235.0615449,76.08144131)
\closepath
}
}
{
\newrgbcolor{curcolor}{0 0 0}
\pscustom[linestyle=none,fillstyle=solid,fillcolor=curcolor]
{
\newpath
\moveto(192.9580999,57.47229231)
\curveto(192.9580999,53.8017391)(190.10198562,50.82616781)(186.5787901,50.82616781)
\curveto(183.05559458,50.82616781)(180.1994803,53.8017391)(180.1994803,57.47229231)
\curveto(180.1994803,61.14284552)(183.05559458,64.11841681)(186.5787901,64.11841681)
\curveto(190.10198562,64.11841681)(192.9580999,61.14284552)(192.9580999,57.47229231)
\closepath
}
}
{
\newrgbcolor{curcolor}{0 0 0}
\pscustom[linestyle=none,fillstyle=solid,fillcolor=curcolor]
{
\newpath
\moveto(183.3891339,135.23194631)
\curveto(183.3891339,131.5613931)(180.53301962,128.58582181)(177.0098241,128.58582181)
\curveto(173.48662858,128.58582181)(170.6305143,131.5613931)(170.6305143,135.23194631)
\curveto(170.6305143,138.90249952)(173.48662858,141.87807081)(177.0098241,141.87807081)
\curveto(180.53301962,141.87807081)(183.3891339,138.90249952)(183.3891339,135.23194631)
\closepath
}
}
{
\newrgbcolor{curcolor}{0 0 0}
\pscustom[linestyle=none,fillstyle=solid,fillcolor=curcolor]
{
\newpath
\moveto(233.7856849,135.89655831)
\curveto(233.7856849,132.2260051)(230.92957062,129.25043381)(227.4063751,129.25043381)
\curveto(223.88317958,129.25043381)(221.0270653,132.2260051)(221.0270653,135.89655831)
\curveto(221.0270653,139.56711152)(223.88317958,142.54268281)(227.4063751,142.54268281)
\curveto(230.92957062,142.54268281)(233.7856849,139.56711152)(233.7856849,135.89655831)
\closepath
}
}
{
\newrgbcolor{curcolor}{0 0 0}
\pscustom[linestyle=none,fillstyle=solid,fillcolor=curcolor]
{
\newpath
\moveto(219.7511999,181.09020431)
\curveto(219.7511999,177.4196511)(216.89508562,174.44407981)(213.3718901,174.44407981)
\curveto(209.84869458,174.44407981)(206.9925803,177.4196511)(206.9925803,181.09020431)
\curveto(206.9925803,184.76075752)(209.84869458,187.73632881)(213.3718901,187.73632881)
\curveto(216.89508562,187.73632881)(219.7511999,184.76075752)(219.7511999,181.09020431)
\closepath
}
}
{
\newrgbcolor{curcolor}{0 0 0}
\pscustom[linestyle=none,fillstyle=solid,fillcolor=curcolor]
{
\newpath
\moveto(217.1994789,111.97051031)
\curveto(217.1994789,108.2999571)(214.34336462,105.32438581)(210.8201691,105.32438581)
\curveto(207.29697358,105.32438581)(204.4408593,108.2999571)(204.4408593,111.97051031)
\curveto(204.4408593,115.64106352)(207.29697358,118.61663481)(210.8201691,118.61663481)
\curveto(214.34336462,118.61663481)(217.1994789,115.64106352)(217.1994789,111.97051031)
\closepath
}
}
{
\newrgbcolor{curcolor}{0 0 0}
\pscustom[linestyle=none,fillstyle=solid,fillcolor=curcolor]
{
\newpath
\moveto(247.8201749,159.82260731)
\curveto(247.8201749,156.1520541)(244.96406062,153.17648281)(241.4408651,153.17648281)
\curveto(237.91766958,153.17648281)(235.0615553,156.1520541)(235.0615553,159.82260731)
\curveto(235.0615553,163.49316052)(237.91766958,166.46873181)(241.4408651,166.46873181)
\curveto(244.96406062,166.46873181)(247.8201749,163.49316052)(247.8201749,159.82260731)
\closepath
}
}
{
\newrgbcolor{curcolor}{0 0 0}
\pscustom[linestyle=none,fillstyle=solid,fillcolor=curcolor]
{
\newpath
\moveto(242.7167249,45.50926731)
\curveto(242.7167249,41.8387141)(239.86061062,38.86314281)(236.3374151,38.86314281)
\curveto(232.81421958,38.86314281)(229.9581053,41.8387141)(229.9581053,45.50926731)
\curveto(229.9581053,49.17982052)(232.81421958,52.15539181)(236.3374151,52.15539181)
\curveto(239.86061062,52.15539181)(242.7167249,49.17982052)(242.7167249,45.50926731)
\closepath
}
}
{
\newrgbcolor{curcolor}{0 0 0}
\pscustom[linestyle=none,fillstyle=solid,fillcolor=curcolor]
{
\newpath
\moveto(257.3891349,107.98283631)
\curveto(257.3891349,104.3122831)(254.53302062,101.33671181)(251.0098251,101.33671181)
\curveto(247.48662958,101.33671181)(244.6305153,104.3122831)(244.6305153,107.98283631)
\curveto(244.6305153,111.65338952)(247.48662958,114.62896081)(251.0098251,114.62896081)
\curveto(254.53302062,114.62896081)(257.3891349,111.65338952)(257.3891349,107.98283631)
\closepath
}
}
{
\newrgbcolor{curcolor}{0 0 0}
\pscustom[linewidth=2.60454035,linecolor=curcolor]
{
\newpath
\moveto(55.165019,185.07788066)
\lineto(204.440859,181.09020566)
}
}
{
\newrgbcolor{curcolor}{0 0 0}
\pscustom[linestyle=none,fillstyle=solid,fillcolor=curcolor]
{
\newpath
\moveto(190.98930092,187.7554078)
\lineto(207.88968189,181.02310889)
\lineto(190.65397378,175.20270459)
\curveto(193.47972185,178.83541845)(193.59944327,183.90568063)(190.98930092,187.7554078)
\closepath
}
}
{
\newrgbcolor{curcolor}{0 0 0}
\pscustom[linewidth=2.60454035,linecolor=curcolor]
{
\newpath
\moveto(67.923639,139.88424066)
\lineto(70.475359,139.88424066)
\lineto(226.768439,81.39835066)
}
}
{
\newrgbcolor{curcolor}{0 0 0}
\pscustom[linestyle=none,fillstyle=solid,fillcolor=curcolor]
{
\newpath
\moveto(216.21710365,92.07723903)
\lineto(230.00781087,80.21287559)
\lineto(211.81616448,80.31651895)
\curveto(215.6700419,82.83228744)(217.43279236,87.58776875)(216.21710365,92.07723903)
\closepath
}
}
{
\newrgbcolor{curcolor}{0 0 0}
\pscustom[linewidth=2.60454035,linecolor=curcolor]
{
\newpath
\moveto(21.9926,112.63512066)
\lineto(222.302929,73.42299066)
}
}
{
\newrgbcolor{curcolor}{0 0 0}
\pscustom[linestyle=none,fillstyle=solid,fillcolor=curcolor]
{
\newpath
\moveto(210.14295648,82.22664915)
\lineto(225.69287046,72.78488549)
\lineto(207.73058797,69.90336709)
\curveto(211.11964222,73.0171708)(212.07853564,77.99737312)(210.14295648,82.22664915)
\closepath
}
}
{
\newrgbcolor{curcolor}{0 0 0}
\pscustom[linewidth=2.60454035,linecolor=curcolor]
{
\newpath
\moveto(29.647772,61.45997066)
\lineto(228.578789,46.41245066)
}
}
{
\newrgbcolor{curcolor}{0 0 0}
\pscustom[linestyle=none,fillstyle=solid,fillcolor=curcolor]
{
\newpath
\moveto(215.46831183,53.72577568)
\lineto(232.02023521,46.17722893)
\lineto(214.52116762,41.20436521)
\curveto(217.5207264,44.69493956)(217.88759042,49.75332888)(215.46831183,53.72577568)
\closepath
}
}
{
\newrgbcolor{curcolor}{0 0 0}
\pscustom[linewidth=2.60454035,linecolor=curcolor]
{
\newpath
\moveto(54.527079,24.90629066)
\lineto(61.544329,28.89396066)
\lineto(235.061549,155.17032066)
}
}
{
\newrgbcolor{curcolor}{0 0 0}
\pscustom[linestyle=none,fillstyle=solid,fillcolor=curcolor]
{
\newpath
\moveto(220.33601708,152.25003917)
\lineto(237.83583884,157.2202468)
\lineto(227.72493355,142.09687633)
\curveto(227.74924669,146.69914436)(224.75226919,150.79060102)(220.33601708,152.25003917)
\closepath
}
}
{
\newrgbcolor{curcolor}{0 0 0}
\pscustom[linestyle=none,fillstyle=solid,fillcolor=curcolor]
{
\newpath
\moveto(14.95321286,200.27087924)
\lineto(8.50376874,200.27087924)
\lineto(7.37358044,197.87859176)
\curveto(7.09512459,197.28983356)(6.95589864,196.85013006)(6.95590215,196.55947997)
\curveto(6.95589864,196.32844796)(7.07669763,196.12536457)(7.31829949,195.95022918)
\curveto(7.55989359,195.77509229)(8.08199093,195.66144012)(8.88459306,195.60927232)
\lineto(8.88459306,195.19565252)
\lineto(3.63904517,195.19565252)
\lineto(3.63904517,195.60927232)
\curveto(4.33517476,195.72106093)(4.78561167,195.86638666)(4.99035727,196.04524995)
\curveto(5.4080336,196.40297397)(5.87075516,197.12960262)(6.37852334,198.22513808)
\lineto(12.238304,210.70080552)
\lineto(12.66826694,210.70080552)
\lineto(18.46662432,198.09099112)
\curveto(18.93342574,197.07743441)(19.35724593,196.41974233)(19.73808617,196.11791289)
\curveto(20.11889381,195.8160816)(20.64918091,195.64653492)(21.32894905,195.60927232)
\lineto(21.32894905,195.19565252)
\lineto(14.75665837,195.19565252)
\lineto(14.75665837,195.60927232)
\curveto(15.42001779,195.63908232)(15.86840727,195.73969244)(16.10182815,195.91110299)
\curveto(16.33522371,196.08251211)(16.45192782,196.29118495)(16.45194083,196.53712214)
\curveto(16.45192782,196.86503527)(16.28813258,197.38299107)(15.96055461,198.09099112)
\closepath
\moveto(14.60924251,201.09811884)
\lineto(11.78377175,207.22416343)
\lineto(8.88459306,201.09811884)
\closepath
}
}
{
\newrgbcolor{curcolor}{0 0 0}
\pscustom[linestyle=none,fillstyle=solid,fillcolor=curcolor]
{
\newpath
\moveto(172.34898143,202.26147547)
\curveto(173.50372627,202.03788969)(174.36774617,201.68016481)(174.94104373,201.18829978)
\curveto(175.73543645,200.50265377)(176.13263991,199.6642361)(176.13265531,198.67304424)
\curveto(176.13263991,197.92032734)(175.87056752,197.19928814)(175.34643736,196.50992449)
\curveto(174.82227797,195.82055686)(174.10362634,195.31750626)(173.19048033,195.00077117)
\curveto(172.27730939,194.68403513)(170.88300239,194.52566734)(169.00755514,194.52566734)
\lineto(161.14537564,194.52566734)
\lineto(161.14537564,194.93928714)
\lineto(161.77189307,194.93928714)
\curveto(162.46802181,194.93928673)(162.9675973,195.14050697)(163.27062104,195.54294847)
\curveto(163.45898302,195.80378851)(163.55316529,196.35900732)(163.55316811,197.20860658)
\lineto(163.55316811,207.00133478)
\curveto(163.55316529,207.9403501)(163.43441374,208.53283192)(163.1969131,208.77878203)
\curveto(162.87750991,209.10668224)(162.40250371,209.27063947)(161.77189307,209.27065422)
\lineto(161.14537564,209.27065422)
\lineto(161.14537564,209.68427402)
\lineto(168.34418375,209.68427402)
\curveto(169.68729712,209.68425886)(170.76425084,209.59482764)(171.57504813,209.41598009)
\curveto(172.80350161,209.14767155)(173.74122937,208.6725682)(174.38823424,207.99066863)
\curveto(175.03521178,207.30874212)(175.35870739,206.52435581)(175.35872202,205.63750734)
\curveto(175.35870739,204.87733087)(175.10482476,204.19728098)(174.59707338,203.59735563)
\curveto(174.08929426,202.99741213)(173.33993102,202.55211919)(172.34898143,202.26147547)
\closepath
\moveto(165.91182196,202.8651368)
\curveto(166.21483798,202.81296025)(166.56085543,202.77383409)(166.94987535,202.74775821)
\curveto(167.33888283,202.72166588)(167.7667979,202.70862383)(168.23362185,202.70863201)
\curveto(169.42931961,202.70862383)(170.328146,202.8260023)(170.93010372,203.06076778)
\curveto(171.53204104,203.2955162)(171.99271516,203.65510422)(172.31212746,204.13953293)
\curveto(172.6315166,204.62394242)(172.79121696,205.15307713)(172.79122903,205.72693865)
\curveto(172.79121696,206.61378703)(172.3940135,207.37022608)(171.59961745,207.99625808)
\curveto(170.80519965,208.62226314)(169.64634831,208.93527241)(168.12305995,208.93528681)
\curveto(167.30407634,208.93527241)(166.56699775,208.85329379)(165.91182196,208.68935072)
\closepath
\moveto(165.91182196,195.62120086)
\curveto(166.86182919,195.41997953)(167.79955695,195.31936941)(168.72500807,195.3193702)
\curveto(170.20734701,195.31936941)(171.33753419,195.62306292)(172.11557298,196.23045165)
\curveto(172.89358899,196.83783697)(173.28260269,197.58868657)(173.28261524,198.48300271)
\curveto(173.28260269,199.07175428)(173.1065228,199.63815199)(172.75437506,200.18219756)
\curveto(172.40220326,200.72623182)(171.82891991,201.15475641)(171.03452329,201.46777261)
\curveto(170.24010606,201.78077493)(169.25733461,201.93727957)(168.08620598,201.93728698)
\curveto(167.57843337,201.93727957)(167.14437598,201.92982697)(166.7840325,201.91492915)
\curveto(166.42367691,201.90001656)(166.13294036,201.87393245)(165.91182196,201.83667676)
\closepath
}
}
{
\newrgbcolor{curcolor}{0 0 0}
\pscustom[linewidth=2.60454035,linecolor=curcolor]
{
\newpath
\moveto(22.953127,114.81860066)
\lineto(169.901389,133.30628066)
}
}
{
\newrgbcolor{curcolor}{0 0 0}
\pscustom[linestyle=none,fillstyle=solid,fillcolor=curcolor]
{
\newpath
\moveto(155.59634599,137.85985934)
\lineto(173.32067084,133.76168486)
\lineto(157.16381766,125.40089329)
\curveto(159.40749789,129.41927123)(158.75882566,134.44929271)(155.59634599,137.85985934)
\closepath
}
}
{
\newrgbcolor{curcolor}{0 0 0}
\pscustom[linewidth=2.60454035,linecolor=curcolor]
{
\newpath
\moveto(53.453129,23.31860066)
\lineto(204.901389,20.30628066)
}
}
{
\newrgbcolor{curcolor}{0 0 0}
\pscustom[linestyle=none,fillstyle=solid,fillcolor=curcolor]
{
\newpath
\moveto(191.40468977,26.8795937)
\lineto(208.35058926,20.26270501)
\lineto(191.15497472,14.32489559)
\curveto(193.95588347,17.97679542)(194.04102498,23.04775614)(191.40468977,26.8795937)
\closepath
}
}
\end{pspicture}

\footnotesize\caption{This picture does {\em not} represent a function.} \label{funcNot}
\end{minipage}
\end{figure}

\begin{rem}
In figures~\ref{func1} and~\ref{funcNot}, the dots represent all of the elements of the
sets $A$ and $B$ and the gray ovals are mainly there to help identify which dots are in which set.
However, in these sorts of diagrams it is more common for the dots to represent only {\em some} of the elements.
You need to let the context help you determine how to properly interpret these diagrams.
\end{rem}


\begin{exa}
Give a formal definition of a function that assigns to an age the number 
of complete decades someone of that age has lived.
For instance, $f(34)=3$ and $f(5)=0$.
Be sure to indicate what the domain and codomain are.
\begin{sol}
It isn't hard to see that the domain and codomain are both $\naturals$.
Thus we want a function $f:\naturals\to\naturals$.
One way to define $f$ is by $f(x)=\lfloor x/10 \rfloor$.
\end{sol}
\end{exa}

\begin{exer}
Give a formal definition of a function that returns the parity of an integer.
That is, it returns $0$ for even numbers and $1$ for odd numbers.
Be sure to indicate what the domain and codomain are.

\noindent
\answerlinetwo
\begin{answer}
$f(x)=x\bmod 2$ works.  The domain is $\BBZ$, and the codomain can be a variety
of things.  $\BBZ$, $\naturals$, and $\{0,1\}$ are the most obvious choices.
Note that we can pick any of these since the only requirement of the codomain is 
that the range is a subset of it.  On the other hand, $\BBR$, $\BBC$ and $\BBQ$
could also all be given as the codomain, but they wouldn't make nearly as much sense.
\end{answer}
\end{exer}

\begin{df}
Let $f : A\to B$ be a function. 
\begin{itemize}
\item
$f$ is said to be {\bf injective} or {\bf one-to-one}\index{function!injective} 
if and only if $f(a)=f(b)$ implies that $a=b$.
In other words, $f$ maps every element of $A$ to a {\em different} element of $B$.
\item
$f$ is said to be {\bf surjective} or {\bf onto} if and only if for every $b\in B$,
there exists some $a\in A$ such that $f(a)=b$.
In other words, every element in $B$ gets mapped to by some element in $A$.
\item
$f$ is said to be {\bf bijective} or a {\bf one-to-one correspondence} 
if it is both injective and surjective.
\end{itemize}
\end{df}

\begin{figure}[ht]
\begin{minipage}{2in}
%LaTeX with PSTricks extensions
%%Creator: 0.48.3.1
%%Please note this file requires PSTricks extensions
\psset{xunit=.5pt,yunit=.5pt,runit=.5pt}
\begin{pspicture}(269.81277466,209.31954956)
{
\newrgbcolor{curcolor}{0.7019608 0.7019608 0.7019608}
\pscustom[linestyle=none,fillstyle=solid,fillcolor=curcolor]
{
\newpath
\moveto(269.8127758,102.66593362)
\curveto(269.8127758,45.9651001)(244.54333166,-0.00000733)(213.3718926,-0.00000733)
\curveto(182.20045354,-0.00000733)(156.9310094,45.9651001)(156.9310094,102.66593362)
\curveto(156.9310094,159.36676713)(182.20045354,205.33187457)(213.3718926,205.33187457)
\curveto(244.54333166,205.33187457)(269.8127758,159.36676713)(269.8127758,102.66593362)
\closepath
}
}
{
\newrgbcolor{curcolor}{0.7019608 0.7019608 0.7019608}
\pscustom[linestyle=none,fillstyle=solid,fillcolor=curcolor]
{
\newpath
\moveto(112.8817668,106.65360862)
\curveto(112.8817668,49.9527751)(87.61232266,3.98766767)(56.4408836,3.98766767)
\curveto(25.26944454,3.98766767)(0.0000004,49.9527751)(0.0000004,106.65360862)
\curveto(0.0000004,163.35444213)(25.26944454,209.31954957)(56.4408836,209.31954957)
\curveto(87.61232266,209.31954957)(112.8817668,163.35444213)(112.8817668,106.65360862)
\closepath
}
}
{
\newrgbcolor{curcolor}{0 0 0}
\pscustom[linestyle=none,fillstyle=solid,fillcolor=curcolor]
{
\newpath
\moveto(34.7512229,153.17647821)
\curveto(34.7512229,149.505925)(31.89510862,146.53035371)(28.3719131,146.53035371)
\curveto(24.84871758,146.53035371)(21.9926033,149.505925)(21.9926033,153.17647821)
\curveto(21.9926033,156.84703141)(24.84871758,159.82260271)(28.3719131,159.82260271)
\curveto(31.89510862,159.82260271)(34.7512229,156.84703141)(34.7512229,153.17647821)
\closepath
}
}
{
\newrgbcolor{curcolor}{0 0 0}
\pscustom[linestyle=none,fillstyle=solid,fillcolor=curcolor]
{
\newpath
\moveto(61.5443199,24.24166321)
\curveto(61.5443199,20.57111)(58.68820562,17.59553871)(55.1650101,17.59553871)
\curveto(51.64181458,17.59553871)(48.7857003,20.57111)(48.7857003,24.24166321)
\curveto(48.7857003,27.91221641)(51.64181458,30.88778771)(55.1650101,30.88778771)
\curveto(58.68820562,30.88778771)(61.5443199,27.91221641)(61.5443199,24.24166321)
\closepath
}
}
{
\newrgbcolor{curcolor}{0 0 0}
\pscustom[linestyle=none,fillstyle=solid,fillcolor=curcolor]
{
\newpath
\moveto(78.1305369,80.06910921)
\curveto(78.1305369,76.398556)(75.27442262,73.42298471)(71.7512271,73.42298471)
\curveto(68.22803158,73.42298471)(65.3719173,76.398556)(65.3719173,80.06910921)
\curveto(65.3719173,83.73966241)(68.22803158,86.71523371)(71.7512271,86.71523371)
\curveto(75.27442262,86.71523371)(78.1305369,83.73966241)(78.1305369,80.06910921)
\closepath
}
}
{
\newrgbcolor{curcolor}{0 0 0}
\pscustom[linestyle=none,fillstyle=solid,fillcolor=curcolor]
{
\newpath
\moveto(36.0270919,61.45996021)
\curveto(36.0270919,57.789407)(33.17097762,54.81383571)(29.6477821,54.81383571)
\curveto(26.12458658,54.81383571)(23.2684723,57.789407)(23.2684723,61.45996021)
\curveto(23.2684723,65.13051341)(26.12458658,68.10608471)(29.6477821,68.10608471)
\curveto(33.17097762,68.10608471)(36.0270919,65.13051341)(36.0270919,61.45996021)
\closepath
}
}
{
\newrgbcolor{curcolor}{0 0 0}
\pscustom[linestyle=none,fillstyle=solid,fillcolor=curcolor]
{
\newpath
\moveto(29.6477809,113.29973121)
\curveto(29.6477809,109.629178)(26.79166662,106.65360671)(23.2684711,106.65360671)
\curveto(19.74527558,106.65360671)(16.8891613,109.629178)(16.8891613,113.29973121)
\curveto(16.8891613,116.97028441)(19.74527558,119.94585571)(23.2684711,119.94585571)
\curveto(26.79166662,119.94585571)(29.6477809,116.97028441)(29.6477809,113.29973121)
\closepath
}
}
{
\newrgbcolor{curcolor}{0 0 0}
\pscustom[linestyle=none,fillstyle=solid,fillcolor=curcolor]
{
\newpath
\moveto(76.8546749,139.88422911)
\curveto(76.8546749,136.2136759)(73.99856062,133.23810461)(70.4753651,133.23810461)
\curveto(66.95216958,133.23810461)(64.0960553,136.2136759)(64.0960553,139.88422911)
\curveto(64.0960553,143.55478231)(66.95216958,146.53035361)(70.4753651,146.53035361)
\curveto(73.99856062,146.53035361)(76.8546749,143.55478231)(76.8546749,139.88422911)
\closepath
}
}
{
\newrgbcolor{curcolor}{0 0 0}
\pscustom[linestyle=none,fillstyle=solid,fillcolor=curcolor]
{
\newpath
\moveto(62.8201919,185.07787621)
\curveto(62.8201919,181.407323)(59.96407762,178.43175171)(56.4408821,178.43175171)
\curveto(52.91768658,178.43175171)(50.0615723,181.407323)(50.0615723,185.07787621)
\curveto(50.0615723,188.74842941)(52.91768658,191.72400071)(56.4408821,191.72400071)
\curveto(59.96407762,191.72400071)(62.8201919,188.74842941)(62.8201919,185.07787621)
\closepath
}
}
{
\newrgbcolor{curcolor}{0 0 0}
\pscustom[linestyle=none,fillstyle=solid,fillcolor=curcolor]
{
\newpath
\moveto(207.6305039,159.15799021)
\curveto(207.6305039,155.487437)(204.77438962,152.51186571)(201.2511941,152.51186571)
\curveto(197.72799858,152.51186571)(194.8718843,155.487437)(194.8718843,159.15799021)
\curveto(194.8718843,162.82854341)(197.72799858,165.80411471)(201.2511941,165.80411471)
\curveto(204.77438962,165.80411471)(207.6305039,162.82854341)(207.6305039,159.15799021)
\closepath
}
}
{
\newrgbcolor{curcolor}{0 0 0}
\pscustom[linestyle=none,fillstyle=solid,fillcolor=curcolor]
{
\newpath
\moveto(218.4753279,20.25398321)
\curveto(218.4753279,16.58343)(215.61921362,13.60785871)(212.0960181,13.60785871)
\curveto(208.57282258,13.60785871)(205.7167083,16.58343)(205.7167083,20.25398321)
\curveto(205.7167083,23.92453641)(208.57282258,26.90010771)(212.0960181,26.90010771)
\curveto(215.61921362,26.90010771)(218.4753279,23.92453641)(218.4753279,20.25398321)
\closepath
}
}
{
\newrgbcolor{curcolor}{0 0 0}
\pscustom[linestyle=none,fillstyle=solid,fillcolor=curcolor]
{
\newpath
\moveto(235.0615449,76.08143721)
\curveto(235.0615449,72.410884)(232.20543062,69.43531271)(228.6822351,69.43531271)
\curveto(225.15903958,69.43531271)(222.3029253,72.410884)(222.3029253,76.08143721)
\curveto(222.3029253,79.75199041)(225.15903958,82.72756171)(228.6822351,82.72756171)
\curveto(232.20543062,82.72756171)(235.0615449,79.75199041)(235.0615449,76.08143721)
\closepath
}
}
{
\newrgbcolor{curcolor}{0 0 0}
\pscustom[linestyle=none,fillstyle=solid,fillcolor=curcolor]
{
\newpath
\moveto(192.9580999,57.47228821)
\curveto(192.9580999,53.801735)(190.10198562,50.82616371)(186.5787901,50.82616371)
\curveto(183.05559458,50.82616371)(180.1994803,53.801735)(180.1994803,57.47228821)
\curveto(180.1994803,61.14284141)(183.05559458,64.11841271)(186.5787901,64.11841271)
\curveto(190.10198562,64.11841271)(192.9580999,61.14284141)(192.9580999,57.47228821)
\closepath
}
}
{
\newrgbcolor{curcolor}{0 0 0}
\pscustom[linestyle=none,fillstyle=solid,fillcolor=curcolor]
{
\newpath
\moveto(183.3891339,135.23194221)
\curveto(183.3891339,131.561389)(180.53301962,128.58581771)(177.0098241,128.58581771)
\curveto(173.48662858,128.58581771)(170.6305143,131.561389)(170.6305143,135.23194221)
\curveto(170.6305143,138.90249541)(173.48662858,141.87806671)(177.0098241,141.87806671)
\curveto(180.53301962,141.87806671)(183.3891339,138.90249541)(183.3891339,135.23194221)
\closepath
}
}
{
\newrgbcolor{curcolor}{0 0 0}
\pscustom[linestyle=none,fillstyle=solid,fillcolor=curcolor]
{
\newpath
\moveto(231.7856849,122.39655421)
\curveto(231.7856849,118.726001)(228.92957062,115.75042971)(225.4063751,115.75042971)
\curveto(221.88317958,115.75042971)(219.0270653,118.726001)(219.0270653,122.39655421)
\curveto(219.0270653,126.06710741)(221.88317958,129.04267871)(225.4063751,129.04267871)
\curveto(228.92957062,129.04267871)(231.7856849,126.06710741)(231.7856849,122.39655421)
\closepath
}
}
{
\newrgbcolor{curcolor}{0 0 0}
\pscustom[linestyle=none,fillstyle=solid,fillcolor=curcolor]
{
\newpath
\moveto(219.7511999,181.09020021)
\curveto(219.7511999,177.419647)(216.89508562,174.44407571)(213.3718901,174.44407571)
\curveto(209.84869458,174.44407571)(206.9925803,177.419647)(206.9925803,181.09020021)
\curveto(206.9925803,184.76075341)(209.84869458,187.73632471)(213.3718901,187.73632471)
\curveto(216.89508562,187.73632471)(219.7511999,184.76075341)(219.7511999,181.09020021)
\closepath
}
}
{
\newrgbcolor{curcolor}{0 0 0}
\pscustom[linestyle=none,fillstyle=solid,fillcolor=curcolor]
{
\newpath
\moveto(200.6994789,103.47050621)
\curveto(200.6994789,99.799953)(197.84336462,96.82438171)(194.3201691,96.82438171)
\curveto(190.79697358,96.82438171)(187.9408593,99.799953)(187.9408593,103.47050621)
\curveto(187.9408593,107.14105941)(190.79697358,110.11663071)(194.3201691,110.11663071)
\curveto(197.84336462,110.11663071)(200.6994789,107.14105941)(200.6994789,103.47050621)
\closepath
}
}
{
\newrgbcolor{curcolor}{0 0 0}
\pscustom[linestyle=none,fillstyle=solid,fillcolor=curcolor]
{
\newpath
\moveto(247.8201749,159.82260321)
\curveto(247.8201749,156.15205)(244.96406062,153.17647871)(241.4408651,153.17647871)
\curveto(237.91766958,153.17647871)(235.0615553,156.15205)(235.0615553,159.82260321)
\curveto(235.0615553,163.49315641)(237.91766958,166.46872771)(241.4408651,166.46872771)
\curveto(244.96406062,166.46872771)(247.8201749,163.49315641)(247.8201749,159.82260321)
\closepath
}
}
{
\newrgbcolor{curcolor}{0 0 0}
\pscustom[linestyle=none,fillstyle=solid,fillcolor=curcolor]
{
\newpath
\moveto(242.7167249,45.50926321)
\curveto(242.7167249,41.83871)(239.86061062,38.86313871)(236.3374151,38.86313871)
\curveto(232.81421958,38.86313871)(229.9581053,41.83871)(229.9581053,45.50926321)
\curveto(229.9581053,49.17981641)(232.81421958,52.15538771)(236.3374151,52.15538771)
\curveto(239.86061062,52.15538771)(242.7167249,49.17981641)(242.7167249,45.50926321)
\closepath
}
}
{
\newrgbcolor{curcolor}{0 0 0}
\pscustom[linestyle=none,fillstyle=solid,fillcolor=curcolor]
{
\newpath
\moveto(257.3891349,107.98283221)
\curveto(257.3891349,104.312279)(254.53302062,101.33670771)(251.0098251,101.33670771)
\curveto(247.48662958,101.33670771)(244.6305153,104.312279)(244.6305153,107.98283221)
\curveto(244.6305153,111.65338541)(247.48662958,114.62895671)(251.0098251,114.62895671)
\curveto(254.53302062,114.62895671)(257.3891349,111.65338541)(257.3891349,107.98283221)
\closepath
}
}
{
\newrgbcolor{curcolor}{0 0 0}
\pscustom[linewidth=2.60454035,linecolor=curcolor]
{
\newpath
\moveto(55.165019,185.07787656)
\lineto(204.440859,181.09020156)
}
}
{
\newrgbcolor{curcolor}{0 0 0}
\pscustom[linestyle=none,fillstyle=solid,fillcolor=curcolor]
{
\newpath
\moveto(190.98930092,187.7554037)
\lineto(207.88968189,181.02310479)
\lineto(190.65397378,175.20270048)
\curveto(193.47972185,178.83541435)(193.59944327,183.90567652)(190.98930092,187.7554037)
\closepath
}
}
{
\newrgbcolor{curcolor}{0 0 0}
\pscustom[linewidth=2.60454035,linecolor=curcolor]
{
\newpath
\moveto(29.009841,153.17647656)
\lineto(192.958099,157.16415656)
}
}
{
\newrgbcolor{curcolor}{0 0 0}
\pscustom[linestyle=none,fillstyle=solid,fillcolor=curcolor]
{
\newpath
\moveto(179.18410296,163.13461445)
\lineto(196.40585538,157.27304735)
\lineto(179.48943747,150.58114587)
\curveto(182.12622862,154.35325333)(181.98719575,159.42302271)(179.18410296,163.13461445)
\closepath
}
}
{
\newrgbcolor{curcolor}{0 0 0}
\pscustom[linewidth=2.60454035,linecolor=curcolor]
{
\newpath
\moveto(67.923639,139.88423656)
\lineto(70.475359,139.88423656)
\lineto(168.768439,136.39834656)
}
}
{
\newrgbcolor{curcolor}{0 0 0}
\pscustom[linestyle=none,fillstyle=solid,fillcolor=curcolor]
{
\newpath
\moveto(155.37566324,143.18088973)
\lineto(172.21654352,136.30110217)
\lineto(154.93061112,130.63159768)
\curveto(157.788009,134.23946957)(157.9520509,139.30849137)(155.37566324,143.18088973)
\closepath
}
}
{
\newrgbcolor{curcolor}{0 0 0}
\pscustom[linewidth=2.60454035,linecolor=curcolor]
{
\newpath
\moveto(21.9926,112.63511656)
\lineto(220.802929,77.92298656)
}
}
{
\newrgbcolor{curcolor}{0 0 0}
\pscustom[linestyle=none,fillstyle=solid,fillcolor=curcolor]
{
\newpath
\moveto(208.46542911,86.47608137)
\lineto(224.20521315,77.35435325)
\lineto(206.30562752,74.10603451)
\curveto(209.63028255,77.288507)(210.48710957,82.28728077)(208.46542911,86.47608137)
\closepath
}
}
{
\newrgbcolor{curcolor}{0 0 0}
\pscustom[linewidth=2.60454035,linecolor=curcolor]
{
\newpath
\moveto(29.647772,61.45996656)
\lineto(205.078789,22.91244656)
}
}
{
\newrgbcolor{curcolor}{0 0 0}
\pscustom[linestyle=none,fillstyle=solid,fillcolor=curcolor]
{
\newpath
\moveto(193.12432408,31.9932011)
\lineto(208.45317252,22.19661336)
\lineto(190.42941929,19.7286061)
\curveto(193.8891299,22.76371148)(194.96220881,27.72056466)(193.12432408,31.9932011)
\closepath
}
}
{
\newrgbcolor{curcolor}{0 0 0}
\pscustom[linewidth=2.60454035,linecolor=curcolor]
{
\newpath
\moveto(54.527079,24.90628656)
\lineto(61.544329,28.89395656)
\lineto(235.061549,155.17031656)
}
}
{
\newrgbcolor{curcolor}{0 0 0}
\pscustom[linestyle=none,fillstyle=solid,fillcolor=curcolor]
{
\newpath
\moveto(220.33601708,152.25003507)
\lineto(237.83583884,157.2202427)
\lineto(227.72493355,142.09687222)
\curveto(227.74924669,146.69914026)(224.75226919,150.79059692)(220.33601708,152.25003507)
\closepath
}
}
{
\newrgbcolor{curcolor}{0 0 0}
\pscustom[linewidth=2.60454035,linecolor=curcolor]
{
\newpath
\moveto(73.027079,83.39217656)
\lineto(71.113289,80.06911656)
\lineto(178.285689,58.80151656)
}
}
{
\newrgbcolor{curcolor}{0 0 0}
\pscustom[linestyle=none,fillstyle=solid,fillcolor=curcolor]
{
\newpath
\moveto(166.14851925,67.63658559)
\lineto(181.67396929,58.15464873)
\lineto(163.70429659,55.31958198)
\curveto(167.10139034,58.42461278)(168.07315702,63.4023192)(166.14851925,67.63658559)
\closepath
}
}
{
\newrgbcolor{curcolor}{0 0 0}
\pscustom[linestyle=none,fillstyle=solid,fillcolor=curcolor]
{
\newpath
\moveto(14.95321312,197.27087955)
\lineto(8.50376898,197.27087955)
\lineto(7.37358067,194.87859206)
\curveto(7.09512482,194.28983385)(6.95589887,193.85013036)(6.95590238,193.55948026)
\curveto(6.95589887,193.32844825)(7.07669786,193.12536486)(7.31829972,192.95022948)
\curveto(7.55989383,192.77509259)(8.08199116,192.66144041)(8.8845933,192.60927262)
\lineto(8.8845933,192.19565282)
\lineto(3.63904539,192.19565282)
\lineto(3.63904539,192.60927262)
\curveto(4.33517498,192.72106123)(4.7856119,192.86638696)(4.9903575,193.04525024)
\curveto(5.40803382,193.40297427)(5.87075538,194.12960292)(6.37852357,195.22513838)
\lineto(12.23830425,207.70080586)
\lineto(12.66826719,207.70080586)
\lineto(18.4666246,195.09099142)
\curveto(18.93342602,194.07743471)(19.35724621,193.41974262)(19.73808644,193.11791318)
\curveto(20.11889409,192.8160819)(20.64918119,192.64653521)(21.32894933,192.60927262)
\lineto(21.32894933,192.19565282)
\lineto(14.75665863,192.19565282)
\lineto(14.75665863,192.60927262)
\curveto(15.42001805,192.63908261)(15.86840753,192.73969273)(16.10182841,192.91110328)
\curveto(16.33522397,193.0825124)(16.45192808,193.29118524)(16.45194109,193.53712244)
\curveto(16.45192808,193.86503556)(16.28813284,194.38299137)(15.96055487,195.09099142)
\closepath
\moveto(14.60924277,198.09811915)
\lineto(11.783772,204.22416376)
\lineto(8.8845933,198.09811915)
\closepath
}
}
{
\newrgbcolor{curcolor}{0 0 0}
\pscustom[linestyle=none,fillstyle=solid,fillcolor=curcolor]
{
\newpath
\moveto(171.34899768,198.26147272)
\curveto(172.50374252,198.03788694)(173.36776242,197.68016206)(173.94105998,197.18829703)
\curveto(174.73545269,196.50265102)(175.13265616,195.66423335)(175.13267156,194.67304149)
\curveto(175.13265616,193.92032459)(174.87058377,193.19928539)(174.34645361,192.50992174)
\curveto(173.82229422,191.82055411)(173.10364259,191.31750351)(172.19049658,191.00076842)
\curveto(171.27732564,190.68403238)(169.88301864,190.52566459)(168.00757139,190.52566459)
\lineto(160.14539189,190.52566459)
\lineto(160.14539189,190.93928439)
\lineto(160.77190932,190.93928439)
\curveto(161.46803805,190.93928398)(161.96761354,191.14050422)(162.27063729,191.54294572)
\curveto(162.45899927,191.80378576)(162.55318154,192.35900457)(162.55318436,193.20860383)
\lineto(162.55318436,203.00133203)
\curveto(162.55318154,203.94034735)(162.43442999,204.53282917)(162.19692935,204.77877928)
\curveto(161.87752616,205.10667949)(161.40251996,205.27063672)(160.77190932,205.27065147)
\lineto(160.14539189,205.27065147)
\lineto(160.14539189,205.68427127)
\lineto(167.34419999,205.68427127)
\curveto(168.68731337,205.68425611)(169.76426709,205.59482489)(170.57506438,205.41597734)
\curveto(171.80351785,205.1476688)(172.74124562,204.67256545)(173.38825049,203.99066588)
\curveto(174.03522803,203.30873937)(174.35872364,202.52435306)(174.35873827,201.63750459)
\curveto(174.35872364,200.87732812)(174.10484101,200.19727823)(173.59708963,199.59735288)
\curveto(173.08931051,198.99740939)(172.33994727,198.55211644)(171.34899768,198.26147272)
\closepath
\moveto(164.91183821,198.86513405)
\curveto(165.21485423,198.8129575)(165.56087168,198.77383134)(165.9498916,198.74775546)
\curveto(166.33889908,198.72166313)(166.76681415,198.70862108)(167.2336381,198.70862926)
\curveto(168.42933586,198.70862108)(169.32816225,198.82599955)(169.93011997,199.06076503)
\curveto(170.53205728,199.29551345)(170.9927314,199.65510147)(171.31214371,200.13953018)
\curveto(171.63153285,200.62393967)(171.79123321,201.15307438)(171.79124527,201.7269359)
\curveto(171.79123321,202.61378428)(171.39402975,203.37022333)(170.59963369,203.99625533)
\curveto(169.8052159,204.62226039)(168.64636456,204.93526966)(167.1230762,204.93528407)
\curveto(166.30409259,204.93526966)(165.567014,204.85329104)(164.91183821,204.68934797)
\closepath
\moveto(164.91183821,191.62119811)
\curveto(165.86184543,191.41997678)(166.7995732,191.31936666)(167.72502431,191.31936745)
\curveto(169.20736326,191.31936666)(170.33755043,191.62306017)(171.11558922,192.2304489)
\curveto(171.89360524,192.83783422)(172.28261894,193.58868382)(172.28263149,194.48299996)
\curveto(172.28261894,195.07175153)(172.10653905,195.63814924)(171.75439131,196.18219481)
\curveto(171.40221951,196.72622907)(170.82893616,197.15475366)(170.03453954,197.46776986)
\curveto(169.24012231,197.78077219)(168.25735085,197.93727682)(167.08622223,197.93728423)
\curveto(166.57844962,197.93727682)(166.14439223,197.92982422)(165.78404875,197.9149264)
\curveto(165.42369316,197.90001381)(165.1329566,197.8739297)(164.91183821,197.83667401)
\closepath
}
}
\end{pspicture}

\footnotesize\caption{Pictorial representation of a {\em one-to-one} function.}
\end{minipage}
\hfill
\begin{minipage}{2in}
%LaTeX with PSTricks extensions
%%Creator: 0.48.3.1
%%Please note this file requires PSTricks extensions
\psset{xunit=.5pt,yunit=.5pt,runit=.5pt}
\begin{pspicture}(269.81277466,209.31954956)
{
\newrgbcolor{curcolor}{0.7019608 0.7019608 0.7019608}
\pscustom[linestyle=none,fillstyle=solid,fillcolor=curcolor]
{
\newpath
\moveto(269.8127758,102.66593362)
\curveto(269.8127758,45.9651001)(244.54333166,-0.00000733)(213.3718926,-0.00000733)
\curveto(182.20045354,-0.00000733)(156.9310094,45.9651001)(156.9310094,102.66593362)
\curveto(156.9310094,159.36676713)(182.20045354,205.33187457)(213.3718926,205.33187457)
\curveto(244.54333166,205.33187457)(269.8127758,159.36676713)(269.8127758,102.66593362)
\closepath
}
}
{
\newrgbcolor{curcolor}{0.7019608 0.7019608 0.7019608}
\pscustom[linestyle=none,fillstyle=solid,fillcolor=curcolor]
{
\newpath
\moveto(112.8817668,106.65360862)
\curveto(112.8817668,49.9527751)(87.61232266,3.98766767)(56.4408836,3.98766767)
\curveto(25.26944454,3.98766767)(0.0000004,49.9527751)(0.0000004,106.65360862)
\curveto(0.0000004,163.35444213)(25.26944454,209.31954957)(56.4408836,209.31954957)
\curveto(87.61232266,209.31954957)(112.8817668,163.35444213)(112.8817668,106.65360862)
\closepath
}
}
{
\newrgbcolor{curcolor}{0 0 0}
\pscustom[linestyle=none,fillstyle=solid,fillcolor=curcolor]
{
\newpath
\moveto(34.7512229,153.17647821)
\curveto(34.7512229,149.505925)(31.89510862,146.53035371)(28.3719131,146.53035371)
\curveto(24.84871758,146.53035371)(21.9926033,149.505925)(21.9926033,153.17647821)
\curveto(21.9926033,156.84703141)(24.84871758,159.82260271)(28.3719131,159.82260271)
\curveto(31.89510862,159.82260271)(34.7512229,156.84703141)(34.7512229,153.17647821)
\closepath
}
}
{
\newrgbcolor{curcolor}{0 0 0}
\pscustom[linestyle=none,fillstyle=solid,fillcolor=curcolor]
{
\newpath
\moveto(61.5443199,24.24166321)
\curveto(61.5443199,20.57111)(58.68820562,17.59553871)(55.1650101,17.59553871)
\curveto(51.64181458,17.59553871)(48.7857003,20.57111)(48.7857003,24.24166321)
\curveto(48.7857003,27.91221641)(51.64181458,30.88778771)(55.1650101,30.88778771)
\curveto(58.68820562,30.88778771)(61.5443199,27.91221641)(61.5443199,24.24166321)
\closepath
}
}
{
\newrgbcolor{curcolor}{0 0 0}
\pscustom[linestyle=none,fillstyle=solid,fillcolor=curcolor]
{
\newpath
\moveto(78.1305369,80.06910921)
\curveto(78.1305369,76.398556)(75.27442262,73.42298471)(71.7512271,73.42298471)
\curveto(68.22803158,73.42298471)(65.3719173,76.398556)(65.3719173,80.06910921)
\curveto(65.3719173,83.73966241)(68.22803158,86.71523371)(71.7512271,86.71523371)
\curveto(75.27442262,86.71523371)(78.1305369,83.73966241)(78.1305369,80.06910921)
\closepath
}
}
{
\newrgbcolor{curcolor}{0 0 0}
\pscustom[linestyle=none,fillstyle=solid,fillcolor=curcolor]
{
\newpath
\moveto(36.0270919,61.45996021)
\curveto(36.0270919,57.789407)(33.17097762,54.81383571)(29.6477821,54.81383571)
\curveto(26.12458658,54.81383571)(23.2684723,57.789407)(23.2684723,61.45996021)
\curveto(23.2684723,65.13051341)(26.12458658,68.10608471)(29.6477821,68.10608471)
\curveto(33.17097762,68.10608471)(36.0270919,65.13051341)(36.0270919,61.45996021)
\closepath
}
}
{
\newrgbcolor{curcolor}{0 0 0}
\pscustom[linestyle=none,fillstyle=solid,fillcolor=curcolor]
{
\newpath
\moveto(29.6477809,113.29973121)
\curveto(29.6477809,109.629178)(26.79166662,106.65360671)(23.2684711,106.65360671)
\curveto(19.74527558,106.65360671)(16.8891613,109.629178)(16.8891613,113.29973121)
\curveto(16.8891613,116.97028441)(19.74527558,119.94585571)(23.2684711,119.94585571)
\curveto(26.79166662,119.94585571)(29.6477809,116.97028441)(29.6477809,113.29973121)
\closepath
}
}
{
\newrgbcolor{curcolor}{0 0 0}
\pscustom[linestyle=none,fillstyle=solid,fillcolor=curcolor]
{
\newpath
\moveto(76.8546749,139.88422911)
\curveto(76.8546749,136.2136759)(73.99856062,133.23810461)(70.4753651,133.23810461)
\curveto(66.95216958,133.23810461)(64.0960553,136.2136759)(64.0960553,139.88422911)
\curveto(64.0960553,143.55478231)(66.95216958,146.53035361)(70.4753651,146.53035361)
\curveto(73.99856062,146.53035361)(76.8546749,143.55478231)(76.8546749,139.88422911)
\closepath
}
}
{
\newrgbcolor{curcolor}{0 0 0}
\pscustom[linestyle=none,fillstyle=solid,fillcolor=curcolor]
{
\newpath
\moveto(62.8201919,185.07787621)
\curveto(62.8201919,181.407323)(59.96407762,178.43175171)(56.4408821,178.43175171)
\curveto(52.91768658,178.43175171)(50.0615723,181.407323)(50.0615723,185.07787621)
\curveto(50.0615723,188.74842941)(52.91768658,191.72400071)(56.4408821,191.72400071)
\curveto(59.96407762,191.72400071)(62.8201919,188.74842941)(62.8201919,185.07787621)
\closepath
}
}
{
\newrgbcolor{curcolor}{0 0 0}
\pscustom[linestyle=none,fillstyle=solid,fillcolor=curcolor]
{
\newpath
\moveto(260.1305049,118.15799021)
\curveto(260.1305049,114.487437)(257.27439062,111.51186571)(253.7511951,111.51186571)
\curveto(250.22799958,111.51186571)(247.3718853,114.487437)(247.3718853,118.15799021)
\curveto(247.3718853,121.82854341)(250.22799958,124.80411471)(253.7511951,124.80411471)
\curveto(257.27439062,124.80411471)(260.1305049,121.82854341)(260.1305049,118.15799021)
\closepath
}
}
{
\newrgbcolor{curcolor}{0 0 0}
\pscustom[linestyle=none,fillstyle=solid,fillcolor=curcolor]
{
\newpath
\moveto(218.4753279,20.25398321)
\curveto(218.4753279,16.58343)(215.61921362,13.60785871)(212.0960181,13.60785871)
\curveto(208.57282258,13.60785871)(205.7167083,16.58343)(205.7167083,20.25398321)
\curveto(205.7167083,23.92453641)(208.57282258,26.90010771)(212.0960181,26.90010771)
\curveto(215.61921362,26.90010771)(218.4753279,23.92453641)(218.4753279,20.25398321)
\closepath
}
}
{
\newrgbcolor{curcolor}{0 0 0}
\pscustom[linestyle=none,fillstyle=solid,fillcolor=curcolor]
{
\newpath
\moveto(250.0615449,66.08143721)
\curveto(250.0615449,62.410884)(247.20543062,59.43531271)(243.6822351,59.43531271)
\curveto(240.15903958,59.43531271)(237.3029253,62.410884)(237.3029253,66.08143721)
\curveto(237.3029253,69.75199041)(240.15903958,72.72756171)(243.6822351,72.72756171)
\curveto(247.20543062,72.72756171)(250.0615449,69.75199041)(250.0615449,66.08143721)
\closepath
}
}
{
\newrgbcolor{curcolor}{0 0 0}
\pscustom[linestyle=none,fillstyle=solid,fillcolor=curcolor]
{
\newpath
\moveto(192.9580999,57.47228821)
\curveto(192.9580999,53.801735)(190.10198562,50.82616371)(186.5787901,50.82616371)
\curveto(183.05559458,50.82616371)(180.1994803,53.801735)(180.1994803,57.47228821)
\curveto(180.1994803,61.14284141)(183.05559458,64.11841271)(186.5787901,64.11841271)
\curveto(190.10198562,64.11841271)(192.9580999,61.14284141)(192.9580999,57.47228821)
\closepath
}
}
{
\newrgbcolor{curcolor}{0 0 0}
\pscustom[linestyle=none,fillstyle=solid,fillcolor=curcolor]
{
\newpath
\moveto(183.3891339,135.23194221)
\curveto(183.3891339,131.561389)(180.53301962,128.58581771)(177.0098241,128.58581771)
\curveto(173.48662858,128.58581771)(170.6305143,131.561389)(170.6305143,135.23194221)
\curveto(170.6305143,138.90249541)(173.48662858,141.87806671)(177.0098241,141.87806671)
\curveto(180.53301962,141.87806671)(183.3891339,138.90249541)(183.3891339,135.23194221)
\closepath
}
}
{
\newrgbcolor{curcolor}{0 0 0}
\pscustom[linestyle=none,fillstyle=solid,fillcolor=curcolor]
{
\newpath
\moveto(219.7511999,181.09020021)
\curveto(219.7511999,177.419647)(216.89508562,174.44407571)(213.3718901,174.44407571)
\curveto(209.84869458,174.44407571)(206.9925803,177.419647)(206.9925803,181.09020021)
\curveto(206.9925803,184.76075341)(209.84869458,187.73632471)(213.3718901,187.73632471)
\curveto(216.89508562,187.73632471)(219.7511999,184.76075341)(219.7511999,181.09020021)
\closepath
}
}
{
\newrgbcolor{curcolor}{0 0 0}
\pscustom[linestyle=none,fillstyle=solid,fillcolor=curcolor]
{
\newpath
\moveto(233.3201749,147.32260321)
\curveto(233.3201749,143.65205)(230.46406062,140.67647871)(226.9408651,140.67647871)
\curveto(223.41766958,140.67647871)(220.5615553,143.65205)(220.5615553,147.32260321)
\curveto(220.5615553,150.99315641)(223.41766958,153.96872771)(226.9408651,153.96872771)
\curveto(230.46406062,153.96872771)(233.3201749,150.99315641)(233.3201749,147.32260321)
\closepath
}
}
{
\newrgbcolor{curcolor}{0 0 0}
\pscustom[linewidth=2.60454035,linecolor=curcolor]
{
\newpath
\moveto(55.165019,185.07787656)
\lineto(204.440859,181.09020156)
}
}
{
\newrgbcolor{curcolor}{0 0 0}
\pscustom[linestyle=none,fillstyle=solid,fillcolor=curcolor]
{
\newpath
\moveto(190.98930092,187.7554037)
\lineto(207.88968189,181.02310479)
\lineto(190.65397378,175.20270048)
\curveto(193.47972185,178.83541435)(193.59944327,183.90567652)(190.98930092,187.7554037)
\closepath
}
}
{
\newrgbcolor{curcolor}{0 0 0}
\pscustom[linewidth=2.60454035,linecolor=curcolor]
{
\newpath
\moveto(29.009841,153.17647656)
\lineto(218.958099,146.66415656)
}
}
{
\newrgbcolor{curcolor}{0 0 0}
\pscustom[linestyle=none,fillstyle=solid,fillcolor=curcolor]
{
\newpath
\moveto(205.55734197,153.43091689)
\lineto(222.4063157,146.57097448)
\lineto(205.12707458,140.88110923)
\curveto(207.98022001,144.49234494)(208.13828994,149.56155648)(205.55734197,153.43091689)
\closepath
}
}
{
\newrgbcolor{curcolor}{0 0 0}
\pscustom[linewidth=2.60454035,linecolor=curcolor]
{
\newpath
\moveto(67.923639,139.88423656)
\lineto(70.475359,139.88423656)
\lineto(168.768439,136.39834656)
}
}
{
\newrgbcolor{curcolor}{0 0 0}
\pscustom[linestyle=none,fillstyle=solid,fillcolor=curcolor]
{
\newpath
\moveto(155.37566324,143.18088973)
\lineto(172.21654352,136.30110217)
\lineto(154.93061112,130.63159768)
\curveto(157.788009,134.23946957)(157.9520509,139.30849137)(155.37566324,143.18088973)
\closepath
}
}
{
\newrgbcolor{curcolor}{0 0 0}
\pscustom[linewidth=2.60454035,linecolor=curcolor]
{
\newpath
\moveto(21.9926,112.63511656)
\lineto(235.802929,68.42298656)
}
}
{
\newrgbcolor{curcolor}{0 0 0}
\pscustom[linestyle=none,fillstyle=solid,fillcolor=curcolor]
{
\newpath
\moveto(223.73691903,77.35499293)
\lineto(239.18591906,67.74899878)
\lineto(221.1941135,65.05796311)
\curveto(224.61597009,68.13568299)(225.62757786,73.10544571)(223.73691903,77.35499293)
\closepath
}
}
{
\newrgbcolor{curcolor}{0 0 0}
\pscustom[linewidth=2.60454035,linecolor=curcolor]
{
\newpath
\moveto(29.647772,61.45996656)
\lineto(205.078789,22.91244656)
}
}
{
\newrgbcolor{curcolor}{0 0 0}
\pscustom[linestyle=none,fillstyle=solid,fillcolor=curcolor]
{
\newpath
\moveto(193.12432408,31.9932011)
\lineto(208.45317252,22.19661336)
\lineto(190.42941929,19.7286061)
\curveto(193.8891299,22.76371148)(194.96220881,27.72056466)(193.12432408,31.9932011)
\closepath
}
}
{
\newrgbcolor{curcolor}{0 0 0}
\pscustom[linewidth=2.60454035,linecolor=curcolor]
{
\newpath
\moveto(54.527079,24.90628656)
\lineto(57.044329,24.89398656)
\lineto(179.561549,55.67034656)
}
}
{
\newrgbcolor{curcolor}{0 0 0}
\pscustom[linestyle=none,fillstyle=solid,fillcolor=curcolor]
{
\newpath
\moveto(164.81158521,58.464609)
\lineto(182.90090008,56.53499426)
\lineto(167.87090428,46.2858008)
\curveto(169.61333423,50.54554259)(168.3624786,55.46054533)(164.81158521,58.464609)
\closepath
}
}
{
\newrgbcolor{curcolor}{0 0 0}
\pscustom[linewidth=2.60454035,linecolor=curcolor]
{
\newpath
\moveto(73.027079,83.39217656)
\lineto(71.113289,80.06911656)
\lineto(245.785689,115.30151656)
}
}
{
\newrgbcolor{curcolor}{0 0 0}
\pscustom[linestyle=none,fillstyle=solid,fillcolor=curcolor]
{
\newpath
\moveto(231.18354994,118.78675245)
\lineto(249.16202752,116.00807122)
\lineto(233.6663987,106.47747691)
\curveto(235.60734331,110.65050893)(234.58915181,115.61892697)(231.18354994,118.78675245)
\closepath
}
}
{
\newrgbcolor{curcolor}{0 0 0}
\pscustom[linestyle=none,fillstyle=solid,fillcolor=curcolor]
{
\newpath
\moveto(14.45320899,195.77087484)
\lineto(8.00376487,195.77087484)
\lineto(6.87357656,193.37858736)
\curveto(6.59512072,192.78982916)(6.45589476,192.35012566)(6.45589828,192.05947557)
\curveto(6.45589476,191.82844356)(6.57669376,191.62536017)(6.81829561,191.45022478)
\curveto(7.05988972,191.27508789)(7.58198706,191.16143572)(8.38458919,191.10926792)
\lineto(8.38458919,190.69564812)
\lineto(3.1390413,190.69564812)
\lineto(3.1390413,191.10926792)
\curveto(3.83517088,191.22105653)(4.2856078,191.36638226)(4.4903534,191.54524555)
\curveto(4.90802972,191.90296957)(5.37075128,192.62959822)(5.87851947,193.72513368)
\lineto(11.73830013,206.20080112)
\lineto(12.16826307,206.20080112)
\lineto(17.96662045,193.59098672)
\curveto(18.43342187,192.57743001)(18.85724206,191.91973793)(19.23808229,191.61790849)
\curveto(19.61888994,191.3160772)(20.14917703,191.14653052)(20.82894518,191.10926792)
\lineto(20.82894518,190.69564812)
\lineto(14.2566545,190.69564812)
\lineto(14.2566545,191.10926792)
\curveto(14.92001392,191.13907792)(15.36840339,191.23968804)(15.60182427,191.41109859)
\curveto(15.83521983,191.58250771)(15.95192395,191.79118055)(15.95193695,192.03711774)
\curveto(15.95192395,192.36503087)(15.7881287,192.88298667)(15.46055074,193.59098672)
\closepath
\moveto(14.10923863,196.59811444)
\lineto(11.28376788,202.72415903)
\lineto(8.38458919,196.59811444)
\closepath
}
}
{
\newrgbcolor{curcolor}{0 0 0}
\pscustom[linestyle=none,fillstyle=solid,fillcolor=curcolor]
{
\newpath
\moveto(174.8489728,199.26146874)
\curveto(176.00371764,199.03788296)(176.86773754,198.68015809)(177.4410351,198.18829305)
\curveto(178.23542782,197.50264705)(178.63263128,196.66422937)(178.63264668,195.67303752)
\curveto(178.63263128,194.92032061)(178.37055889,194.19928142)(177.84642873,193.50991776)
\curveto(177.32226934,192.82055013)(176.60361771,192.31749953)(175.6904717,192.00076444)
\curveto(174.77730076,191.6840284)(173.38299376,191.52566062)(171.50754651,191.52566062)
\lineto(163.64536701,191.52566062)
\lineto(163.64536701,191.93928042)
\lineto(164.27188444,191.93928042)
\curveto(164.96801318,191.93928)(165.46758867,192.14050024)(165.77061241,192.54294174)
\curveto(165.95897439,192.80378178)(166.05315666,193.35900059)(166.05315948,194.20859985)
\lineto(166.05315948,204.00132806)
\curveto(166.05315666,204.94034337)(165.93440511,205.5328252)(165.69690448,205.7787753)
\curveto(165.37750128,206.10667551)(164.90249508,206.27063275)(164.27188444,206.27064749)
\lineto(163.64536701,206.27064749)
\lineto(163.64536701,206.68426729)
\lineto(170.84417512,206.68426729)
\curveto(172.18728849,206.68425213)(173.26424221,206.59482091)(174.0750395,206.41597337)
\curveto(175.30349298,206.14766482)(176.24122074,205.67256147)(176.88822561,204.9906619)
\curveto(177.53520315,204.30873539)(177.85869876,203.52434908)(177.85871339,202.63750061)
\curveto(177.85869876,201.87732414)(177.60481613,201.19727425)(177.09706475,200.5973489)
\curveto(176.58928563,199.99740541)(175.8399224,199.55211247)(174.8489728,199.26146874)
\closepath
\moveto(168.41181333,199.86513007)
\curveto(168.71482935,199.81295352)(169.0608468,199.77382736)(169.44986672,199.74775148)
\curveto(169.8388742,199.72165915)(170.26678927,199.7086171)(170.73361322,199.70862528)
\curveto(171.92931098,199.7086171)(172.82813737,199.82599557)(173.43009509,200.06076106)
\curveto(174.03203241,200.29550947)(174.49270653,200.65509749)(174.81211883,201.13952621)
\curveto(175.13150797,201.62393569)(175.29120833,202.1530704)(175.2912204,202.72693192)
\curveto(175.29120833,203.6137803)(174.89400487,204.37021936)(174.09960882,204.99625136)
\curveto(173.30519102,205.62225641)(172.14633968,205.93526568)(170.62305132,205.93528009)
\curveto(169.80406771,205.93526568)(169.06698912,205.85328706)(168.41181333,205.68934399)
\closepath
\moveto(168.41181333,192.62119414)
\curveto(169.36182056,192.4199728)(170.29954832,192.31936268)(171.22499944,192.31936347)
\curveto(172.70733838,192.31936268)(173.83752556,192.62305619)(174.61556435,193.23044492)
\curveto(175.39358036,193.83783024)(175.78259406,194.58867985)(175.78260662,195.48299599)
\curveto(175.78259406,196.07174755)(175.60651417,196.63814527)(175.25436643,197.18219084)
\curveto(174.90219463,197.72622509)(174.32891128,198.15474968)(173.53451466,198.46776589)
\curveto(172.74009743,198.78076821)(171.75732598,198.93727284)(170.58619735,198.93728025)
\curveto(170.07842474,198.93727284)(169.64436735,198.92982024)(169.28402387,198.91492242)
\curveto(168.92366828,198.90000983)(168.63293173,198.87392573)(168.41181333,198.83667003)
\closepath
}
}
{
\newrgbcolor{curcolor}{0 0 0}
\pscustom[linestyle=none,fillstyle=solid,fillcolor=curcolor]
{
\newpath
\moveto(94.8793099,64.50000021)
\curveto(94.8793099,60.829447)(92.02319562,57.85387571)(88.5000001,57.85387571)
\curveto(84.97680458,57.85387571)(82.1206903,60.829447)(82.1206903,64.50000021)
\curveto(82.1206903,68.17055341)(84.97680458,71.14612471)(88.5000001,71.14612471)
\curveto(92.02319562,71.14612471)(94.8793099,68.17055341)(94.8793099,64.50000021)
\closepath
}
}
{
\newrgbcolor{curcolor}{0 0 0}
\pscustom[linestyle=none,fillstyle=solid,fillcolor=curcolor]
{
\newpath
\moveto(92.3793099,120.00000021)
\curveto(92.3793099,116.329447)(89.52319562,113.35387571)(86.0000001,113.35387571)
\curveto(82.47680458,113.35387571)(79.6206903,116.329447)(79.6206903,120.00000021)
\curveto(79.6206903,123.67055341)(82.47680458,126.64612471)(86.0000001,126.64612471)
\curveto(89.52319562,126.64612471)(92.3793099,123.67055341)(92.3793099,120.00000021)
\closepath
}
}
{
\newrgbcolor{curcolor}{0 0 0}
\pscustom[linewidth=2.60454035,linecolor=curcolor]
{
\newpath
\moveto(87.004679,120.47799656)
\lineto(89.556399,120.47799656)
\lineto(207.849479,174.99210656)
}
}
{
\newrgbcolor{curcolor}{0 0 0}
\pscustom[linestyle=none,fillstyle=solid,fillcolor=curcolor]
{
\newpath
\moveto(192.83718979,175.01462548)
\lineto(210.97173946,176.45852086)
\lineto(198.09280504,163.61018289)
\curveto(199.01867735,168.11842232)(196.88174096,172.7179217)(192.83718979,175.01462548)
\closepath
}
}
{
\newrgbcolor{curcolor}{0 0 0}
\pscustom[linewidth=2.60454035,linecolor=curcolor]
{
\newpath
\moveto(88.004679,64.47799656)
\lineto(90.556399,64.47799656)
\lineto(222.349479,141.99210656)
}
}
{
\newrgbcolor{curcolor}{0 0 0}
\pscustom[linestyle=none,fillstyle=solid,fillcolor=curcolor]
{
\newpath
\moveto(207.40968899,140.51833289)
\lineto(225.31004599,143.76239591)
\lineto(213.77574655,129.69446512)
\curveto(214.24770123,134.27253469)(211.66300109,138.63615898)(207.40968899,140.51833289)
\closepath
}
}
\end{pspicture}

\footnotesize\caption{Pictorial representation of an {\em onto} function.}
\end{minipage}
\hfill
\begin{minipage}{2in}
%LaTeX with PSTricks extensions
%%Creator: 0.48.3.1
%%Please note this file requires PSTricks extensions
\psset{xunit=.5pt,yunit=.5pt,runit=.5pt}
\begin{pspicture}(269.81277466,209.31954956)
{
\newrgbcolor{curcolor}{0.7019608 0.7019608 0.7019608}
\pscustom[linestyle=none,fillstyle=solid,fillcolor=curcolor]
{
\newpath
\moveto(269.8127758,102.66593362)
\curveto(269.8127758,45.9651001)(244.54333166,-0.00000733)(213.3718926,-0.00000733)
\curveto(182.20045354,-0.00000733)(156.9310094,45.9651001)(156.9310094,102.66593362)
\curveto(156.9310094,159.36676713)(182.20045354,205.33187457)(213.3718926,205.33187457)
\curveto(244.54333166,205.33187457)(269.8127758,159.36676713)(269.8127758,102.66593362)
\closepath
}
}
{
\newrgbcolor{curcolor}{0.7019608 0.7019608 0.7019608}
\pscustom[linestyle=none,fillstyle=solid,fillcolor=curcolor]
{
\newpath
\moveto(112.8817668,106.65360862)
\curveto(112.8817668,49.9527751)(87.61232266,3.98766767)(56.4408836,3.98766767)
\curveto(25.26944454,3.98766767)(0.0000004,49.9527751)(0.0000004,106.65360862)
\curveto(0.0000004,163.35444213)(25.26944454,209.31954957)(56.4408836,209.31954957)
\curveto(87.61232266,209.31954957)(112.8817668,163.35444213)(112.8817668,106.65360862)
\closepath
}
}
{
\newrgbcolor{curcolor}{0 0 0}
\pscustom[linestyle=none,fillstyle=solid,fillcolor=curcolor]
{
\newpath
\moveto(34.7512229,153.17647821)
\curveto(34.7512229,149.505925)(31.89510862,146.53035371)(28.3719131,146.53035371)
\curveto(24.84871758,146.53035371)(21.9926033,149.505925)(21.9926033,153.17647821)
\curveto(21.9926033,156.84703141)(24.84871758,159.82260271)(28.3719131,159.82260271)
\curveto(31.89510862,159.82260271)(34.7512229,156.84703141)(34.7512229,153.17647821)
\closepath
}
}
{
\newrgbcolor{curcolor}{0 0 0}
\pscustom[linestyle=none,fillstyle=solid,fillcolor=curcolor]
{
\newpath
\moveto(61.5443199,24.24166321)
\curveto(61.5443199,20.57111)(58.68820562,17.59553871)(55.1650101,17.59553871)
\curveto(51.64181458,17.59553871)(48.7857003,20.57111)(48.7857003,24.24166321)
\curveto(48.7857003,27.91221641)(51.64181458,30.88778771)(55.1650101,30.88778771)
\curveto(58.68820562,30.88778771)(61.5443199,27.91221641)(61.5443199,24.24166321)
\closepath
}
}
{
\newrgbcolor{curcolor}{0 0 0}
\pscustom[linestyle=none,fillstyle=solid,fillcolor=curcolor]
{
\newpath
\moveto(78.1305369,80.06910921)
\curveto(78.1305369,76.398556)(75.27442262,73.42298471)(71.7512271,73.42298471)
\curveto(68.22803158,73.42298471)(65.3719173,76.398556)(65.3719173,80.06910921)
\curveto(65.3719173,83.73966241)(68.22803158,86.71523371)(71.7512271,86.71523371)
\curveto(75.27442262,86.71523371)(78.1305369,83.73966241)(78.1305369,80.06910921)
\closepath
}
}
{
\newrgbcolor{curcolor}{0 0 0}
\pscustom[linestyle=none,fillstyle=solid,fillcolor=curcolor]
{
\newpath
\moveto(36.0270919,61.45996021)
\curveto(36.0270919,57.789407)(33.17097762,54.81383571)(29.6477821,54.81383571)
\curveto(26.12458658,54.81383571)(23.2684723,57.789407)(23.2684723,61.45996021)
\curveto(23.2684723,65.13051341)(26.12458658,68.10608471)(29.6477821,68.10608471)
\curveto(33.17097762,68.10608471)(36.0270919,65.13051341)(36.0270919,61.45996021)
\closepath
}
}
{
\newrgbcolor{curcolor}{0 0 0}
\pscustom[linestyle=none,fillstyle=solid,fillcolor=curcolor]
{
\newpath
\moveto(29.6477809,113.29973121)
\curveto(29.6477809,109.629178)(26.79166662,106.65360671)(23.2684711,106.65360671)
\curveto(19.74527558,106.65360671)(16.8891613,109.629178)(16.8891613,113.29973121)
\curveto(16.8891613,116.97028441)(19.74527558,119.94585571)(23.2684711,119.94585571)
\curveto(26.79166662,119.94585571)(29.6477809,116.97028441)(29.6477809,113.29973121)
\closepath
}
}
{
\newrgbcolor{curcolor}{0 0 0}
\pscustom[linestyle=none,fillstyle=solid,fillcolor=curcolor]
{
\newpath
\moveto(76.8546749,139.88422911)
\curveto(76.8546749,136.2136759)(73.99856062,133.23810461)(70.4753651,133.23810461)
\curveto(66.95216958,133.23810461)(64.0960553,136.2136759)(64.0960553,139.88422911)
\curveto(64.0960553,143.55478231)(66.95216958,146.53035361)(70.4753651,146.53035361)
\curveto(73.99856062,146.53035361)(76.8546749,143.55478231)(76.8546749,139.88422911)
\closepath
}
}
{
\newrgbcolor{curcolor}{0 0 0}
\pscustom[linestyle=none,fillstyle=solid,fillcolor=curcolor]
{
\newpath
\moveto(62.8201919,185.07787621)
\curveto(62.8201919,181.407323)(59.96407762,178.43175171)(56.4408821,178.43175171)
\curveto(52.91768658,178.43175171)(50.0615723,181.407323)(50.0615723,185.07787621)
\curveto(50.0615723,188.74842941)(52.91768658,191.72400071)(56.4408821,191.72400071)
\curveto(59.96407762,191.72400071)(62.8201919,188.74842941)(62.8201919,185.07787621)
\closepath
}
}
{
\newrgbcolor{curcolor}{0 0 0}
\pscustom[linestyle=none,fillstyle=solid,fillcolor=curcolor]
{
\newpath
\moveto(254.1305049,124.15799021)
\curveto(254.1305049,120.487437)(251.27439062,117.51186571)(247.7511951,117.51186571)
\curveto(244.22799958,117.51186571)(241.3718853,120.487437)(241.3718853,124.15799021)
\curveto(241.3718853,127.82854341)(244.22799958,130.80411471)(247.7511951,130.80411471)
\curveto(251.27439062,130.80411471)(254.1305049,127.82854341)(254.1305049,124.15799021)
\closepath
}
}
{
\newrgbcolor{curcolor}{0 0 0}
\pscustom[linestyle=none,fillstyle=solid,fillcolor=curcolor]
{
\newpath
\moveto(218.4753279,20.25398321)
\curveto(218.4753279,16.58343)(215.61921362,13.60785871)(212.0960181,13.60785871)
\curveto(208.57282258,13.60785871)(205.7167083,16.58343)(205.7167083,20.25398321)
\curveto(205.7167083,23.92453641)(208.57282258,26.90010771)(212.0960181,26.90010771)
\curveto(215.61921362,26.90010771)(218.4753279,23.92453641)(218.4753279,20.25398321)
\closepath
}
}
{
\newrgbcolor{curcolor}{0 0 0}
\pscustom[linestyle=none,fillstyle=solid,fillcolor=curcolor]
{
\newpath
\moveto(258.0615449,77.08143721)
\curveto(258.0615449,73.410884)(255.20543062,70.43531271)(251.6822351,70.43531271)
\curveto(248.15903958,70.43531271)(245.3029253,73.410884)(245.3029253,77.08143721)
\curveto(245.3029253,80.75199041)(248.15903958,83.72756171)(251.6822351,83.72756171)
\curveto(255.20543062,83.72756171)(258.0615449,80.75199041)(258.0615449,77.08143721)
\closepath
}
}
{
\newrgbcolor{curcolor}{0 0 0}
\pscustom[linestyle=none,fillstyle=solid,fillcolor=curcolor]
{
\newpath
\moveto(192.9580999,57.47228821)
\curveto(192.9580999,53.801735)(190.10198562,50.82616371)(186.5787901,50.82616371)
\curveto(183.05559458,50.82616371)(180.1994803,53.801735)(180.1994803,57.47228821)
\curveto(180.1994803,61.14284141)(183.05559458,64.11841271)(186.5787901,64.11841271)
\curveto(190.10198562,64.11841271)(192.9580999,61.14284141)(192.9580999,57.47228821)
\closepath
}
}
{
\newrgbcolor{curcolor}{0 0 0}
\pscustom[linestyle=none,fillstyle=solid,fillcolor=curcolor]
{
\newpath
\moveto(221.8891339,101.23194221)
\curveto(221.8891339,97.561389)(219.03301962,94.58581771)(215.5098241,94.58581771)
\curveto(211.98662858,94.58581771)(209.1305143,97.561389)(209.1305143,101.23194221)
\curveto(209.1305143,104.90249541)(211.98662858,107.87806671)(215.5098241,107.87806671)
\curveto(219.03301962,107.87806671)(221.8891339,104.90249541)(221.8891339,101.23194221)
\closepath
}
}
{
\newrgbcolor{curcolor}{0 0 0}
\pscustom[linestyle=none,fillstyle=solid,fillcolor=curcolor]
{
\newpath
\moveto(219.7511999,181.09020021)
\curveto(219.7511999,177.419647)(216.89508562,174.44407571)(213.3718901,174.44407571)
\curveto(209.84869458,174.44407571)(206.9925803,177.419647)(206.9925803,181.09020021)
\curveto(206.9925803,184.76075341)(209.84869458,187.73632471)(213.3718901,187.73632471)
\curveto(216.89508562,187.73632471)(219.7511999,184.76075341)(219.7511999,181.09020021)
\closepath
}
}
{
\newrgbcolor{curcolor}{0 0 0}
\pscustom[linestyle=none,fillstyle=solid,fillcolor=curcolor]
{
\newpath
\moveto(247.8201749,159.82260321)
\curveto(247.8201749,156.15205)(244.96406062,153.17647871)(241.4408651,153.17647871)
\curveto(237.91766958,153.17647871)(235.0615553,156.15205)(235.0615553,159.82260321)
\curveto(235.0615553,163.49315641)(237.91766958,166.46872771)(241.4408651,166.46872771)
\curveto(244.96406062,166.46872771)(247.8201749,163.49315641)(247.8201749,159.82260321)
\closepath
}
}
{
\newrgbcolor{curcolor}{0 0 0}
\pscustom[linewidth=2.60454035,linecolor=curcolor]
{
\newpath
\moveto(55.165019,185.07787656)
\lineto(204.440859,181.09020156)
}
}
{
\newrgbcolor{curcolor}{0 0 0}
\pscustom[linestyle=none,fillstyle=solid,fillcolor=curcolor]
{
\newpath
\moveto(190.98930092,187.7554037)
\lineto(207.88968189,181.02310479)
\lineto(190.65397378,175.20270048)
\curveto(193.47972185,178.83541435)(193.59944327,183.90567652)(190.98930092,187.7554037)
\closepath
}
}
{
\newrgbcolor{curcolor}{0 0 0}
\pscustom[linewidth=2.60454035,linecolor=curcolor]
{
\newpath
\moveto(29.009841,153.17647656)
\lineto(237.958099,125.66415656)
}
}
{
\newrgbcolor{curcolor}{0 0 0}
\pscustom[linestyle=none,fillstyle=solid,fillcolor=curcolor]
{
\newpath
\moveto(225.27284134,133.69245479)
\lineto(241.38123258,125.23867118)
\lineto(223.63357914,121.24273103)
\curveto(226.82187886,124.56179823)(227.46836835,129.59210071)(225.27284134,133.69245479)
\closepath
}
}
{
\newrgbcolor{curcolor}{0 0 0}
\pscustom[linewidth=2.60454035,linecolor=curcolor]
{
\newpath
\moveto(67.923639,139.88423656)
\lineto(70.475359,139.88423656)
\lineto(207.768439,102.89834656)
}
}
{
\newrgbcolor{curcolor}{0 0 0}
\pscustom[linestyle=none,fillstyle=solid,fillcolor=curcolor]
{
\newpath
\moveto(196.25240311,112.52903977)
\lineto(211.10559233,102.02525495)
\lineto(192.98602626,100.40412391)
\curveto(196.58409351,103.27385809)(197.88815597,108.17501232)(196.25240311,112.52903977)
\closepath
}
}
{
\newrgbcolor{curcolor}{0 0 0}
\pscustom[linewidth=2.60454035,linecolor=curcolor]
{
\newpath
\moveto(21.9926,112.63511656)
\lineto(243.802929,77.92298656)
}
}
{
\newrgbcolor{curcolor}{0 0 0}
\pscustom[linestyle=none,fillstyle=solid,fillcolor=curcolor]
{
\newpath
\moveto(231.31663631,86.25736291)
\lineto(247.21470444,77.41439084)
\lineto(229.37513411,73.85118011)
\curveto(232.64319683,77.09173998)(233.41181087,82.10483525)(231.31663631,86.25736291)
\closepath
}
}
{
\newrgbcolor{curcolor}{0 0 0}
\pscustom[linewidth=2.60454035,linecolor=curcolor]
{
\newpath
\moveto(29.647772,61.45996656)
\lineto(205.078789,22.91244656)
}
}
{
\newrgbcolor{curcolor}{0 0 0}
\pscustom[linestyle=none,fillstyle=solid,fillcolor=curcolor]
{
\newpath
\moveto(193.12432408,31.9932011)
\lineto(208.45317252,22.19661336)
\lineto(190.42941929,19.7286061)
\curveto(193.8891299,22.76371148)(194.96220881,27.72056466)(193.12432408,31.9932011)
\closepath
}
}
{
\newrgbcolor{curcolor}{0 0 0}
\pscustom[linewidth=2.60454035,linecolor=curcolor]
{
\newpath
\moveto(54.527079,24.90628656)
\lineto(61.544329,28.89395656)
\lineto(235.061549,155.17031656)
}
}
{
\newrgbcolor{curcolor}{0 0 0}
\pscustom[linestyle=none,fillstyle=solid,fillcolor=curcolor]
{
\newpath
\moveto(220.33601708,152.25003507)
\lineto(237.83583884,157.2202427)
\lineto(227.72493355,142.09687222)
\curveto(227.74924669,146.69914026)(224.75226919,150.79059692)(220.33601708,152.25003507)
\closepath
}
}
{
\newrgbcolor{curcolor}{0 0 0}
\pscustom[linewidth=2.60454035,linecolor=curcolor]
{
\newpath
\moveto(73.027079,83.39217656)
\lineto(71.113289,80.06911656)
\lineto(178.285689,58.80151656)
}
}
{
\newrgbcolor{curcolor}{0 0 0}
\pscustom[linestyle=none,fillstyle=solid,fillcolor=curcolor]
{
\newpath
\moveto(166.14851925,67.63658559)
\lineto(181.67396929,58.15464873)
\lineto(163.70429659,55.31958198)
\curveto(167.10139034,58.42461278)(168.07315702,63.4023192)(166.14851925,67.63658559)
\closepath
}
}
{
\newrgbcolor{curcolor}{0 0 0}
\pscustom[linestyle=none,fillstyle=solid,fillcolor=curcolor]
{
\newpath
\moveto(14.95321286,197.27087979)
\lineto(8.50376874,197.27087979)
\lineto(7.37358044,194.87859231)
\curveto(7.09512459,194.28983411)(6.95589864,193.85013062)(6.95590215,193.55948052)
\curveto(6.95589864,193.32844851)(7.07669763,193.12536512)(7.31829949,192.95022974)
\curveto(7.55989359,192.77509284)(8.08199093,192.66144067)(8.88459306,192.60927287)
\lineto(8.88459306,192.19565308)
\lineto(3.63904517,192.19565308)
\lineto(3.63904517,192.60927287)
\curveto(4.33517476,192.72106148)(4.78561167,192.86638721)(4.99035727,193.0452505)
\curveto(5.4080336,193.40297452)(5.87075516,194.12960317)(6.37852334,195.22513863)
\lineto(12.238304,207.70080607)
\lineto(12.66826694,207.70080607)
\lineto(18.46662432,195.09099167)
\curveto(18.93342574,194.07743496)(19.35724593,193.41974288)(19.73808617,193.11791344)
\curveto(20.11889381,192.81608215)(20.64918091,192.64653547)(21.32894905,192.60927287)
\lineto(21.32894905,192.19565308)
\lineto(14.75665837,192.19565308)
\lineto(14.75665837,192.60927287)
\curveto(15.42001779,192.63908287)(15.86840727,192.73969299)(16.10182815,192.91110354)
\curveto(16.33522371,193.08251266)(16.45192782,193.2911855)(16.45194083,193.53712269)
\curveto(16.45192782,193.86503582)(16.28813258,194.38299163)(15.96055461,195.09099167)
\closepath
\moveto(14.60924251,198.09811939)
\lineto(11.78377175,204.22416398)
\lineto(8.88459306,198.09811939)
\closepath
}
}
{
\newrgbcolor{curcolor}{0 0 0}
\pscustom[linestyle=none,fillstyle=solid,fillcolor=curcolor]
{
\newpath
\moveto(171.34899768,198.26147272)
\curveto(172.50374252,198.03788694)(173.36776242,197.68016206)(173.94105998,197.18829703)
\curveto(174.73545269,196.50265102)(175.13265616,195.66423335)(175.13267156,194.67304149)
\curveto(175.13265616,193.92032459)(174.87058377,193.19928539)(174.34645361,192.50992174)
\curveto(173.82229422,191.82055411)(173.10364259,191.31750351)(172.19049658,191.00076842)
\curveto(171.27732564,190.68403238)(169.88301864,190.52566459)(168.00757139,190.52566459)
\lineto(160.14539189,190.52566459)
\lineto(160.14539189,190.93928439)
\lineto(160.77190932,190.93928439)
\curveto(161.46803805,190.93928398)(161.96761354,191.14050422)(162.27063729,191.54294572)
\curveto(162.45899927,191.80378576)(162.55318154,192.35900457)(162.55318436,193.20860383)
\lineto(162.55318436,203.00133203)
\curveto(162.55318154,203.94034735)(162.43442999,204.53282917)(162.19692935,204.77877928)
\curveto(161.87752616,205.10667949)(161.40251996,205.27063672)(160.77190932,205.27065147)
\lineto(160.14539189,205.27065147)
\lineto(160.14539189,205.68427127)
\lineto(167.34419999,205.68427127)
\curveto(168.68731337,205.68425611)(169.76426709,205.59482489)(170.57506438,205.41597734)
\curveto(171.80351785,205.1476688)(172.74124562,204.67256545)(173.38825049,203.99066588)
\curveto(174.03522803,203.30873937)(174.35872364,202.52435306)(174.35873827,201.63750459)
\curveto(174.35872364,200.87732812)(174.10484101,200.19727823)(173.59708963,199.59735288)
\curveto(173.08931051,198.99740939)(172.33994727,198.55211644)(171.34899768,198.26147272)
\closepath
\moveto(164.91183821,198.86513405)
\curveto(165.21485423,198.8129575)(165.56087168,198.77383134)(165.9498916,198.74775546)
\curveto(166.33889908,198.72166313)(166.76681415,198.70862108)(167.2336381,198.70862926)
\curveto(168.42933586,198.70862108)(169.32816225,198.82599955)(169.93011997,199.06076503)
\curveto(170.53205728,199.29551345)(170.9927314,199.65510147)(171.31214371,200.13953018)
\curveto(171.63153285,200.62393967)(171.79123321,201.15307438)(171.79124527,201.7269359)
\curveto(171.79123321,202.61378428)(171.39402975,203.37022333)(170.59963369,203.99625533)
\curveto(169.8052159,204.62226039)(168.64636456,204.93526966)(167.1230762,204.93528407)
\curveto(166.30409259,204.93526966)(165.567014,204.85329104)(164.91183821,204.68934797)
\closepath
\moveto(164.91183821,191.62119811)
\curveto(165.86184543,191.41997678)(166.7995732,191.31936666)(167.72502431,191.31936745)
\curveto(169.20736326,191.31936666)(170.33755043,191.62306017)(171.11558922,192.2304489)
\curveto(171.89360524,192.83783422)(172.28261894,193.58868382)(172.28263149,194.48299996)
\curveto(172.28261894,195.07175153)(172.10653905,195.63814924)(171.75439131,196.18219481)
\curveto(171.40221951,196.72622907)(170.82893616,197.15475366)(170.03453954,197.46776986)
\curveto(169.24012231,197.78077219)(168.25735085,197.93727682)(167.08622223,197.93728423)
\curveto(166.57844962,197.93727682)(166.14439223,197.92982422)(165.78404875,197.9149264)
\curveto(165.42369316,197.90001381)(165.1329566,197.8739297)(164.91183821,197.83667401)
\closepath
}
}
\end{pspicture}

\footnotesize\caption{Pictorial representation of an {\em bijective} function.}
\end{minipage}
\end{figure}

\begin{exa}
For each of the following functions from
$\integers$ to $\integers$,
we determine whether or not they are
one-to-one and onto.
\begin{enumerate}[(a)]
\item
Let $f(x)=x+2$.
Notice that if $f(a)=f(b)$, then $a+2=b+2$ so $a=b$.  Thus, $f$ is one-to-one.
Also notice that for any $b\in\integers$, $f(b-2)=b-2+2=b$, so $f$ is onto.
\item
Let $g(x)=x^2$.
Since $g(1)=g(-1)=1$, $g$ is not one-to-one.  Also notice that there is no integer $a$ such that
$g(a)=a^2=5$, so $g$ is not onto.
\item
Let $h(x)=2x$.
If $h(a)=h(b)$, then $2a=2b$ so $a=b$.  Thus, $h$ is one-to-one.
But there is no integer $a$ such that $h(a)=2a=3$, so $a$ is not onto.
\item
Let $r(x)=\lfloor x/2\rfloor$.
Notice that $r(0)=\lfloor 0/2\rfloor=\lfloor 0\rfloor = 0$ and 
$r(1)=\lfloor 1/2\rfloor=\lfloor 0\rfloor = 0$, so $r$ is not one-to-one.
But for any integer $b$, $r(2b)=\lfloor 2b/2\rfloor=\lfloor b\rfloor = b$, so $r$ is onto. 
\end{enumerate}
\end{exa}

The functions in the previous exercise were specifically chosen to demonstrate that 
all four possibilities of being or not being one-to-one and onto 
(one-to-one and onto, one-to-one and not onto, not one-to-one but onto, and not one-to-one or onto)
are possible.



The following theorem should come as no surprise if you take a few minutes to think about it
(and you {\em should} take a few minutes to think about it until you are convinced it is correct).

\begin{thm}\label{thm:size_domain_image_injections_surjections}
Let $f:A\to B$ be a function, and let $A$ and $B$ be finite.
\begin{enumerate}
\item
If $f$ is one-to-one, then $|A| \leq |B|$. 
\item
If $f$ is onto, then $|A|\geq |B|$. 
\item
If $f$ is bijective, then $|A| = |B|$.
\end{enumerate}
\end{thm}

\newpage
\begin{exer}
Let's test your understanding of the material so far. Answer each of the following
true/false questions, giving a very brief justification/counterexample.
\begin{enumerate}[(a)]
\itemtf 
If $f:A\to B$ is onto, then the domain and range are not only the same size, but they are the same set.
\bigskip
\itemtf
If $f:A\to A$, then $f$ must be one-to-one and onto.
\bigskip
\itemtf
If $f:A\to B$ is both one-to-one and onto, then $A$ and $B$ have the same number of elements.
\bigskip
\itemtf
Let $f(1)=2$ and $f(1)=3$.  Then $f$ is a valid function.
\bigskip
\itemtf
Let $f:\reals\to \reals$ be defined by $f(x)=x^3$.  Then $f$ is one-to-one and onto.
\bigskip
\itemtf
Let $f:\reals^{+}\to \reals$ be defined by $f(x)=\sqrt{x}$.  Then $f$ is a function that is neither 
one-to-one nor onto.
\bigskip
\itemtf
The range of a function is always a subset of the codomain.
\bigskip
\itemtf
A function that is one-to-one is guaranteed to be onto.
\bigskip
\itemtf
Let $a,b\in\integers$, with $a\neq 0$, and define $f:\integers\to\integers$ by $f(x)=ax+b$.
Then $f$ is one-to-one and onto.
\bigskip
\itemtf
Let $a,b\in\integers$, with $a\neq 0$, and define $f:\naturals\to\naturals$  by $f(x)=ax+b$.
Then $f$ is one-to-one and onto.
\bigskip
\itemtf
Let $a,b\in\reals$, with $a\neq 0$, and define $f:\reals\to\reals$ by $f(x)=ax+b$.
Then $f$ is one-to-one and onto.
\end{enumerate}
%\vspace*{.25in}
\begin{answer}
(a) F.  Consider $f(x)=\lfloor x \rfloor$ from $\reals$ to $\integers$.
(b) F.  Consider $f(x)=x^2$ from $\reals$ to $\reals$ which is not one-to-one.
(c) T.  See Theorem~\ref{thm:size_domain_image_injections_surjections}.
(d) F.  $f$ maps $1$ to two different values, so it isn't a function.
(e) T.  We previously showed it was onto, and it isn't difficult to see that it is one-to-one.
(f) F. $f$ is not onto, but it is one-to-one.
(g) T.  By definition of range, it is a subset of the codomain.
(h) F.  We have seen several counter examples to this.
(i) F.  If $a=2$ and $b=0$, the odd numbers are not in the range.
(j) F.  Same counterexample as the previous question.
(k) T.  The proof is similar to several previous proofs.
\end{answer}
\end{exer}


\begin{df}
Let $f$ be a one-to-one correspondence from $A$ to $B$.
The {\bf inverse} of $f$, denoted by $f^{-1}$, is the function such that
$f^{-1}(b)=a$ whenever $f(a)=b$.

\noindent
A function that has an inverse is called {\bf invertible}.
Said another way, a function is {\bf invertible} if and only if it is one-to-one and onto.
\end{df}
\begin{rem}
It is important to note that the function $f^{-1}$ is not the same thing as $1/f$.
This is an unfortunate case when a notation can be interpreted in two different ways.
That is, in some cases, $a^{-1}$ means the inverse function and in other cases it means $1/a$.
Usually the context will help you determine which one is the correct interpretation.
\end{rem}

\begin{proc}
One method of finding the inverse of a function is to replace $f(x)$ (or whatever the name of the
function is) with $y$ and solve for $x$ (or whatever the variable is).
Finally, replace $y$ with $x$ and you have the inverse.
However, it is important to note that this only works if
$f$ is a one-to-one correspondence, so you typically need to verify
that first.
\end{proc}

\begin{exa}
Let $f:\integers\to\integers$ be defined by $f(x)=x+2$.
Notice that $f$ is a one-to-one correspondence, so it has an inverse.
We let $y=x+2$.  Solving for $x$, we get $x=y-2$.
Thus, $f^{-1}(x)=x-2$.
\end{exa}

\begin{exa}
Let $f:\reals\to\reals$ be defined by $f(x)=x^2$.  Then $f$ does not have an inverse since it is not one-to-one.
\end{exa}

\begin{exa}
Let $f:\reals\to\reals$ be defined by $f(x)=x^3$.
We leave it to the reader to prove that $f$ is one-to-one and onto.
Given that, we can find it's inverse.

Let $y=x^3$.  Taking the third root of both sides, we obtain
$\sqrt[3]{y}=\sqrt[3]{x^3}=x$.  Or $x=\sqrt[3]{y}$.
Thus, the inverse of $f$ is given by $f^{-1}(x)=\sqrt[3]{x}$.
\end{exa}
Notice that the previous
example works 
because $\sqrt[3]{x^3}=x$ (similarly for any odd power).
However,
it does {\em not} work for squares since $\sqrt{x^2}=|x|$
(similar for any even power). The fact that the absolute value
shows up should clue you into the fact that $x^2$ is not one-to-one,
so it can't be invertible.

\begin{exer}
Let $f(x)=7x+2$ be a function over $\reals$. 
You can assume that $f$ is a one-to-one correspondence.
Find $f^{-1}$.
\vspace*{1.25in}
\begin{answer}
Let $y=7x+2$.  Then $7x=y-2$, so $x=(y-2)/7$.  Thus, $f^{-1}(x)=(x-2)/7$
(or $\frac{x}{7}-\frac{2}{7}$).
\end{answer}
\end{exer}

\begin{df}
Let $g$ be a function from $A$ to $B$ and $f$ a function from $B$ to $C$.
The {\bf composition} of $f$ and $g$, denoted by $f\circ g$, is defined as
$(f\circ g)(x)=f(g(x))$ for any $x\in A$.
\end{df}

In other words, to compose $f$ with $g$, we {\em first} compute $g(x)$.  
Then we plug in $g(x)$ into the formula for $f$.

\begin{rem}
Look closely at the notation.  $f\circ g$ has $f$ before $g$, so it might seem like it should be
$g(f(x))$--in other words, apply $f$ first, then then $g$. But that is {\em not} how it is defined.

Also notice that to compose $f$ with $g$, it is necessary that the range of $g$ is a subset of the domain
of $f$ since otherwise it would be impossible to compute.
\end{rem}

\begin{exa}
Let $f$ and $g$ be functions on $\integers$ defined by $f(x)=x^2$ and $g(x)=2x-5$.  Compute
$f\circ g$ and $g\circ f$, simplifying your answers.
\begin{sol}
\begin{eqnarray*}
(f\circ g)(x) &=& f(g(x)) = f(2x-5) = (2x-5)^2 = 4x^2-20x+25.\\
(g\circ f)(x) &=& g(f(x)) = g(x^2) = 2x^2-5.
\end{eqnarray*}
\end{sol}
\end{exa}

Notice that in the previous example, $f\circ g \neq g\circ f$.  In other words, the order in which we
compose functions matters since the result 
usually not the same (although occasionally it is).

\begin{exer}
Let $f$ and $g$ be functions on $\reals$ defined by $f(x)=\lfloor x \rfloor$ and $g(x)=x/2$.  Compute
$f\circ g$ and $g\circ f$, simplifying your answers.
\begin{eqnarray*}
(f\circ g)(x) &=& \blankline{3}\\
(g\circ f)(x) &=& \blankline{3}
\end{eqnarray*}
\begin{answer}
$(f\circ g)(x)=f(x/2)=\lfloor x/2\rfloor$, and
$(g\circ f)(x)=g(\lfloor x\rfloor) = (\lfloor x\rfloor)/2$.
\end{answer}
\end{exer}



\begin{df}
We define the {\bf identity function}, $\iota_A:A\to A$,
by $\iota_A(x)=x$. 

\noindent
The subscript can be omitted if the domain/codomain is clear.
\end{df}

\begin{thm}\label{thm:inveIdent}
Let $f$ be an invertible function from $A$ to $B$.  Then $f\circ f^{-1}=\iota_B$
and $f^{-1}\circ f=\iota_A$.
\begin{pf}
Let $a\in A$ and define $b=f(a)$. Then by definition, $f^{-1}(b)=a$, 
so $(f^{-1}\circ f)(a)=f^{-1}(f(a))=f^{-1}(b)=a$.
Thus, $f^{-1}\circ f=\iota_A$.

Conversely, if $b\in B$ and we define $a=f^{-1}(b)$, then 
$(f\circ f^{-1})(b)=f(f^{-1}(b))=f(a)=b$.
Thus,  $f\circ f^{-1}=\iota_B$.
\end{pf}
\end{thm}

\begin{exa}
Prove or disprove that $f(x)=2x+1$ and $g(x)=2x-1$, defined over the real numbers, are inverses.
\begin{sol}
Notice that $(f\circ g)(x)=f(2x-1)=2(2x-1)+1=4x-1\neq x$.
According to Theorem~\ref{thm:inveIdent}, this implies that $f$ and $g$ are not inverses.

\end{sol}
\end{exa}

\begin{exer}
Let's test your understanding of the material so far. Answer each of the following
true/false questions, giving a very brief justification/counterexample.
\begin{enumerate}[(a)]
\itemtf
Let $a,b\in\integers$, with $a\neq 0$, and define $f:\integers\to\integers$ by $f(x)=ax+b$.
Then $f$ is invertible.
\bigskip
\itemtf
Let $a,b\in\integers$, with $a\neq 0$, and define $f:\naturals\to\naturals$ by $f(x)=ax+b$.
Then $f$ is invertible.
\bigskip
\itemtf
Let $a,b\in\reals$, with $a\neq 0$, and define $f:\reals\to\reals$ by $f(x)=ax+b$.
Then $f$ is invertible.
\bigskip
\itemtf
If $f(x)=x^2$, then $f^{-1}(x)=1/x^2$.
\bigskip
\itemtf
Let $n$ be a positive integer. Then the function $\sqrt[n]{x}$ is invertible on $\reals$.
\bigskip
\itemtf
Let $n$ be a positive integer. Then the function $\sqrt[n]{x}$ is invertible on $\naturals$.
\bigskip
\itemtf
Let $n$ be a positive integer. Then the function $\sqrt[n]{x}$ is invertible on $\reals^{+}$ 
(the positive real numbers).
\bigskip
\itemtf
Let $f$ and $g$ be functions on $\BBZ^+$ defined by $f(x)=x^2$ and $g(x)=1/x$.
Then $f\circ g = g\circ f$.
\bigskip
\itemtf
Let $f$ and $g$ be functions on $\integers$ defined by $f(x)=(x+1)^2$ and $g(x)=x+1$.
Then $f\circ g = g\circ f$.
\bigskip
\itemtf
Let $f(x)=\lfloor x\rfloor$ and $g(x)=\lceil x \rceil$ be defined on the real numbers.
Then $f\circ g$ = $g\circ f$.
\bigskip
\itemtf
Let $f(x)=\lfloor x\rfloor$ and $g(x)=\lceil x \rceil$ be defined on the real numbers.
Then $f$ and $g$ are inverses of each other.
\bigskip
\itemtf
Let $f(x)=x^2$ and $g(x)=\sqrt{x}$ be defined over the positive real numbers.  Then $f$ and $g$ are
inverses of each other.
\end{enumerate}
\vspace*{.25in}
\begin{answer}
(a) F.  $f$ might not be onto--e.g. if $a=2$ and $b=0$.
(b) F. Same reason as the previous question.
(c) T. Since over the reals, $f$ is one-to-one and onto.
(d) F. There are several problems.  First, $x^2$ may not even have an inverse depending on 
the domain (which was not specified).  Second, even if it had an inverse, it certainly wouldn't
be $1/x^2$.  That's its reciprocal, not its inverse.  Its inverse would be $\sqrt{x}$ (again,
assuming the domain was chosen so that it is invertible).
(e) F.  This is only true if $n$ is odd.
(f) F.  $\sqrt{2}\not\in\naturals$, so not only is it not invertible, it can't even be defined on $\naturals$.
(g) T.  The $n$th root of a positive number is defined for all positive real numbers, so the function is
well defined.  It is not too difficult to convince yourself that the function is both one-to-one and onto
when restricted to positive numbers, so it is invertible.
(h) T.  In both cases you get $1/x^2$.
(i) F.  $(f\circ g)(x)=f(x+1)=(x+1+1)^2=(x+2)^2=x^2+4x+4$, and $(g\circ f)(x)=g((x+1)^2)=(x+1)^2+1=x^2+2x+2$,
which are clearly not the same.
(j) F. $(f\circ g)(x)=\lceil x\rceil$, and $(g\circ f)(x)=\lfloor x\rfloor$. 
(We'll leave it to you to see why this is the case.)
(k) F.  Certainly not.  $f(3.5)=3$, but $g(3)=3$, not $3.5$.
(l) T. With the restricted domain, they are indeed inverses.
\end{answer}
\end{exer}


\newpage
\subsection{Function Proofs}\label{section:functionProofs}

In this section we give more in depth examples
of proving things about functions.

\begin{proc}
To show that a function $f$ is one-to-one, you just need to show that whenever $f(a)=f(b)$, then $a=b$.
\end{proc}

\begin{exa}
Let $f(x)=2x-3$ be a function on the integers.  Show that $f$ is one-to-one.
\begin{sol}  
Let $a,b\in\integers$ and assume that $f(a)=f(b)$.  Then $2a-3=2b-3$.  Adding $3$
to both sides, we get $2a=2b$.  Dividing both sides by two, we obtain $a=b$.  Therefore, $f(x)=2x-3$
is one-to-one.
\end{sol}
\end{exa}

\begin{ques}
Previously we mentioned that `working both sides' was not an appropriate proof technique.
Why is it O.K. in the previous example?

\noindent
\answerlinetwo
\begin{answer}
We never said it was {\em always} wrong to work both sides of an equation.  If you are working
on an equation that you know to be true, there is absolutely nothing wrong with it.
It is a problem only when you are starting with something you don't know to be true.
In this case, we know that $2a-3=2b-3$ is true given the assumption made.  Therefore, we are
free to `work both sides'.
\end{answer}
\end{ques}

\begin{exer}
Prove that $f(x)=5x$ is one-to-one over the real numbers.

\noindent
\pflinethree
\begin{answer}
Let $a,b\in \BBR$.  If $f(a)=f(b)$, then $5a=5b$.  Dividing both sides by $5$, we get $a=b$.
Thus, $f$ is one-to-one.
\end{answer}
\end{exer}

\begin{proc}
To show that a function $f$ is {\bf not} one-to-one, we simply need to find two values $a\neq b$ in the
domain such that $f(a)=f(b)$. That is, we just need to show that there are two different numbers in the domain
that are mapped to the same value in the codomain.
\end{proc}

\begin{exa}\label{exa:xsquarenotoneone}
Let $f(x)=x^2$ be a function on the integers.  Show that $f$ is not one-to-one.
\begin{sol}
Notice that $f(-1)=f(1)=1$.  Thus, $f(x)$ is not one-to-one.
\end{sol}
\end{exa}

\begin{exer}
Let $f(x)=\lfloor x \rfloor$ be a function on $\BBR$.  Prove that $f$ is not one-to-one.

\noindent
\pflinetwo
\begin{answer}
Notice that $f(4.5)=f(4)=4$, so clearly $f$ is not one-to-one.
(Your proof may involve different numbers, but should be this simple.)
\end{answer}
\end{exer}

\begin{proc}
To show that a function $f$ is onto, we need to show that for an arbitrary $b\in B$, there is some
$a\in A$ such that $f(a)=b$.  That is, show that every value in $B$ is mapped to by $f$.
\end{proc}

\begin{exa}
Let $f(x)=x^3$ be a function on the real numbers.  Show that $f$ is onto.
\begin{sol}
Let $b\in \reals$.  
Then $f\left(\sqrt[3]{b}\right)=\left(\sqrt[3]{b}\right)^3 = b^{3/3}=b$.  
Since every $b\in\reals$ is mapped to (from $\sqrt[3]{b}$), $f$ is onto.
\end{sol}
\end{exa}

\begin{exer}
Let $f(x)=2x+1$ be a function on $\reals$.  Show that $f$ is onto.

\noindent
\pflinefour
\begin{answer}
Notice that if $y=2x+1$, then $y-1=2x$ and $x=(y-1)/2$.
Let $b\in\reals$.  Then $f((b-1)/2)=2((b-1)/2)+1 = b-1+1=b$.
Thus, every $b\in\reals$ is mapped to by $f$, so $f$ is onto.
\end{answer}
\end{exer}

\begin{proc}
To show that a function $f$ is {\bf not} onto, we just need to find some $b\in B$ such that there is no
$a\in A$ with $f(a)=b$.  In other words, we just need to find one value that isn't mapped to by $f$.
\end{proc}

\begin{exa}
Let $f(x)=x^3$ be a function on the integers.  Show that $f$ is not onto.
\begin{sol}
There is no integer $a$ such that $a^3=2$.  In other words, $2$ is not mapped to.
Thus, $f(x)$ is not onto.
\end{sol}
\end{exa}

\begin{exer}
Let $f(x)=\lfloor x \rfloor$ be a function on $\reals$.  Prove that $f$ is not onto.

\noindent
\pflinetwo
\begin{answer}
Since the floor of any number is an integer, there is no $a$ such that $f(a)=4.5$ (for instance).
Thus, $f$ is not onto.
\end{answer}
\end{exer}

It is important to remember that whether or not a function is one-to-one or onto might depend
on the domain/codomain over which the function is defined.
For instance, notice that in the last two examples we used the same function but on different domains/codomains.
In one case the function was onto, and in the other case it wasn't.

\newpage
\begin{exer}
Consider the function $f(x)=x^2$.
\begin{enumerate}[(a)]
\item
Prove or disprove that $f(x)=x^2$ is one-to-one on $\BBZ$.

\noindent
\answerlinefour
\item
Prove or disprove that $f(x)=x^2$ is one-to-one on $\BBR$.

\noindent
\answerlinefour
\item
Prove or disprove that $f(x)=x^2$ is one-to-one on $\naturals$.

\noindent
\answerlinefour
\end{enumerate}
\begin{answer}
(a) {\bf $f$ is not one-to-one.}  See Example~\ref{exa:xsquarenotoneone} for a proof.
(b) The same proof from Example~\ref{exa:xsquarenotoneone} works over the reals.
But I guess it doesn't hurt to repeat it:  Since $f(-1)=f(1)=1$, {\bf $f$ is not one-to-one.}
(c) Let $a,b\in\naturals$.  If $f(a)=f(b)$, 
that means $a^2=b^2$. Taking the square root of both sides, we obtain $\sqrt{a^2}=\sqrt{b^2}$, or
$|a|=|b|$  (if you didn't remember that
$\sqrt{x^2}=|x|$, you do now).  
But since $a,b\in\naturals$, $|a|=a$ and $|b|=b$.  Thus, $a=b$.  Thus, {\bf $f$ is one-to-one.}
\end{answer}
\end{exer}

\begin{exer}
Let $f(x)=3x-5$ be a function over $\reals$.  
Prove that $f$ has an inverse and then find it.

\vspace*{2in}
\begin{answer}
If $f(a)=f(b)$, $3a-5=3b-5$.  Subtracting 5 from both sides and then dividing both sides by $3$,
we get $a=b$.  Thus, $f$ is one-to-one.  
If $b\in\reals$, notice that $f((b+5)/3)=3((b+5)/3)-5 = b+5-5=b$, so there is some value that maps
to $b$.  Therefore, $f$ is onto.  Since $f$ is one-to-one and onto, it has an inverse.
To find the inverse, we
let $y=3x-5$.  Then $3x=y+5$, so $x=(y+5)/3$.  Thus, $f^{-1}(x)=(x+5)/3$
(or $\frac{x}{3}+\frac{5}{3}$).
\end{answer}
\end{exer}


\begin{exer}
Determine which of the following functions from $\integers$ to $\integers$ is one-to-one and/or onto.
Prove your answers.
\begin{enumerate}[(a)]
\item
$f(x)=x-7$

\noindent
\answerlinefour
\item
$g(x)=x^4$

\noindent
\answerlinefour
\item
$h(x)=3x$

\noindent
\answerlinefour
\item
$r(x)=\lfloor x/2\rfloor$

\noindent
\answerlinefour
\end{enumerate}
\begin{answer}
(a)
Notice that if $f(a)=f(b)$, then $a-7=b-7$ so $a=b$.  Thus, $f$ is one-to-one.
Also notice that for any $b\in\integers$, $f(b+7)=b+7-7=b$, so $f$ is onto.
(b) Since $g(1)=g(-1)=1$, $g$ is not one-to-one.  Also notice that there is no integer $a$ such that
$g(a)=a^4=5$, so $g$ is not onto.
(c) If $h(a)=h(b)$, then $3a=3b$ so $a=b$.  Thus, $h$ is one-to-one.
But there is no integer $a$ such that $h(a)=3a=1$, so $a$ is not onto.
(d) Notice that $r(0)=\lfloor 0/2\rfloor=\lfloor 0\rfloor = 0$ and 
$r(1)=\lfloor 1/2\rfloor=\lfloor 0\rfloor = 0$, so $r$ is not one-to-one.
But for any integer $b$, $r(2b)=\lfloor 2b/2\rfloor=\lfloor b\rfloor = b$, so $r$ is onto. 
\end{answer}
\end{exer}


\newpage
\begin{exa}
Let $f$ be a function from $B$ to $C$, and $g$ be a function
from $A$ to $B$.  If both $f$ and $g$ are one-to-one, prove
that $f\circ g$ is one-to-one.
\begin{quote}
\noindent 
{\bf Direct Proof:}

\noindent
For any distinct elements $x,y\in A$,
$g(x)\neq g(y)$, since $g$ is one-to-one.
Since $f$ is also one-to-one, then $f(g(x))\neq f(g(y))$, 
which is the same as $(f\circ g)(x) \neq (f\circ g)(y)$.
Therefore $f\circ g$ is one-to-one.\hfill$\square$

\bigskip
\noindent 
{\bf Proof by Contradiction:}

\noindent
Assume $f\circ g$ is not one-to-one.  Then there exist distinct elements
$x,y\in A$ such that $(f\circ g)(x)=(f\circ g)(y)$. This is equivalent
$f(g(x))=f(g(y))$.  Since $f$ is one-to-one, it must be the case that 
$g(x)=g(y)$. But $x\neq y$, and $g$ is one-to-one, so $g(x)\neq g(y)$.
This is a contradiction.  Therefore $f\circ g$ is one-to-one.\hfill$\square$
\end{quote}
\end{exa}

%%%%%%%%%%%%%%%%%%%%%%%%%%%%%%%%%%%%%%%%%%%
\newpage
\section{Partitions and Equivalence Relations}\label{section:parts}
%%%%%%%%%%%%%%%%%%%%%%%%%%%%%%%%%%%%%%%%%%%
Partitions and equivalence relations are not only fun and
interesting to learn about, they have various applications,
including some related to software testing that we will explore later.  

\begin{df}
Let $S \neq \varnothing$ be a set. A {\bf partition}\index{partition}\index{set!partition}
of $S$ is a collection of non-empty, pairwise  disjoint
subsets  of $S$ whose union is $S$.
\end{df}

\begin{exa}\label{exa:parityPartex}
Define $\BBE=\{2k:k\in \BBZ\}$ and $\BBO=\{2k+1:k\in \BBZ\}$.  Clearly $\BBE$ is the set of
even integers and $\BBO$ is the set of odd integers.
Since $\BBE\cap\BBO=\varnothing$ and  $\BBE\cup\BBO=\BBZ$,
$\{\BBE, \BBO\}$ is a partition of $\BBZ$.
Put another way, we can partition the integers based on parity.
\label{ex:mod2}
\end{exa}

\begin{exa}\label{exa:socksPart}
We can partition the socks in our sock drawer by color.  In other words, we put all of the
black socks in one set, the white ones in another, the green ones in another, etc.
For simplicity, we can put all of the multi-color socks in a single set.
\end{exa}


\begin{exa}\label{exa:peoplePart}
We can partition the set of all humans by putting each person into a set based on the first
letter of their first name.  So {\em Adam} and {\em Adele} go into set $A$ 
and {\em Zeek} goes into set $Z$, for instance.  The sets in the partition are 
$A, B, \ldots Z$.\footnote{For simplicity, we assume everyone's name is written using the Roman alphabet.}
\end{exa}

\begin{exa}\label{exa:weirdPart}
Let $A=\{1, 5, 8\}$, $B=\{2, 3\}$, $C=\{4\}$, $D=\{6, 9\}$, and $E=\{7, 10, 11, 12\}$.
Then the sets $A$, $B$, $C$, $D$, and $E$ form a partition of the set
$\{1, 2, 3, 4, 5, 6, 7, 8, 9, 10, 11, 12\}$.
\end{exa}



\begin{exa}
When choosing test cases for the \verb+factorial+ method in Example~\ref{exa:factPart}, 
we thought about 3 subsets of $\BBZ$: $\{0\}$, $\BBZ^{+}$, and $\BBZ^{-}$.  
These cases form a partition of $\BBZ$ since they are disjoint and $\BBZ=\{0\}\cup\BBZ^{+}\cup\BBZ^{-}$.
This is good since it means we covered at least one value of the different types, and we didn't `overtest'
any of the cases by unknowingly duplicating values from the same case.
\end{exa}

\begin{exer}\label{exer:maxPart}
You must decide on test cases for a method \verb+int maximum(int a,int b)+ that returns the
maximum of its arguments.  How would you partition the possible inputs into sets such that if
it is correct for one (or a few) tests of cases from that set, it is probably correct for the
rest of the cases in that set?  Notice that the set of inputs is $\BBZ\times\BBZ$.

\noindent
\answerlinethree
\begin{answer}
The following three cases probably make the most sense: When $a=b$, when $a<b$ and when $a>b$.
These make sense because these are likely different cases in the code.
Mathematically, we can think of it as follows.  The possible inputs are from the set $\BBZ\times\BBZ$.
The partition we have in mind is $A=\{(a,a) : a\in \BBZ\}$, $B=\{(a,b) : a,b\in \BBZ, a<b\}$, and
$C=\{(a,b) : a,b\in \BBZ, a>b\}$.  Convince yourself that these sets form a partition of $\BBZ\times\BBZ$.
That is, they are all disjoint from each other and $\BBZ\times\BBZ =A\cup B\cup C$.  

Alternatively, you might have thought in terms of $a$ and/or $b$ being positive, negative, or $0$.
Although that may make some sense, given that we are comparing $a$ and $b$ with each other, it probably
doesn't matter exactly what values $a$ and $b$ have (i.e. whether they are positive, negative, or $0$),
but what values they have {\em relative to each other}.  That is why the first answer is much better.
With that being said, it wouldn't hurt to include several tests for each of our three cases that involve
various combinations of positive, negative, and zero values.
\end{answer}
\end{exer}

Most of the partitions we talk about will be based on some meaningful characteristic of the
elements of a set--like parity, color, or sign.  
But this is not inherent in the definition.
For instance, the sets in the partition from Example~\ref{exa:weirdPart} do not seem to have
any significant meaning.
Some, like the one in Example~\ref{exa:parityPartex}, will have a precise mathematical definition.  
Others, like the one in Examples~\ref{exa:socksPart} will not.

\begin{exer}
Define a partition on $\BBZ$ that contains more than one subset.

\noindent
\answerlinethree
\begin{answer}
Did you define two or more subsets of $\BBZ$? 
Are they all non-empty?
Do none of them intersect with each other?
If you take the union of all of them, do you get $\BBZ$?  
If so, your answer is correct!  If not, try again.
\end{answer}
\end{exer}

\begin{exa}\label{exa:mod3partition}
Let $3\BBZ =\{3k:k\in \BBZ\}$, 
$3\BBZ +1=\{3k+1:k\in \BBZ\}$, and 
$3\BBZ +2=\{3k+2:k\in \BBZ\}$.\footnote{The notation in this example may seem a bid odd at first.
How are you supposed to interpret ``$3\BBZ +1$''?  Is this 3 times the set $\BBZ$ plus 1?
What does it mean to do algebra with sets and numbers? 
I won't get into all of the technical details, but here is a short answer.
You can think of ``$3\BBZ +1$'' as just a name.  
Sure, it may seem like an odd name, but why can't we name a set whatever we want?
Some people name their kids {\em Jon Blake Cusack 2.0} and get away with it.
You can also think of ``$3\BBZ +1$'' as describing how to create the set---by taking
every element from $\BBZ$, multiplying it by 3, and then adding 1.
Thus, you can think of  ``$3\BBZ +1$'' as being both an algebraic expression and a name.
}   
Since 
$$(3\BBZ)\cup (3\BBZ + 1) \cup (3\BBZ + 2) = \BBZ \mbox{ and}$$ 
$$ (3\BBZ)\cap (3\BBZ + 1) = \varnothing, \ (3\BBZ)\cap (3\BBZ + 2) = \varnothing,
(3\BBZ + 1)\cap (3\BBZ + 2) = \varnothing,$$ 
  $\{3\BBZ, 3\BBZ
+ 1, 3\BBZ + 2\}$ is a partition of $\BBZ$. \label{ex:mod3}
\end{exa}

\begin{exer}
Let $\BBI=\BBR \setminus \BBQ$ (the set of irrational numbers).
Prove that $\{\BBQ,\BBI\}$ is a partition of $\BBR$.

\noindent
\pflinethree
\begin{answer}
Since $\BBR = \BBQ \cup \BBI$ and $\BBQ \cap \BBI = \varnothing$, $\{\BBQ,\BBI\}$ is a partition of $\BBR$.
Hopefully this comes as no surprise.   
\end{answer}  
\end{exer}

Recall that when a list of number is given between parentheses (e.g. $(1,2,3)$), it typically denotes
an ordered list.  That is, the order that the element are listed matters.  So, for instance, $(1,2)$ 
and $(2,1)$ are not the same thing.

Next we will develop an alternative way of thinking about partitions: equivalence relations.
After defining some terms and providing a few examples, we will make the connection between
partitions and equivalence relations more clear.
 
\begin{df}
Let $A, B$ be sets. A {\bf relation} (or {\bf binary relation}) 
from $A$ to $B$ is a subset of the Cartesian product $A\times B$.
\index{relation}
\index{set!relation}

Given a relation $R$, we say that $x$ is {\bf related to} $y$ 
if $(x, y)\in R$.  We sometimes write this as  $x  R y.$  
An alternative notation is $x\sim y$.

If $R$ is a relation from $A$ to $A$, we sometimes say $R$ is a relation {\bf on} $A$.
\end{df}

\begin{exa}
Let $A$ be the set of all students at this school and $B$ be the set of all courses at this school.
We can define a relation $R$ by saying that $x R y$ if student $x$ has taken course $y$.
Said another way, 
we can define $R$ by saying that $(x,y)\in R$ if student $x$ has taken course $y$.
\end{exa}


\begin{exa}
We can define a relation $R=\{(a,a^2) : a\in \BBZ\}$.
That is, $x$ is related to $y$ if $y=x^2$.
\end{exa}

\begin{exa}\label{exa:parityRelation}
We can define a relation on $\BBZ$ by saying that $x$ is related to $y$ if they have the same parity.
Thus, $(2,0)$, $(234,-342)$, $(3,17)$ are all in $R$, but $(2,127)$ is not.
\end{exa}

\begin{ques}
Define $R=\{(a,b) : a,b\in \BBZ \mbox{ and } a<b\}$.  Is $R$ a relation?  Explain.

\noindent
\answerlinetwo
\begin{answer}
$R$ is a subset of $\BBZ\times\BBZ$, so it is a relation.
By the way, this relation should look familiar.  
Did you read the solution to Exercise~\ref{exer:maxPart}?
\end{answer}
\end{ques}

\begin{ques}
Is $\{(1,2),(345,7), (43,8675309), (11,11)\}$ a relation on $\BBZ^+$?  Explain.

\noindent
\answerlinetwo
\begin{answer}
Is it a subset of $\BBZ^+\times\BBZ^+$?  It is.  So it is a relation on $\BBZ^+$.
\end{answer}
\end{ques}

\begin{df}
\index{relation!reflexive}
\index{reflexive relation}
A relation $R$ on set $A$ is said to be
{\bf reflexive} if 
for all $x\in A$,
$x R x$ (or $(x,x)\in R$).
\end{df}


\begin{exer}\label{exa:propsExa}
Let $P$ be the set of all people.  
Which of the following relations on $P$ are reflexive?  Explain why or why not.
%\footnote{You
%may notice that the notation used to describe these is not consistent.  That is intentional.
%Since you will likely see different notations in different places, it is good to get used to them %all.}
\begin{lenum}
\item $T= \{(a,b) : a,b\in P \mbox{ and $a$ is taller than $b$}\}$
\item $N$ is the relation with $a$ related to $b$ iff $a$'s name starts with the same letter as $b$'s name.
\item $C$ is the relation defined by $(a,b)\in C$ if $a$ and $b$ have been to the same city.
\item $K = \{(a,b) : a,b\in P \mbox{ and $a$ does not know who $b$ is}\}$
\item $R = \{(\mbox{Barack Obama}, \mbox{George W. Bush})\}$.
\end{lenum}
\begin{lenum}
\item $T$: \blanklinefill  

\blanklinefill
\item $N$: \blanklinefill 

\blanklinefill
\item $C$: \blanklinefill

\blanklinefill
\item $K$: \blanklinefill

\blanklinefill
\item $R$: \blanklinefill

\blanklinefill
\end{lenum}
\begin{answer}
(a) $T$ is {\bf not} reflexive since you cannot be taller than yourself.
(b) $N$ is reflexive because everybody's name starts with the same letter as their name does.
(c) $C$ is reflexive because everybody have been to the same city as they have been in.
(d) $K$ is {\bf not} reflexive because you know who you are, 
so it is not the case that you don't know who you are.  That is, $(a,a)\not\in K$ for any $a$.
(e) $R$ is {\bf not} reflexive because $(\mbox{Donald Knuth}, \mbox{Donald Knuth})$ 
(for instance) is not in the relation.
\end{answer}
\end{exer}



\begin{df}
\index{relation!symmetric}
\index{symmetric relation}
A relation $R$ on set $A$ is said to be
{\bf symmetric} if 
for all $x,y\in A$,
$x R y$ implies $y  R x$
(or $(x,y)\in R$ implies $(y,x)\in R$).
\end{df}


\begin{exer}
Which of the relations from Example~\ref{exa:propsExa} are symmetric?  Explain why or why not.
\begin{lenum}
\item $T$: \blanklinefill  

\blanklinefill
\item $N$: \blanklinefill 

\blanklinefill
\item $C$: \blanklinefill

\blanklinefill
\item $K$: \blanklinefill

\blanklinefill
\item $R$: \blanklinefill

\blanklinefill
\end{lenum}
\begin{answer}
(a) $T$ is {\bf not} symmetric since if $a$ is taller than $b$, $b$ is clearly not taller than $a$.
(b) $N$ is symmetric since if $a$'s name starts with the same letter as $b$'s name, clearly $b$'s name
starts with the same letter as $a$'s name.
(c) $C$ is symmetric since it is worded such that it doesn't distinguish between the first and second
item in the pair.  In other words, if $a$ and $b$ have been to the same city, then $b$ and $a$ have 
been to the same city.
(d) $K$ is {\bf not} symmetric since $(\mbox{David Letterman},\mbox{Chuck Cusack})\in K$, but
$(\mbox{Chuck Cusack},\mbox{David Letterman})\not\in K$.
(e) $R$ is not symmetric since $(\mbox{Barack Obama}, \mbox{George W. Bush})\in R$, but
(George W. Bush, Barack Obama)$\not\in R$.
\end{answer}
\end{exer}

\begin{df}
\index{relation!anti-symmetric}
\index{anti-symmetric relation}
A relation $R$ on set $A$ is said to be
{\bf anti-symmetric} if
for all $x,y\in A$,
$x R y \mbox{ and } y  R x $ implies   $x = y$ 
(or $(x,y)\in R \mbox{ and }  (y,x)\in R$ implies $x=y$).
\end{df}

\begin{ques}
Let $R$ be a relation on $\BBZ$.
\begin{lenum}
\item
If $(1,1)\in R$, can you tell whether or not $R$ is anti-symmetric? Explain.

\noindent
\answerlinetwo
\item
What if $(1,2)$ and $(2,1)$ are both in $R$?  Can you tell whether or not
$R$ is anti-symmetric?

\noindent
\answerlinetwo
\end{lenum}
\begin{answer}
(a) Just knowing that $(1,1)\in R$ is not enough to tell either way.
(b) On the other hand, if $(1,2)$ and $(2,1)$ are both in $R$, it is clearly
{\bf not} anti-symmetric.
\end{answer}
\end{ques}

\begin{ques}
An alternative definition of {\em anti-symmetric} is  that 
if $x\neq y$, then $(x,y)$ and $(y,x)$ are not both in the relation.
Why is this definition equivalent?

\noindent
\answerlinetwo
\begin{answer}
This is just the contrapositive of the original definition.
\end{answer}
\end{ques}

\begin{rem}
The definition of anti-symmetric is sometimes misunderstood.  Let's call elements of the form $(x,x)$ {\bf diagonal} elements
and elements of the form $(x,y)$ where $x\neq y$ {\bf off-diagonal} elements.\footnote{These terms
come from thinking about the elements of a relation as elements in a matrix indexed by the members of
the set.  If this doesn't make sense, don't worry too much about it.}  
Then the definition of anti-symmetric is only dealing with off-diagonal elements.
It is saying nothing about the diagonal elements.
In other words, it is {\bf not} saying that $(x,x)\in R$ for any, let alone all, values of $x$.
But it also isn't saying $(x,x)\not\in R$.
It is simply saying that the {\em only} way for both $(x,y)$ and $(y,x)$ to be in $R$ is if $x=y$.

The alternative definition given in the previous question may help a little.  Notice that the
definition there starts with `if $x\neq y$...'  So what does the definition say about the case
$x=y$?  Nothing.  It never mentions it.

You could redefine it as follows: $R$ is anti-symmetric if for all non-diagonal elements $(x,y)\in R$,
$(y,x)\not\in R$.  But that can be problematic if you forget that $x\neq y$ is required.
\end{rem}
\begin{exer}
Which of the relations from Example~\ref{exa:propsExa} are anti-symmetric?  Explain why or why not.
\begin{lenum}
\item $T$: \blanklinefill  

\blanklinefill
\item $N$: \blanklinefill 

\blanklinefill
\item $C$: \blanklinefill

\blanklinefill
\item $K$: \blanklinefill

\blanklinefill
\item $R$: \blanklinefill

\blanklinefill
\end{lenum}
\begin{answer}
(a) $T$ is anti-symmetric since whenever $a\neq b$, if $a$ is taller than $b$, then $b$ is not taller than $a$,
so if $(a,b)\in T$, then $(b,a)\not\in T$.  
(b) $N$ is {\bf not} anti-symmetric since $(\mbox{Bono},\mbox{Boy George})$ and $(\mbox{Boy George},\mbox{Bono})$ 
are both in $N$.
(c) $C$ is {\bf not} anti-symmetric since $(\mbox{Bono},\mbox{The Edge})$ and $(\mbox{The Edge},\mbox{Bono})$ 
are both in $C$ (since they have played many concerts together, they have certainly been in the same
city at least once).
(d) Since
both $(\mbox{Dirk Benedict},\mbox{Jon Blake Cusack 2.0})$ and
$(\mbox{Jon Blake Cusack 2.0},\mbox{Dirk Benedict})$ are in $K$,
$K$ is {\bf not} anti-symmetric.
%(d) $K$ is {\bf not} anti-symmetric because both $(\mbox{Brian Yurk},\mbox{Chuck Cusack})$ and
%$(\mbox{Chuck Cusack},\mbox{Brian Yurk})$ are in $K$.
(e) $R$ is anti-symmetric since it only contains one element, 
$(\mbox{Barack Obama}, \mbox{George W. Bush})$, and
(George W. Bush, Barack Obama)$\not\in R$.
\end{answer}
\end{exer}


\begin{ques}
Answer each of the following.  Include a brief justification/example.
\begin{enumerate}[(a)]
\item If a relation is not symmetric, is it anti-symmetric?

\noindent
\answerlinetwo
\item If a relation is not anti-symmetric, is it symmetric?

\noindent
\answerlinetwo
\item Can a relation be both symmetric and anti-symmetric?

\noindent
\answerlinetwo
\end{enumerate}
\begin{answer}
(a) No.  The relation $R=\{(1,2),(2,1),(1,3)\}$ is neither symmetric ($(3,1)\not\in R$)
nor anti-symmetric ($(1,2)$ and $(2,1)$ are both in $R$).
(b) No.  For example, $R$ from answer (a) is not anti-symmetric, but isn't symmetric either.
(c) Yes.  If you answered incorrectly, don't worry.  You get to think about why the answer
is `yes' in the next exercise.
\end{answer}
\end{ques}


\begin{exer}
Give an example of a relation on any set of your choice that is both symmetric and
anti-symmetric.  Justify your answer.

\noindent
\answerlinetwo
\begin{answer}
Many answers will work, but they all have the same thing in common: They only contain
`diagonal' elements (but not necessarily all of the diagonal elements).
For instance, let $R=\{(a,a): a\in \BBZ\}$.
Go back to the definitions for symmetric and anti-symmetric and verify that this is indeed both.
Another examples is $R=\{(\mbox{Ken}, \mbox{Ken})\}$ on the set of English words.
\end{answer}
\end{exer}



\begin{df}
\index{relation!transitive}
\index{transitive relation}
A relation $R$ on set $A$ is said to be
{\bf transitive} if
for all $x,y,z\in A$,
$x R y\mbox{ and }  y R z$ implies $x R z$
(or $((x,y)\in R \mbox{ and }  (y,z)\in R)$ implies $(x,z)\in R$).
\end{df}

\begin{exer}
Which of the relations from Example~\ref{exa:propsExa} are transitive?  Explain why or why not.
\begin{lenum}
\item $T$: \blanklinefill  

\blanklinefill
\item $N$: \blanklinefill 

\blanklinefill
\item $C$: \blanklinefill

\blanklinefill
\item $K$: \blanklinefill

\blanklinefill
\item $R$: \blanklinefill

\blanklinefill
\end{lenum}
\begin{answer}
(a) $T$ is transitive since if $a$ is taller than $b$, and $b$ is taller than $c$, clearly $a$ is taller
than $c$.  In other words $(a,b)\in R$ and $(b,c)\in R$ implies that $(a,c)\in R$.
(b) $N$ is transitive because if $a$'s name starts with the same letter as $b$'s name, and
$b$'s name starts with the same letter as $c$'s name, clearly it is the same letter in all of them,
so $a$'s name starts with the same letter as $c$'s.
(c) $C$ is {\bf not} transitive.  You might think a similar argument as in (a) and (b) works here, but
it doesn't.  The proof from (b) works because names start with a single letter, so
transitivity holds.  But if $(a,b)\in C$ and $(b,c)\in C$, it might be because $a$ and $b$ have both
been to Chicago, and $b$ and $c$ have both been to New York, but that $a$ has never been to New York.
In this case, $(a,c)\not\in C$.  So $C$ is not transitive.
(d) $K$ is not transitive.  For instance, $(\mbox{David Letterman},\mbox{Chuck Cusack})\in K$
and $(\mbox{Chuck Cusack},\mbox{David Letterman's son})\in K$, but $(\mbox{David Letterman},\mbox{David Letterman's son})\not\in K$ since I sure hope he knows his own son.
(e) $R$ is transitive since there isn't even an $a,b,c\in R$ such that $(a,b)$ and $(b,c)$ are both in $R$,
so it holds vacuously.
\end{answer}
\end{exer}



\begin{df}
A relation which is reflexive, symmetric and transitive is
called an {\bf equivalence relation}.\index{equivalence relation}
\index{relation!equivalence} 
\end{df}


\begin{exa}
Let $S = $\{All Human Beings\}, and define 
the the relation $M$ by $(a,b)\in M$ if $a$ has the same (biological) 
mother\footnote{The important assumption we are making is that each person has exactly one mother.}
 as $b$.
Show that $M$ is an equivalence relation.
\begin{pf}
({\bf Reflexive}) 
$a$ has the same mother as $a$, 
so $(a,a)\in M$ and 
$M$ is reflexive.

\noindent
({\bf Symmetric})
If $a$ has the same mother as $b$, then $b$ clearly has the same mother as $a$.
Thus, $(a,b)\in M$ implies $(b,a)\in M$, so $M$ is symmetric.

\noindent
({\bf Transitive})
If $a$ has the same mother as $b$, and $b$ has the same mother as $c$, 
then clearly $a$ has the same mother as $c$.  In other words, $(a,b)\in M$ and
$(b,c)\in M$ implies that $(a,c)\in M$, so $M$ is transitive.

\noindent
Since $M$ is reflexive, symmetric, and transitive, it is an equivalence relation.
\end{pf}
\end{exa}

\begin{exer}
Which of the relations from Example~\ref{exa:propsExa} are equivalence relations?  Explain why or why not.
\begin{lenum}
\item $T$: \blanklinefill  

\blanklinefill
\item $N$: \blanklinefill 

\blanklinefill
\item $C$: \blanklinefill

\blanklinefill
\item $K$: \blanklinefill

\blanklinefill
\item $R$: \blanklinefill

\blanklinefill
\end{lenum}
\begin{answer}
(a) $T$ is {\bf not} an equivalence relation since it is not symmetric.
(b) $N$ is an equivalence relation since it is reflexive, symmetric, and transitive. 
(c) $C$ is {\bf not} an equivalence relation since it is not transitive.
(d) $K$ is {\bf not} an equivalence relation since it is not reflexive, symmetric, or transitive.  
This one isn't even close!
(e) $R$ is {\bf not} an equivalence relation since it is not reflexive. 
\end{answer}
\end{exer}

\begin{df}
A relation which is
reflexive, anti-symmetric and transitive is called a {\bf partial
order}.\index{partial order}
\end{df}


\begin{exer}
Which of the relations from Example~\ref{exa:propsExa} are partial orders?  Explain why or why not.
\begin{lenum}
\item $T$: \blanklinefill  

\blanklinefill
\item $N$: \blanklinefill 

\blanklinefill
\item $C$: \blanklinefill

\blanklinefill
\item $K$: \blanklinefill

\blanklinefill
\item $R$: \blanklinefill

\blanklinefill
\end{lenum}
\begin{answer}
(a) $T$ is a {\bf not} partial order because it is not reflexive.
(b) $N$ is {\bf not} a partial order since it is not anti-symmetric.
(c) $C$ is {\bf not} a partial order since it is not anti-symmetric or transitive.
(d) $K$ is {\bf not} a partial order since it is not reflexive, anti-symmetric, or transitive.
(e) $R$ is {\bf not} a partial order since it is not reflexive.
\end{answer}
\end{exer}

\begin{exer}
Let $X$ be a collection of sets. Let $R$ be the relation on $X$
such that $A$ is related to $B$ if $A \subseteq B$.
Prove that $R$ is a partial order on $X$.
\begin{pf}
({\bf Reflexive}) \blanklinefill

\noindent
({\bf Anti-symmetric}) \blanklinefill

\noindent
\blanklinefill

\noindent
({\bf Transitive}) \blanklinefill

\noindent
\blanklinefill

\noindent
\blanklinefill
\end{pf}
\begin{answer}
In the following, $A$, $B$, and $C$ are elements of $X$.  As such, they are sets.

\noindent
({\bf Reflexive}) 
Since $A\subseteq A$, $(A,A)\in R$, so $R$ is reflexive.

\noindent
({\bf Anti-symmetric}) 
If $(A,B)\in R$ and $(B,A)\in R$, then we know that 
$A\subseteq B$ and $B\subseteq A$.  By Theorem~\ref{sets:containtmentproof}, this implies
that $A=B$.  Therefore $R$ is anti-symmetric.

\noindent
({\bf Transitive}) 
If $(A,B)\in R$ and $(B,C)\in R$, then $A\subseteq B$ and $B\subseteq C$.  But the definition
of $\subseteq$ implies that $A\subseteq C$, so $(A,C)\in R$, and $R$ is transitive.

\noindent
Since $R$ is reflexive, anti-symmetric, and transitive, it is a partial order.
\end{answer}
\end{exer}

Labeling the lines of these proofs with what property we are proving isn't strictly necessary.
However, it does make the proofs a little easier to read.


\begin{exer}
Consider the relation 
$R=\{(1,2), (1,3), (1,5), (2,2), (3,5), (5,5)\}$ on the set $\{1,2,3,4,5\}$.
Prove or disprove each of the following.
\begin{enumerate}[(a)]
\item $R$ is reflexive

\noindent \answerlinetwo
\item $R$ is symmetric

\noindent \answerlinetwo
\item $R$ is anti-symmetric

\noindent \answerlinetwo
\item $R$ is transitive

\noindent \answerlinetwo
\item $R$ is an equivalence relation

\noindent \answerlinetwo
\item $R$ is a partial order

\noindent \answerlinetwo
\end{enumerate}
\begin{answer}
(a) Since $(1,1)\not\in R$, $R$ is not reflexive.
(b) Since $(1,2)\in R$, but $(2,1)\not\in R$, $R$ is not symmetric.
(c) A careful examination of the elements reveals that it {\em is} anti-symmetric.
(d) A careful examination of the elements reveals that it {\em is} transitive.
(e) Since it is not reflexive or symmetric, it is not an equivalence relation.
(f) Since it is not reflexive, it is not a partial order.
\end{answer}
\end{exer}

Next we show that congruence modulo $n$ (See Definition~\ref{df:congr}) is an equivalence relation.

\begin{thm}
Let $n$ be a positive integer.  Let $R$ be the relation on the set of integers
defined by $R=\{(a,b):a\equiv b \pmod n\}$. Then $R$ is an equivalence relation.
\begin{pf}
We need to show that $R$ is reflexive, symmetric, and transitive.

\noindent
({\bf Reflexive})
Clearly $a-a=0\cdot n$, so $a\equiv a\pmod n$.  Thus, $R$ is reflexive.

\noindent
({\bf Symmetric})
Assume $(a,b)\in R$.   Then $a\equiv b\pmod n$, which implies $a-b=kn$ for some integer $k$.  
So $b-a=(-k)n$, and since $-k$ is an integer, $b\equiv a\pmod n$.  Therefore, $(b,a)\in R$.  
Thus, $R$ is symmetric.

\noindent
({\bf Transitive})
Assume $(a,b), (b,c)\in R$.  Then 
$a\equiv b\pmod n$ and
$b\equiv c\pmod n$.
Thus,  $a-b=kn$ for some integer $k$ and 
$b-c=ln$ for some integer $l$. 
Given these, we can see that 
$$a-c = (a-b) + (b-c) = kn + ln = (k+l)n.$$
Since $k+l$ is an integer,  $a\equiv c \pmod n$. 
Thus $(a,c)\in R$, so  $R$ is transitive.
\end{pf}
\end{thm}

Notice that if we let $n=2$ in the previous theorem, 
we essentially have the relation from Example~\ref{exa:parityRelation}.

\begin{fib}
 Let $R$ be the relation on the set of ordered pairs of positive integers (that is, $\BBZ^+\times\BBZ^+$)
such that $( (a,b), (c,d) ) \in R$ if and only if $ad=bc$.
Show that $R$ is an equivalence relation.\footnote{In this example, $R$ is a relation on
a set of ordered pairs.  Thus, the elements of $R$ are ordered pairs of ordered pairs.  
Don't let this confuse you.  The elements of a relation are always ordered pairs.
What each part of the pair is depends on the underlying set.  If it is the set of animals, then the elements
of the relation are ordered pairs of animals.  If it is $\BBZ$, then the elements of the relation are
ordered pairs of integers.  And if it is $\BBZ^+\times\BBZ^+$, then the elements of the relation are
ordered pairs of ordered pairs of positive integers.
}
\begin{pf}
We need to show that $R$ is reflexive, symmetric, and transitive.

\noindent
({\bf Reflexive})
Since $ab=ba$ for all positive integers, $\blankline{1.5} \in R$ for all
$(a,b)$.  Thus $R$ is reflexive.

\noindent
({\bf Symmetric})
Assume  $( (a,b),(c,d) ) \in R$.  Then we know that $ad=\blankline{.5}.$
We can rearrange this as $cb=\blankline{.5}$.
Thus,  $\blankline{1}\in R$, so $R$ is \blankline{1.5}.

\noindent
({\bf Transitive})
Assume that $( (a,b),(c,d)) \in R$ and $( (c,d),(e,f) ) \in R$.  
Then we know that \blankline{1.5} and \blankline{1.5}.
Solving the second for $c$, we get $c=\blankline{1}$.
Plugging it into the first we get $ad=\blankline{1}$.  
Multiplying both sides by $f$, and canceling the $d$ on both sides yields
\blankline{1.5}.  
Thus, $\blankline{1.5} \in R$, so $R$ is transitive.
\end{pf}
\begin{answer}
$( (a,b), (a,b))$;
$bc$;
$da$;
$((c,d),(a,b))$;
symmetric;
$ad=bc$; 
$cf=de$;
$de/f$;
$b(de/f)$;
$af=be$;
$( (a,b),(e,f) )$
\end{answer}
\end{fib}

\begin{df}
\index{equivalence class}
Let $R$ be an equivalence relation on a set $S$. Then
the {\bf equivalence class of \boldmath$a$}, denoted by \boldmath$[a]$\unboldmath, is the subset of $S$ 
containing all of the elements that are related to $a$.
More formally, 
$$[a] = \{x\in S: x R a\}.$$
If $x\in[a]$, we say that $x$ is a {\bf representative} of the equivalence class $[a]$.
Note that {\em any} element of an equivalence class can serve as a representative.
\end{df}

\begin{exa}
The equivalence class of $3$ modulo $8$ is 
$[3]=\{8k+3 : k\in\BBZ\}$.
Notice that 
$[11]=\{8k+11 : k\in\BBZ\}=\{8k+3 : k\in\BBZ\}=[3]$.
In fact, $[3]=[8l+3]$ for all integers $l$.
In other words, any element of the form $8l+3$, where $l$ is an integer,
can serve as a representative of $[3]$.
Further, we can call this class $[3]$, $[11]$, $[19]$, etc.  It doesn't really
matter since they all represent the same set of integers.  
Of course, $[3]$ is the most logical choice.
\end{exa}

\begin{exa}
Notice that if our relation is congruence modulo $3$, we can define three
equivalence classes:
\begin{eqnarray*}
{[}0{]}&=&\{3k : k\in\BBZ\},\\
{[}1{]}&=&\{3k+1 : k\in\BBZ\}, \mbox{ and}\\
{[}2{]}&=&\{3k+2 : k\in\BBZ\}.\\
\end{eqnarray*}
It isn't too difficult to see that $\BBZ=[1]\cup[2]\cup[3]$,
and that these three sets are disjoint.
In other words, the equivalence classes $\left\{[1],[2],[3]\right\}$ form a partition of $\BBZ$.  
As we will see shortly, this is not a coincidence.
\end{exa}


\begin{lem}\label{lem:equiv}
Let $R$ be an equivalence relation on a set $S$. Then
two equivalence classes are either identical or disjoint.
\label{lem:equiv_classes}
\begin{pf}
Let $a, b\in S$, and assume $[a]\cap [b] \neq
\varnothing$ (that is, that they are not disjoint).  
We need to show that $[a] = [b]$. 
First, let $x\in [a]\cap [b]$ (which exists since $[a]\cap [b] \neq \varnothing$).
Then $x R a$ and $x R b$, so by symmetry $a R x$ and by transitivity $a R b$.

Now let $y\in [a]$.  Then $y R a$. Since we just showed that $a R b$,  
then $y R b$ by transitivity.  
Thus $y\in[b]$.
Therefore $[a]\subseteq [b]$.  

A symmetric argument proves that $[b]\subseteq [a]$.
Therefore, $[a]=[b]$.
%Santos' proof.
%Suppose that $x\in [a]\cap [b]$. Now
%$x\in [a] \rightarrow x R a  \rightarrow a R x$, by symmetry. Similarly,
%$x\in [b] \rightarrow x R b.$ By transitivity $$(a R x)\AND (x R
%b)  \rightarrow a R b.$$Now, if $y\in [b]$ then $b R y$. Again by
%transitivity, $a R y$. This means that $y\in [a]$. We have shown
%that $y\in [b] \rightarrow y\in [a]$ and so $[b]\subseteq [a]$. In a
%similar fashion, we may prove that $[a]\subseteq [b]$. This
%establishes the result.
\end{pf}
\end{lem}

Let's bring together some of the examples of partitions with examples of equivalence
relations and classes.

\begin{exa}
We just saw that congruence modulo $3$ is an equivalence relation with three
equivalence classes,  $\{3k : k\in\BBZ\}$, $\{3k+1 : k\in\BBZ\}$, and $\{3k+2 : k\in\BBZ\}$.
In Example~\ref{exa:mod3partition}, we defined a partition of $\BBZ$ using these same
three subsets.
\end{exa}

\begin{exa}
In Example~\ref{exa:parityRelation} we defined a relation on $\BBZ$ based on parity.
It is not difficult to see that the equivalence classes of that relation are $[0]=\BBE$
and $[1]=\BBO$.  Notice these are the same subsets we used to partition $\BBZ$ in 
Example~\ref{exa:parityPartex}.
\end{exa}

\begin{exa}
In Example ~\ref{exa:peoplePart} we defined a partition of people according to the first
letter of their first name.  The sets in the partition were $A, B, \ldots, Z$.

We can define an equivalence relation on the set of all people by saying $a$ is related to $b$ if
$a$'s name starts with the same letter of the alphabet as $b$'s name.
In a series of previous exercises, you proved that this defines an equivalence relation.
Notice that the equivalence classes are the sets $A, B, \ldots, Z$ 
(which we can think of as, for instance $[Adam], [Betty], \ldots, [Zeek]$).
Again, these are the same sets that we used to partition people into in Example~\ref{exa:peoplePart}.
\end{exa}

In these examples, there seems to be a connection between the equivalence classes of the 
relation and the sets in a partition.  As the next theorem illustrates, this is no coincidence.  


\begin{thm}\label{thm:equivPart}
Let $S \neq \varnothing$ be a set. Every equivalence
relation on $S$ induces a partition of $S$ and vice-verse. 
\label{thm:equiv_relation_yields_partition}
\begin{pf}
By Lemma \ref{lem:equiv_classes}, if $R$ is an equivalence
relation on $S$ then
$$ S = \bigcup _{a\in S} [a],$$ and $[a]\cap [b] = \varnothing$ if $a$ is not related to $b$.
This proves the first half of the theorem.

Conversely, let $$ S = \bigcup _{\alpha} S_\alpha, \  \mbox{where}\ \
S_\alpha \cap S_\beta = \varnothing\ \ {\rm if}\ \alpha \neq
\beta,$$ be a partition of $S$. We define the relation
$R$ on $S$ by letting $a R b$ if and only if
they belong to the same $S_\alpha$. Since the $S_\alpha$ are
mutually disjoint, it is clear that $R$ is an equivalence
relation on $S$ and that for $a\in S_\alpha,$ we have
$[a] = S_\alpha$.
\end{pf}
\end{thm}

%Put in simple terms,
Equivalence classes of an equivalence relation and partitions of sets are
essentially the same thing.  The main difference is in how we look at it.  When thinking about
equivalence relations/classes, the focus is on what it means for two things to be related.
When thinking about partitions, the focus is on what it means for an element to be in a particular
subset of the partition.

\begin{exa}\label{exa:modex}
In light of Theorem~\ref{thm:equivPart}, 
we can say that the relation defined by congruence modulo $4$ 
partitions the set of integers into precisely $4$ equivalence classes: $[0]$, $[1]$, $[2]$, and $[3]$.
That is, given any integer, it is contained in one (and only one) of these classes.

More generally, if $n>2$, $\BBZ$ can be partitioned into $n$ sets, $[0],[1],\ldots,[n-1]$, each of
which is an equivalence class of the relation defined by congruence modulo $n$.

When we think about the partition, we are focused on the concept that each number $x$ goes into
one of the $n$ subsets based on the value $x\bmod n$.
On the other hand, when we think about the relation of congruence modulo $n$, we are focused on
the idea that $x$ and $y$ are in the same equivalence class iff $x\equiv y\pmod n$.
\end{exa}

For more discussion and some applications of partitions and equivalence relations, see Section~\ref{section:funAlgs}.
